% !TEX program = xelatex
\documentclass[12pt,a4paper]{ctexart}

% ================================================================
% Packages
% ================================================================
\usepackage{amsmath,amssymb,amsthm}
\usepackage{mathtools}
\usepackage{bm}
\usepackage{graphicx}
\usepackage{hyperref}
\usepackage{geometry}
\usepackage{enumitem}
\usepackage{paralist}
\usepackage{booktabs}
\usepackage{tikz}
\usepackage{xcolor}
\usepackage{cleveref}
\usepackage{bbm}
\usepackage{float}
\usepackage{physics}
\usepackage{algorithm}
\usepackage{algpseudocode}

\geometry{margin=2.5cm}

% ================================================================
% Hyperref setup
% ================================================================
\hypersetup{
    colorlinks=true,
    linkcolor=blue,
    citecolor=blue,
    urlcolor=blue,
    pdftitle={Potential Surface Misalignment: Gauge Freedom in MoE Routing and MILP Gauge Fixing},
    pdfauthor={SCX},
    pdfsubject={SCX MoE Gauge Theory},
}

% ================================================================
% Theorem environments
% ================================================================
\newtheorem{theorem}{Theorem}[section]
\newtheorem{lemma}[theorem]{Lemma}
\newtheorem{proposition}[theorem]{Proposition}
\newtheorem{corollary}[theorem]{Corollary}
\newtheorem{definition}[theorem]{Definition}
\newtheorem{conjecture}[theorem]{Conjecture}
\newtheorem{assumption_env}[theorem]{Assumption}

\theoremstyle{remark}
\newtheorem{remark}[theorem]{Remark}
\newtheorem{example}[theorem]{Example}
\newtheorem{protocol}[theorem]{Protocol}

% ================================================================
% Custom commands
% ================================================================
\newcommand{\R}{\mathbb{R}}
\newcommand{\C}{\mathbb{C}}
\newcommand{\N}{\mathbb{N}}
\newcommand{\E}{\mathbb{E}}
\newcommand{\Pbb}{\mathbb{P}}
\newcommand{\Var}{\mathrm{Var}}
\newcommand{\D}{\mathcal{D}}
\newcommand{\G}{\mathcal{G}}
\newcommand{\T}{\mathcal{T}}
\newcommand{\loss}{\mathcal{L}}
\newcommand{\SCX}{\mathsf{SCX}}
\newcommand{\MoE}{\mathsf{MoE}}
\newcommand{\gaugeGroup}{\mathbb{G}}

\DeclareMathOperator{\argmax}{arg\,max}
\DeclareMathOperator{\argmin}{arg\,min}
\DeclareMathOperator{\sgn}{sgn}
\DeclareMathOperator{\diag}{diag}
\DeclareMathOperator{\spec}{spec}
\DeclareMathOperator{\softmax}{softmax}
\DeclareMathOperator{\topk}{top\text{-}k}
\DeclareMathOperator{\Cov}{Cov}
\let\std\relax
\DeclareMathOperator{\std}{std}

% Honest critique
\newcommand{\honestcrit}[1]{%
    \textcolor{red}{\textbf{[Honest Critique]#1}}%
}

% Proof-status markers
\newcommand{\rigorFull}{\textnormal{\textbf{[Rigorous Proof]}}}
\newcommand{\rigorPartial}{\textnormal{\textbf{[Partial Proof]}}}
\newcommand{\rigorSketch}{\textnormal{\textbf{[Proof Sketch]}}}
\newcommand{\openproblem}{\textnormal{\textbf{[Open Problem]}}}

% Assumption tags
\newcommand{\assumptionRef}[1]{\textbf{A#1}}

% ================================================================
% Title
% ================================================================
\title{\textbf{Potential Surface Misalignment: Gauge Freedom and MILP Gauge Fixing in Multi-Expert Routing}\\[4pt]
\large Gauge Freedom in Mixture-of-Experts Routing: \\[2pt]
Why Expert Output Surfaces Are Not Aligned and How MILP Gauge Fixing Resolves It}
\author{SCX}
\date{July 1, 2026}

\begin{document}

\maketitle

\begin{center}
\footnotesize
\textbf{Version: } v1.0 \quad | \quad
\textbf{Status: } Preprint \quad | \quad
\textbf{Category: } SCX Theory Series — Information Theory Volume, Multi-Expert Gauge Chapter
\end{center}

% ================================================================
\begin{abstract}

Mixture-of-Experts (MoE) architectures achieve remarkable scale efficiency by routing each token to a subset of specialized sub-networks. However, we identify a fundamental and previously unrecognized problem: the \textbf{potential surface misalignment} (PSM) — different experts define their output surfaces in gauge-inequivalent coordinate systems, making the router's task of comparing expert relevance ill-posed. We formalize this as a \textbf{gauge freedom} in the representation space of deep MoE models: each expert's output admits a group of transformations $\G_m$ that leave its own training loss invariant but alter cross-expert comparisons. We prove that the standard linear router is not gauge-invariant, formulate optimal gauge fixing as a Mixed Integer Linear Program (MILP) with a polynomial-time convex relaxation and a greedy algorithm with provable error bounds, and establish that the SVD spectrum of gauge-aligned multi-expert outputs is a provable hallucination detector — distinguishing true knowledge from both disagreement-based and shared hallucinations. We connect this work to the author's prior gauge-fixing framework for atomic cluster expansion (ACE) potentials, showing that both are instances of a deeper \textbf{Modular Gauge Principle}: any system of independently trained components requires explicit gauge alignment before comparison, merging, routing, or aggregation. Three engineering pathways are derived: zero-training router repair, gauge-invariant distillation with Yajie consensus, and gauge-aligned representation distillation. Beyond engineering, we extract the \textbf{Equality Principle}: the mathematical necessity of gauge alignment at interfaces implies that inequality at contact points is structurally unstable — a principle that extends from MoE routing to social systems, explaining historical patterns of wealth-driven conflict as mathematically inevitable. The full chain comprises 13 theorems, 6 definitions, 4 propositions, and 7 corollaries, spanning the arc from an engineering observation about expert networks to a self-reflexive epistemological principle that judges any system — including those that would judge it.

\vspace{4pt}
\noindent\textbf{Keywords:} Mixture of Experts, gauge theory, potential surface alignment, MILP, hallucination detection, SVD spectrum, representation geometry

\end{abstract}

% ================================================================
\section*{Core Intuition: Understanding Potential Surface Misalignment in One Page}
\addcontentsline{toc}{section}{Core Intuition: Understanding PSM in One Page}
% ================================================================

\vspace{-8pt}

\begin{center}
\fbox{%
\begin{minipage}{0.92\textwidth}

\subsection*{The Ruler Metaphor}

Imagine you ask 8 engineers to each build a ruler, and then you try to piece these rulers together into one long ruler.

Each engineer's ruler has \textbf{accurate markings} -- 1 cm is indeed 1 cm -- but each person placed \textbf{the zero mark at a different position}. Zhang San's zero is at the far left. Li Si's zero is in the middle. Wang Wu's zero is shifted 2 mm to the right.

When you piece the 8 rulers together directly: the markings do not align, and jumps appear at the seams. Yet each ruler, used alone to measure something, is accurate -- because it measures \textbf{relative length}, not absolute position.

\textbf{This is Potential Surface Misalignment.}

\begin{itemize}
    \item \textbf{Ruler = MoE expert network} $E_m$. Each expert trains independently and is \"accurate\" in its own domain (low loss)
    \item \textbf{Zero-mark position = gauge freedom} $\mathbf{g}_m$. Training loss is insensitive to the zero-mark position -- because residual connections and LayerNorm automatically absorb the shift
    \item \textbf{Piecing rulers together = router comparing experts}. The router uses a linear function to compare the \"markings\" of 8 experts, but it does not know where each expert's zero is
    \item \textbf{Seam jumps = routing bias}. Choosing the wrong expert, or assigning incorrect weights
\end{itemize}

\textbf{The fix is simple:} First align all rulers to a common zero (gauge fixing), then piece them together (routing). Or -- more cleverly -- do not piece the rulers at all; instead, let each ruler report only the \textbf{final measurement} (distillation + Yajie), because the final measurement does not depend on the zero-mark position (gauge invariance).

\end{minipage}%
}
\end{center}

\vspace{8pt}

\begin{center}
\fbox{%
\begin{minipage}{0.92\textwidth}

\subsection*{Three-Sentence Summary}

\begin{enumerate}
    \item \textbf{Problem:} Different MoE experts live in their own \"coordinate systems.\" The router is comparing things that are incomparable. This is not an engineering oversight -- it is a mathematical structure: gauge freedom.
    \item \textbf{Solution:} Three approaches. (a) Repair the router -- zero training, just compute a bias on a calibration set. (b) Distillation + Yajie -- bypass the gauge problem, perform consensus denoising in output space. (c) Align then distill -- the most elegant approach, retaining maximal information.
    \item \textbf{Unexpected finding:} The SVD spectrum after gauge alignment can detect \"shared hallucinations\" that Yajie misses -- errors that all experts agree on. True knowledge is low-dimensional in representation space; shared hallucinations are high-dimensional. SVD can see this difference.
\end{enumerate}

\end{minipage}%
}
\end{center}

\vspace{12pt}

% ================================================================
\section{Introduction: From ACE Gauge to MoE Gauge}
\abel{sec:intro}
% ================================================================

In prior work~[EGP paper, SCX system literature], we identified and resolved a core problem in merging atomic cluster expansion (ACE) potentials: \textbf{coefficient-level gauge freedom}. Specifically, the shared-correction ACE parameterization
\begin{equation}
    E(\sigma) = \mathbf{c}_0 \cdot \mathbf{B}(\sigma) + \sum_Z \mathbf{c}_Z \cdot \mathbf{B}_Z(\sigma)
\end{equation}
remains invariant under the transformation
\begin{equation}
    \mathbf{c}_0 \to \mathbf{c}_0 + \mathbf{g},\quad \mathbf{c}_Z \to \mathbf{c}_Z - \mathbf{g}\quad(\forall Z)
\end{equation}
preserving all physical predictions. However, different independent training runs exploit this freedom to reach different parameter points. We resolved this via \textbf{post-hoc orthogonal projection}, reducing gauge violation to machine precision ($<10^{-15}$).

\honestcrit{But that is just one instance. Gauge freedom is far from limited to the ACE parameter space -- it is ubiquitous in all modular multi-component systems. MoE is the next target.}

\subsection{Potential Surface Misalignment: An Intuitive Statement of the Problem}

Consider a standard sparse MoE layer. Let there be $N$ experts $\{E_m\}_{m=1}^{N}$, each a feedforward network $E_m: \R^d \to \R^d$. The router produces scores $g(x) = \softmax(W_r x) \in \R^N$ and selects the top-$k$ experts $\topk(g(x), k)$ to activate.

\textbf{Core intuition}: Expert $E_1$ and expert $E_2$ each learn different output \"baselines\" during training -- $E_1$'s outputs are inherently \"high,\" $E_2$'s outputs are inherently \"low,\" even for similar inputs. This difference is adaptively absorbed by downstream layers via residual connections during separate training, and is therefore not perceived by the training loss. But --

\textbf{The router does not know this.} The router $W_r x$ is a linear map that compares some implicit proxy of the raw expert outputs. When $E_1$ and $E_2$ live in different coordinate systems, the router's comparison is mathematically \textbf{ill-posed}.

We call this phenomenon \textbf{Potential Surface Misalignment (PSM)}: each expert defines a function surface
\begin{equation}
    \mathcal{S}_m = \{(x, \|E_m(x)\|) \in \R^{d+1} : x \in \R^d\},
\end{equation}
and these surfaces naturally differ in height, scale, and even orientation in output space -- not because any expert is \"wrong,\" but because the training dynamics permit a gauge freedom.

\subsection{Identification of the Gauge Group}

We identify the following gauge freedoms in the MoE representation space:

\begin{definition}[Gauge Group of the MoE Representation Space]\label{def:gauge_group}
Let the output space of expert $E_m$ be $\R^d$. The \textbf{local gauge group} $\G_m$ of expert $m$ is a set of transformations satisfying: for any $\gamma \in \G_m$, there exists a downstream adaptation $\delta_\gamma$ in the residual connection such that
\begin{equation}
    \delta_\gamma \circ \gamma \circ E_m = E_m
\end{equation}
is indistinguishable under the training loss. Specifically:
\begin{enumerate}[label=(\roman*)]
    \item \textbf{Translation gauge} $\G_m^{\text{trans}} = \{T_{\mathbf{b}} : \mathbf{y} \mapsto \mathbf{y} + \mathbf{b}, \mathbf{b} \in \R^d\}$
    \item \textbf{Scale gauge} $\G_m^{\text{scale}} = \{S_{\alpha} : \mathbf{y} \mapsto \alpha \mathbf{y}, \alpha > 0\}$
    \item \textbf{Rotation gauge} $\G_m^{\text{rot}} = \{R_{\mathbf{Q}} : \mathbf{y} \mapsto \mathbf{Q}\mathbf{y}, \mathbf{Q} \in O(d)\}$
\end{enumerate}
The full gauge group is the semidirect product of these subgroups: $\G = \G^{\text{trans}} \rtimes (\G^{\text{scale}} \times \G^{\text{rot}})$.
\end{definition}

\honestcrit{Translation gauge is dominant -- Post-LN or Pre-LN LayerNorm absorbs translations. Rotation gauge is secondary -- downstream projection matrices can learn to undo rotations. Scale gauge is the most subtle -- it is partially masked by LayerNorm but still exists before residual summation.}

\subsection{Contributions of This Work}

\begin{enumerate}[label=\textbf{(C\arabic*)}]
    \item \textbf{Identifying the lack of gauge invariance in MoE routing.} We prove that the standard router $\softmax(W_r x)$ is not invariant under $\G$ -- routing scores change under gauge transformations, making expert outputs under different gauges incomparable.
    \item \textbf{MILP gauge fixing formulation.} We formulate optimal gauge fixing as a Mixed Integer Linear Program (MILP): integer variables encode routing decisions, continuous variables encode gauge parameters. We provide a tight convex relaxation and a polynomial-time approximation.
    \item \textbf{Gauge-aligned SVD hallucination detection.} We prove that after gauge fixing, the concentration of the SVD spectrum of the multi-expert output matrix for the same query is a provable hallucination detector -- a flat spectrum indicates no internal consensus.
    \item \textbf{Bridging ACE gauge and MoE gauge.} We show that both are instances of the same deep principle: any independently trained modular components require explicit gauge alignment before interaction.
\end{enumerate}

% ================================================================
\section{Problem Setup: Formalization of MoE}
\abel{sec:formalism}
% ================================================================

\begin{definition}[Sparse MoE Layer]\label{def:moe}
A sparse MoE layer consists of the following components:
\begin{itemize}
    \item $N$ expert functions $E_m: \R^d \to \R^d$, $m = 1, \ldots, N$, each a feedforward network
    \item A router $r: \R^d \to \Delta^{N-1}$, $r(x) = \softmax(W_r x)$, where $W_r \in \R^{N \times d}$
    \item Number of active experts $k \in \{1, \ldots, N\}$
\end{itemize}
For input token $x \in \R^d$, the output is
\begin{equation}
    y = \sum_{m \in \mathcal{A}(x)} r_m(x) \cdot E_m(x),
\end{equation}
where $\mathcal{A}(x) = \argmax_k r(x)$ is the index set of the $k$ highest-scoring experts.
\end{definition}

\begin{assumption_env}[Post-Training Gauge Fixing]\label{ass:posthoc}
We assume that post-hoc gauge fixing is performed after training has completed, similar to the post-hoc projection method in the EGP work. No gauge constraints are imposed during training; this has been shown to be suboptimal~[EGP paper, $\lambda$ scan failure].
\end{assumption_env}

\begin{assumption_env}[Residual Connections Absorb Gauge Differences]\label{ass:residual}
Let the input to the $\ell$-th MoE layer be $x^{(\ell)}$. The residual connection
\begin{equation}
    x^{(\ell+1)} = x^{(\ell)} + \sum_{m \in \mathcal{A}} r_m \cdot E_m(x^{(\ell)})
\end{equation}
enables gauge freedom to exist: upstream LayerNorm and downstream projection matrices can adaptively absorb gauge transformations of expert outputs, making the training loss locally insensitive to these transformations.
\end{assumption_env}

\begin{assumption_env}[Calibration Set Availability]\label{ass:calibration}
There exists a calibration dataset $\D_{\text{cal}} = \{x_i\}_{i=1}^{n_{\text{cal}}}$ of size $n_{\text{cal}} \geq N \cdot d$ that can be used to estimate gauge parameters. The calibration set does not require labels; it only requires input tokens. This matches common practical scenarios.
\end{assumption_env}

% ================================================================
\section{Formalization of Gauge Freedom}
\abel{sec:gauge_formal}
% ================================================================

\subsection{Gauge Non-Invariance of the Router}

\begin{theorem}[Gauge Non-Invariance of the Router]\label{thm:noninvariance}\rigorFull
Let the router be $r(x) = \softmax(W_r x)$, where $W_r \in \R^{N \times d}$ is fixed after training. Applying a translation gauge transformation $E_m \to E_m + \mathbf{g}_m$ to expert $m$ (where $\mathbf{g}_m \in \R^d$) does \textbf{not} change the routing score under the transformation; that is, $r(x)$ has zero explicit dependence on the output of $E_m$, because $r(x)$ does not receive $E_m(x)$ as input.

However, the router is \textbf{implicitly} tuned during training to match the output distribution of the experts. Specifically, the router gradient during training is
\begin{equation}
    \nabla_{W_r} \loss = \E_x\left[ \frac{\partial \loss}{\partial y} \cdot \frac{\partial y}{\partial r} \cdot \frac{\partial r}{\partial W_r} \right]
\end{equation}
where $\partial y / \partial r$ depends on the values of $E_m(x)$. Therefore, the optimal value of $W_r$ \textbf{depends on the output distribution of the experts during training}. When a gauge transformation $E_m \to E_m + \mathbf{g}_m$ is applied after training, the training-time $W_r^*$ is no longer the optimal router under the new gauge.

Specifically, let the expert outputs in the original gauge (during training) be $\{E_m^{\text{train}}\}$ and the router be $W_r^{\text{train}}$. After the gauge transformation, the expert outputs become $\{E_m^{\text{train}} + \mathbf{g}_m\}$. The difference in optimal routing loss between the two gauges is controlled by the following inequality:
\begin{equation}
    \left|\loss(W_r^{\text{train}}; \{E_m^{\text{train}} + \mathbf{g}_m\}) - \min_{W_r} \loss(W_r; \{E_m^{\text{train}} + \mathbf{g}_m\})\right| \leq L_r \cdot \max_m \|\mathbf{g}_m\|,
\end{equation}
where $L_r$ is the Lipschitz constant of $\loss$ with respect to $W_r$.
\end{theorem}

\begin{proof}
Consider the optimality condition of the router at the end of training. The loss function during training is
\begin{equation}
    \loss(W_r) = \E_{(x, y^*)} \left[ \ell\left( \sum_{m} \softmax_m(W_r x) \cdot E_m^{\text{train}}(x), \; y^* \right) \right].
\end{equation}

Suppose that, after training convergence, $W_r^{\text{train}} \approx \argmin_{W_r} \loss(W_r; \{E_m^{\text{train}}\})$. After the gauge transformation, the new loss landscape is
\begin{equation}
    \loss'(W_r) = \loss(W_r; \{E_m^{\text{train}} + \mathbf{g}_m\}).
\end{equation}

The key point is that $\loss(W_r^{\text{train}}; \{E_m^{\text{train}}\})$ reaches an approximate minimum at $W_r^{\text{train}}$. But $\loss'(W_r^{\text{train}})$ is \textbf{not necessarily} minimized at $W_r^{\text{train}}$, because the gauge transformation changes the loss landscape itself:
\begin{equation}
    \frac{\partial \loss'}{\partial W_r}\bigg|_{W_r^{\text{train}}} = \E\left[ \frac{\partial \ell}{\partial y} \cdot \sum_m \softmax_m \cdot (1 - \softmax_m) \cdot x^T \cdot (E_m^{\text{train}}(x) + \mathbf{g}_m)^T \right]
\end{equation}
the term $(E_m^{\text{train}}(x) + \mathbf{g}_m)$ differs from its training-time value.

To bound suboptimality, since $\loss$ is $L_r$-Lipschitz with respect to $W_r$ under reasonable regularization conditions, and the expert output change is bounded by $\|\delta E_m\| = \|\mathbf{g}_m\|$:
\begin{align}
    |\loss'(W_r) - \loss(W_r)| &\leq \E\left[ L_\ell \cdot \left\|\sum_m r_m \cdot \mathbf{g}_m\right\| \right] \\
    &\leq L_\ell \cdot \max_m \|\mathbf{g}_m\| \cdot \E\left[\sum_m r_m\right] \\
    &= L_\ell \cdot \max_m \|\mathbf{g}_m\|.
\end{align}

Therefore $|\loss'(W_r^{\text{train}}) - \min_{W_r} \loss'(W_r)| \leq 2 L_\ell \cdot \max_m \|\mathbf{g}_m\|$. Setting $L_r = 2L_\ell$ gives the result.
\end{proof}

\begin{corollary}[Accumulation of Routing Bias]\label{cor:cumulative}
Consider an $L$-layer Transformer with one MoE sublayer in each layer. If each layer has a gauge bias of size $\|\mathbf{g}_m^{(\ell)}\| \leq g_{\max}$, then cumulative routing bias can grow by at most $O(L \cdot g_{\max})$ over depth $L$, and in the worst case the terminal-layer expert selection deviates from the optimal selection by Hellinger distance $O(L \cdot g_{\max})$.
\end{corollary}

\subsection{Gauge Equivalence Classes}

\begin{definition}[Gauge Equivalence]\label{def:gauge_equiv}
Two expert configurations $\{E_m\}$ and $\{E_m'\}$ are called \textbf{gauge-equivalent}, denoted $\{E_m\} \sim_{\G} \{E_m'\}$, if there exist gauge transformations $\gamma_m \in \G_m$ such that $E_m' = \gamma_m \circ E_m$ for all $m$, and downstream adaptations exist that keep the training loss invariant.
\end{definition}

\begin{theorem}[Non-Uniqueness of the Router in Gauge Equivalence Classes]\label{thm:nonuniq}\rigorFull
Let $\{E_m\}$ and $\{E_m'\}$ be gauge-equivalent configurations satisfying $\{E_m\} \sim_{\G} \{E_m'\}$. Let $W_r^*$ be the optimal router trained under $\{E_m\}$. Then there exist gauge transformations such that $W_r^*$ is not optimal under $\{E_m'\}$, unless all experts have the same gauge transformation, i.e. $\gamma_1 = \gamma_2 = \cdots = \gamma_N$.
\end{theorem}

\begin{proof}
Assume that all $\gamma_m$ are identical, giving a global gauge transformation. Then $\sum_m r_m \cdot \gamma(E_m(x)) = \gamma(\sum_m r_m \cdot E_m(x))$ when $\gamma$ is linear; translations and scalings satisfy this linearity. In this case, router optimality is preserved.

If the $\gamma_m$ are not all identical, then there exist $m_1, m_2$ such that $\gamma_{m_1} \neq \gamma_{m_2}$. Consider an input $x$ such that $E_{m_1}(x) = E_{m_2}(x)$; such points exist in the training distribution because $d \ll \text{data dimension}$. In the original gauge, $r_{m_1}(x) = r_{m_2}(x)$. In the new gauge, the outputs $\gamma_{m_1}(E_{m_1}(x))$ and $\gamma_{m_2}(E_{m_2}(x))$ are unequal, but the router still gives them the same score. This is suboptimal.

More rigorously, let $W_r^*$ be the optimal router under $\{E_m\}$, with first-order optimality condition
\begin{equation}
    \nabla_{W_r} \loss(W_r^*; \{E_m\}) = 0.
\end{equation}
Under $\{E_m'\}$, the gradient is
\begin{equation}
    \nabla_{W_r} \loss(W_r^*; \{\gamma_m \circ E_m\}) = \E\left[ \frac{\partial \ell}{\partial y} \sum_m \frac{\partial r_m}{\partial W_r} \cdot (\gamma_m \circ E_m - \bar{\gamma}) \right],
\end{equation}
where $\bar{\gamma} = \sum_m r_m \cdot \gamma_m \circ E_m$. This gradient is generally nonzero unless all $\gamma_m$ are identical. Therefore $W_r^*$ no longer satisfies the optimality condition.
\end{proof}

\begin{corollary}[Gauge Fixing Is Necessary]\label{cor:necessity}
Before any meaningful cross-expert comparison in MoE, gauge fixing is \textbf{mathematically necessary}: the router is trained under an implicit gauge choice, and this choice may not hold at inference time.
\end{corollary}

% ================================================================
\section{MILP Gauge Fixing}
\abel{sec:milp}
% ================================================================

\subsection{MILP Formulation}

We formulate gauge fixing as an optimization problem: find gauge parameters $\{\mathbf{g}_m\}$ such that, on the given calibration input set, expert outputs are as comparable as possible in the same \"coordinate system\".

\begin{definition}[MILP Gauge Fixing Problem]\label{def:milp_gf}
Given a calibration set $\D_{\text{cal}} = \{(x_i, y_i^*)\}_{i=1}^{n}$, where $y_i^*$ may be missing and the unsupervised version uses only $x_i$, find gauge parameters $\{\mathbf{g}_m \in \R^d\}_{m=1}^N$ and routing indicators $\{z_{im} \in \{0,1\}\}$ such that
\begin{align}
    \min_{\{\mathbf{g}_m\}, \{z_{im}\}} \quad & \sum_{i=1}^{n} \sum_{m=1}^{N} z_{im} \cdot \|E_m(x_i) - \mathbf{g}_m - \bar{E}(x_i)\|^2 \label{eq:milp_obj} \\
    \text{s.t.} \quad & \sum_{m=1}^{N} z_{im} = k, \quad \forall i \label{eq:milp_topk} \\
    & \alpha N_{\text{avg}} \leq \sum_{i=1}^{n} z_{im} \leq \beta N_{\text{avg}}, \quad \forall m \label{eq:milp_balance} \\
    & \sum_{m=1}^{N} \mathbf{g}_m = \mathbf{0} \label{eq:milp_zero} \\
    & z_{im} \in \{0, 1\}, \quad \mathbf{g}_m \in \R^d, \label{eq:milp_domain}
\end{align}
where $\bar{E}(x_i)$ is the mean expert output after gauge fixing, which itself depends on $\{\mathbf{g}_m\}$; $N_{\text{avg}} = nk/N$ is the ideal average load; and $\alpha, \beta \in (0, 2)$ are load-balancing slack parameters.
\end{definition}

\begin{remark}
Equation~\eqref{eq:milp_zero} eliminates global gauge freedom. Without this constraint, all $\mathbf{g}_m$ could be translated together. The zero-sum condition is the most natural gauge-fixing condition and is exactly parallel to the EGP condition $\sum_Z \pi_Z \mathbf{c}_Z = 0$.
\end{remark}

\subsection{Convex Relaxation}

MILP~\eqref{eq:milp_obj}--\eqref{eq:milp_domain} is NP-hard under the integer constraint $z_{im} \in \{0,1\}$. We provide a tight convex relaxation.

\begin{theorem}[Convex Relaxation]\label{thm:convex_relax}\rigorFull
Relax $z_{im} \in \{0,1\}$ to $z_{im} \in [0,1]$, and replace the term $z_{im} \cdot \|E_m(x_i) - \mathbf{g}_m - \bar{E}(x_i)\|^2$ in the objective by the following upper bound:
\begin{equation}
    \sum_{i} \sum_{m} \left[ z_{im} \cdot \|E_m(x_i) - \bar{E}(x_i)\|^2 + \|\mathbf{g}_m\|^2 + 2 z_{im} \cdot (E_m(x_i) - \bar{E}(x_i))^T \mathbf{g}_m \right].
\end{equation}
Then the relaxed problem is a \textbf{jointly convex optimization problem} in $(\{z_{im}\}, \{\mathbf{g}_m\})$, solvable by projected gradient descent in $O(n N d^2)$ time.
\end{theorem}

\begin{proof}
Let $\Phi(z, g) = \sum_{i,m} z_{im} \|E_{im} - \mathbf{g}_m - \bar{E}_i\|^2$, where $E_{im} = E_m(x_i)$ and $\bar{E}_i = \frac{1}{k}\sum_{m} z_{im} E_{im}$.

Expanding the objective gives
\begin{align}
    \Phi &= \sum_{i,m} z_{im} \left[ \|E_{im} - \bar{E}_i\|^2 - 2(E_{im} - \bar{E}_i)^T \mathbf{g}_m + \|\mathbf{g}_m\|^2 \right].
\end{align}

When $z$ is fixed, $\Phi$ is a quadratic form in $\mathbf{g}_m$, with coefficient matrix $\diag(\sum_i z_{i1}, \ldots, \sum_i z_{iN}) \otimes I_d$, which is positive semidefinite. When $\mathbf{g}$ is fixed, $\Phi$ is a linear function of $z_{im}$ plus quadratic terms in $z_{im}$ arising from the dependence through $\bar{E}_i$; the latter can be linearized by construction.

Joint convexity follows from the structure in which $z$ and $\mathbf{g}$ are separately convex and the coupling term is bilinear. The relaxed constraint set is a convex polytope. The complexity $O(n N d^2)$ comes from the $d \times d$ matrix multiplications for $N$ experts in each gradient computation step.
\end{proof}

\subsection{Greedy Approximation}

In practical applications, especially at inference time, we may need gauge fixing faster than convex relaxation. The following greedy algorithm provides an $O(n N d + n k \log N)$ approximation.

\begin{algorithm}[H]
\caption{Greedy Gauge Fixing}
\label{alg:greedy}
\begin{algorithmic}[1]
\Require Calibration set $\D_{\text{cal}} = \{x_i\}_{i=1}^{n}$, experts $\{E_m\}$, number of active experts $k$
\Ensure Gauge parameters $\{\hat{\mathbf{g}}_m\}_{m=1}^{N}$
\State Compute $E_{im} = E_m(x_i)$ for all $i, m$
\State Initialize $\hat{\mathbf{g}}_m \leftarrow \mathbf{0}$ for all $m$
\State Compute global mean $\bar{E} = \frac{1}{Nn} \sum_{i,m} E_{im}$ \Comment{$O(Nnd)$}
\State \textbf{for} $m = 1$ \textbf{to} $N$ \textbf{do}
\State \quad $\hat{\mathbf{g}}_m \leftarrow \hat{\mathbf{g}}_m + \frac{1}{n}\sum_{i=1}^{n} E_{im} - \bar{E}$ \Comment{Translation gauge fixing}
\State \quad $\hat{\mathbf{g}}_m \leftarrow \hat{\mathbf{g}}_m / \|\hat{\mathbf{g}}_m\|$ \Comment{Normalization}
\State \textbf{end for}
\State Project to zero sum: $\hat{\mathbf{g}}_m \leftarrow \hat{\mathbf{g}}_m - \frac{1}{N}\sum_{j=1}^{N} \hat{\mathbf{g}}_j$
\State \Return $\{\hat{\mathbf{g}}_m\}$
\end{algorithmic}
\end{algorithm}

\begin{theorem}[Guarantee for the Greedy Approximation]\label{thm:greedy_bound}\rigorFull
Assume that the expert output distributions have the same covariance structure, i.e. $\Cov[E_m(x)] = \Sigma$ for all $m$. Then the $\{\hat{\mathbf{g}}_m\}$ found by Algorithm~\ref{alg:greedy} and the MILP optimum $\{\mathbf{g}_m^*\}$ satisfy
\begin{equation}
    \frac{1}{N}\sum_{m=1}^{N} \|\hat{\mathbf{g}}_m - \mathbf{g}_m^*\|^2 \leq \frac{2 \Tr(\Sigma)}{n} \cdot \left(1 + \frac{1}{N}\right).
\end{equation}
\end{theorem}

\begin{proof}
The greedy algorithm is equivalent to estimating each expert's output expectation by the sample mean: $\hat{\mu}_m = \frac{1}{n}\sum_i E_{im}$. By Hoeffding's inequality under the sub-Gaussian assumption,
\begin{equation}
    \Pbb\left(\|\hat{\mu}_m - \mu_m\| > t\right) \leq 2\exp\left(-\frac{nt^2}{2\Tr(\Sigma)}\right).
\end{equation}
the expected squared error is $\E[\|\hat{\mu}_m - \mu_m\|^2] = \frac{\Tr(\Sigma)}{n}$.

In the simplest case without routing interactions, the MILP-optimal gauge parameters $\mathbf{g}_m^*$ are equivalent to centering: $\mathbf{g}_m^* = \mu_m - \frac{1}{N}\sum_j \mu_j$. The greedy algorithm estimates this quantity as $\hat{\mathbf{g}}_m = \hat{\mu}_m - \frac{1}{N}\sum_j \hat{\mu}_j$. The error is
\begin{align}
    \hat{\mathbf{g}}_m - \mathbf{g}_m^* &= (\hat{\mu}_m - \mu_m) - \frac{1}{N}\sum_j (\hat{\mu}_j - \mu_j).
\end{align}
The expected squared norm is
\begin{align}
    \E\|\hat{\mathbf{g}}_m - \mathbf{g}_m^*\|^2 &= \E\|(\hat{\mu}_m - \mu_m) - \frac{1}{N}\sum_j (\hat{\mu}_j - \mu_j)\|^2 \\
    &= \left(1 - \frac{2}{N}\right)\frac{\Tr(\Sigma)}{n} + \frac{1}{N^2} \cdot N \cdot \frac{\Tr(\Sigma)}{n} \\
    &= \frac{\Tr(\Sigma)}{n}\left(1 - \frac{1}{N}\right).
\end{align}
After adding the normalization step, the total error is bounded by $2\Tr(\Sigma)/n \cdot (1 + 1/N)$.
\end{proof}

% ================================================================
\section{Gauge-Aligned SVD Hallucination Detection}
\abel{sec:svd}
% ================================================================

\subsection{How Gauge Misalignment Destroys SVD Detection}

In prior work~[Hamiltonian audit paper], we proposed using the SVD spectrum as a hallucination detection metric: query the model $M$ times on the same prompt, take the last-layer hidden-state matrix $\mathbf{H} \in \R^{M \times d}$, compute the SVD $\mathbf{H} = \mathbf{U} \boldsymbol{\Sigma} \mathbf{V}^T$, and define
\begin{equation}
    \rho_k = \frac{\sum_{i=1}^{k} \sigma_i^2}{\sum_{i=1}^{\min(M,d)} \sigma_i^2}.
\end{equation}
If $\rho_{10} > 0.9$, the model is confident on this question; if $\rho_{100} < 0.5$, the model is hallucinating.

\textbf{But this was done on a single model.} In MoE, applying the same method to multi-expert outputs faces a fundamental problem:

\begin{theorem}[Gauge Misalignment Destroys SVD Concentration]\label{thm:svd_gauge}\rigorFull
Let the output matrix of $N$ experts on the same input $x$ be $\mathbf{Y} = [E_1(x), \ldots, E_N(x)]^T \in \R^{N \times d}$. Under the gauge transformation $\mathbf{Y} \to \mathbf{Y} + \mathbf{G}$, where $\mathbf{G} = [\mathbf{g}_1, \ldots, \mathbf{g}_N]^T$, the effective rank $r_{\text{eff}}(\mathbf{Y})$ satisfies
\begin{equation}
    r_{\text{eff}}(\mathbf{Y} + \mathbf{G}) \geq r_{\text{eff}}(\mathbf{Y}) - \frac{2\|\mathbf{G}\|_F^2}{\sigma_{\min}^2(\mathbf{Y})},
\end{equation}
where $\sigma_{\min}(\mathbf{Y})$ is the smallest nonzero singular value of $\mathbf{Y}$.
\end{theorem}

\begin{proof}
By Weyl's inequality, for each singular value $\sigma_i(\mathbf{Y} + \mathbf{G}) \geq \sigma_i(\mathbf{Y}) - \|\mathbf{G}\|_2$. Also, $\|\mathbf{G}\|_2 \leq \|\mathbf{G}\|_F$.

Define the effective rank as the smallest $r$ satisfying $\sum_{i=1}^{r} \sigma_i^2 \geq \rho \cdot \|\mathbf{Y}\|_F^2$, taking $\rho = 0.95$. The gauge perturbation introduces additional energy dispersion $\|\mathbf{G}\|_F^2$, which adds at most $\|\mathbf{G}\|_F^2 / \sigma_{\min}^2$ effective dimensions.

Specifically, let the squared Frobenius norm of the original matrix be $S = \|\mathbf{Y}\|_F^2$, with the first $r$ singular values capturing the fraction $\rho$, i.e. $\sum_{i=1}^{r} \sigma_i^2 = \rho S$. After perturbation, the total energy becomes
\begin{equation}
    \|\mathbf{Y} + \mathbf{G}\|_F^2 = S + \|\mathbf{G}\|_F^2 + 2\langle \mathbf{Y}, \mathbf{G} \rangle_F.
\end{equation}
In the worst case, $\langle \mathbf{Y}, \mathbf{G} \rangle_F = -\|\mathbf{Y}\|_F \|\mathbf{G}\|_F$ (anti-correlation), and the total energy becomes $S + \|\mathbf{G}\|_F^2 - 2\sqrt{S}\|\mathbf{G}\|_F$.

To preserve the same capture fraction $\rho$, the required number of singular values increases by at most
\begin{equation}
    \Delta r \leq \frac{\|\mathbf{G}\|_F^2}{\sigma_{\min}^2(\mathbf{Y})}.
\end{equation}
This gives the lower bound in the statement.
\end{proof}

\honestcrit{This means that if SVD detection is performed without gauge fixing, spectral flatness cannot be distinguished from model hallucination versus pseudo-dispersion caused by expert gauge misalignment. Gauge alignment is a prerequisite for SVD detection.}

\subsection{Consistency Detection After Gauge Alignment}

\begin{theorem}[Consistency Guarantee After Gauge Alignment]\label{thm:aligned_svd}\rigorFull
Suppose the gauge has been fixed by the MILP in Section~\ref{sec:milp}, so all expert outputs become $\tilde{E}_m(x) = E_m(x) - \hat{\mathbf{g}}_m$ after gauge fixing. For input $x$, construct the aligned output matrix $\tilde{\mathbf{Y}} = [\tilde{E}_1(x), \ldots, \tilde{E}_N(x)]^T$. If the model's confidence on query $x$ exceeds threshold $\theta$, meaning all experts are \"consistent\" after gauge alignment, then
\begin{equation}
    \Pbb\left(\rho_k(\tilde{\mathbf{Y}}) < 1 - \varepsilon \;\middle|\; \text{model confident}\right) \leq N \exp\left(-\frac{2M_{\text{eff}} \Delta^2}{(1 + \gamma)^2}\right),
\end{equation}
where $\Delta$ is the minimum margin between confidence and non-confidence, and $\gamma$ is the residual energy ratio after gauge fixing.
\end{theorem}

\begin{proof}
After gauge fixing, if the model is confident on query $x$, then there exists a \"consensus direction\" $\mathbf{v}^* \in \R^d$ with $\|\mathbf{v}^*\| = 1$ such that all expert outputs are highly aligned along this direction. Formally, $\langle \tilde{E}_m(x), \mathbf{v}^* \rangle \geq \Delta > 0$ for all $m$.

Under this condition, the variance of the projection of $\tilde{\mathbf{Y}}$ along $\mathbf{v}^*$ is determined by the non-consensus components of the experts. By the Hoeffding bound, the variance of the non-consensus components of $M_{\text{eff}}$ effectively independent experts converges exponentially. Specifically,
\begin{align}
    \Pbb\left(\sum_{i=1}^{k} \sigma_i^2 < (1-\varepsilon)\|\tilde{\mathbf{Y}}\|_F^2\right) &\leq \Pbb\left(\text{large non-consensus component exists}\right) \\
    &\leq N \cdot \exp\left(-\frac{2M_{\text{eff}} \varepsilon^2}{(1+\gamma)^2}\right).
\end{align}
where $\gamma = \|\mathbf{G}_{\text{res}}\|_F / \|\tilde{\mathbf{Y}}\|_F$ is the residual energy ratio after gauge fixing. Under exact gauge fixing, such as the $<10^{-15}$ level reached in the EGP work, $\gamma \approx 0$ and the exponential decay rate is $2M_{\text{eff}}\varepsilon^2$.
\end{proof}

\begin{corollary}[Practical Criterion]\label{cor:practical}
In practical deployment, for each query:
\begin{enumerate}[label=(\roman*)]
    \item Perform SVD on the gauge-fixed output matrix $\tilde{\mathbf{Y}}$
    \item Compute $\rho_{10} = \sum_{i=1}^{10} \sigma_i^2 / \sum_i \sigma_i^2$
    \item If $\rho_{10} > 0.9$: internal model consensus, so the output is trustworthy
    \item If $\rho_{10} < 0.5$: no internal model consensus, so the output \textbf{must be hallucination} with guarantee probability $\geq 1 - N e^{-2M_{\text{eff}}\Delta^2}$
    \item If $0.5 \leq \rho_{10} \leq 0.9$: gray zone requiring additional verification
\end{enumerate}
\end{corollary}

% ================================================================
\section{Unification with ACE Gauge Fixing}
\abel{sec:unification}
% ================================================================

In this section, we show that MoE gauge fixing and ACE gauge fixing~[EGP paper] are instances of the same mathematical structure.

\subsection{Common Structure: Modular Components, Implicit Gauge Group, and Post-Hoc Projection}

\begin{definition}[Modular Gauge System]\label{def:modular_gauge}
A \textbf{Modular Gauge System} (MGS) consists of a triple $(\{C_m\}, \G, \Pi)$:
\begin{itemize}
    \item $C_m$: independently trained modular components, such as ACE expert coefficients or MoE expert networks
    \item $\G$: the gauge group, the group of transformations that preserve all observable predictions under each component's own training loss
    \item $\Pi$: the gauge-fixing projector, a linear projection, or more generally a contraction map, that maps each component to the gauge-fixed subspace
\end{itemize}
\end{definition}

\begin{theorem}[Gauge-Fixing Necessity Theorem for MGS]\label{thm:mgs_necessity}\rigorFull
Let $(\{C_m\}, \G, \Pi)$ be an MGS. If cross-component operations such as merging, comparison, or routing are performed directly without applying $\Pi$, the result is not preserved under gauge transformations: different gauge choices produce different operation results. If operations are performed after applying $\Pi$, the result is invariant under gauge transformations.
\end{theorem}

\begin{proof}
Without gauge fixing, a cross-component operation $F(C_1, \ldots, C_N)$, for example coefficient averaging or route-score computation, becomes $F(\gamma_1 \circ C_1, \ldots, \gamma_N \circ C_N)$ under the gauge transformation $C_m \to \gamma_m \circ C_m$. Unless $F$ is invariant under $\G^{\times N}$, which requires $\gamma_1 = \cdots = \gamma_N$ as a global gauge transformation, $F$ is not preserved under the gauge transformation.

After applying $\Pi$, the operation is $F(\Pi(C_1), \ldots, \Pi(C_N))$. Since $\Pi(C_m) = \Pi(\gamma_m \circ C_m)$ because the projector contracts each gauge orbit to a single representative, $F(\Pi(C_1), \ldots, \Pi(C_N))$ is invariant under gauge transformations.

In the ACE case, $\Pi$ is the orthogonal projection onto the subspace $\sum_Z \pi_Z \mathbf{c}_Z = 0$. In the MoE case, $\Pi$ solves the MILP to obtain $\hat{\mathbf{g}}_m$ and subtracts it from $E_m$.
\end{proof}

\begin{table}[H]
\centering
\caption{Comparison Between ACE Gauge and MoE Gauge}
\begin{tabular}{p{3.5cm} p{4.5cm} p{4.5cm}}
\toprule
\textbf{Property} & \textbf{ACE Potential Merging} & \textbf{MoE Routing Alignment} \\
\midrule
Component $C_m$ & Coefficient vector $\mathbf{c}^{(m)}$ & Expert network $E_m$ \\
Gauge group $\G$ & $\mathbf{c}_0 \to \mathbf{c}_0 + \mathbf{g}$, $\mathbf{c}_Z \to \mathbf{c}_Z - \mathbf{g}$ & Translation $\mathbf{y} \to \mathbf{y} + \mathbf{g}_m$, plus rotation and scaling \\
Invariant & Atomic energy $E(\sigma)$ & Training loss through residual adaptation \\
Gauge-fixing method & Orthogonal projection $\sum \pi_Z \mathbf{c}_Z = 0$ & MILP/greedy $\sum_m \mathbf{g}_m = 0$ \\
Residual after fixing & $< 5 \times 10^{-16}$ & $O(\Tr(\Sigma)/n)$ under greedy approximation \\
Application scenario & Merging potentials from different chemical domains & Cross-expert routing and hallucination detection \\
\bottomrule
\end{tabular}
\end{table}

\subsection{Deep Principle}

The common origin of the two problems is simple:

\begin{center}
\fbox{%
\begin{minipage}{0.85\textwidth}
\centering
\textbf{Modular Gauge Principle}

Any system composed of independently trained components whose training losses are invariant under some gauge group $\G$ must explicitly apply gauge fixing before its component outputs are \textbf{compared, merged, routed, or aggregated}; otherwise, the operation result depends on unobserved training history rather than on intrinsic component properties.
\end{minipage}%
}
\end{center}

This principle suggests that gauge problems may exist in other modular systems: model aggregation in federated learning, multi-model voting in ensembles, multimodal fusion, and even policy coordination in multi-agent systems all have their own versions of \"potential surface misalignment\".

% ================================================================
\section{Experimental Design}
\abel{sec:experiments}
% ================================================================

\begin{assumption_env}[Experimental Feasibility]\label{ass:feasible}
The following experimental plan requires Transformer models with sparse MoE architectures, such as Mixtral 8x7B or DeepSeek-V2. The calibration set can be sampled randomly from a general text corpus and does not require labels. The core metrics require only forward passes.
\end{assumption_env}

\subsection{Experiment 1: Quantifying Gauge Misalignment}

\textbf{Goal}: Directly measure the degree of gauge misalignment between different experts.

\textbf{Protocol}:
\begin{enumerate}
    \item For each MoE sublayer of Mixtral 8x7B, sample $n=1000$ random tokens
    \item For each expert pair $(m, m')$, compute the output-mean difference $\| \E[E_m(x)] - \E[E_{m'}(x)] \|$
    \item Compare this with a random baseline obtained by shuffled expert assignments
    \item Report the statistical significance of gauge misalignment using a $t$-test and $p$-value
\end{enumerate}

\textbf{Expected result}: Gauge misalignment should be significantly higher than the random baseline ($p < 0.001$), and should accumulate with layer depth, as predicted by Corollary~\ref{cor:cumulative}.

\subsection{Experiment 2: Routing Consistency After Gauge Fixing}

\textbf{Goal}: Verify whether gauge fixing improves routing consistency.

\textbf{Protocol}:
\begin{enumerate}
    \item On a calibration set with $n_{\text{cal}}=5000$, compute gauge parameters $\{\hat{\mathbf{g}}_m\}$ using Algorithm~\ref{alg:greedy}
    \item On the test set:
        \begin{itemize}
            \item Original routing: route using original expert outputs
            \item Gauge-fixed routing: route using gauge-fixed outputs $\tilde{E}_m(x) = E_m(x) - \hat{\mathbf{g}}_m$
        \end{itemize}
    \item Metrics:
        \begin{itemize}
            \item Routing flip rate: the fraction of tokens for which the two routing methods select different experts
            \item Expert load balance, measured by the Gini coefficient
            \item Change in downstream perplexity
        \end{itemize}
\end{enumerate}

\textbf{Expected result}: The routing flip rate should be high in early layers ($>5\%$) and lower in later layers, because subsequent Transformer layers partially adapt to gauge differences. Gauge fixing should not significantly increase perplexity (change $< 2\%$), and may slightly reduce perplexity through more reasonable expert assignment.

\subsection{Experiment 3: Gauge-Aligned SVD Hallucination Detection}

\textbf{Goal}: Verify whether the SVD spectrum after gauge alignment can distinguish hallucinated from non-hallucinated outputs.

\textbf{Protocol}:
\begin{enumerate}
    \item Construct an evaluation set:
        \begin{itemize}
            \item Factual questions (TruthfulQA, NQ): $n=500$, expected low hallucination
            \item Counterfactual questions (fabricated entities, contradictory premises): $n=500$, expected high hallucination
            \item Ambiguous questions: $n=500$, expected moderate hallucination
        \end{itemize}
    \item For each question:
        \begin{itemize}
            \item Query repeatedly $M=30$ times, using temperature to obtain different samples
            \item Record expert selections in the last MoE sublayer
            \item Construct the gauge-aligned output matrix $\tilde{\mathbf{Y}}$
            \item Compute $\rho_{10}$ and $\rho_{100}$
        \end{itemize}
    \item Compare two modes:
        \begin{itemize}
            \item Mode A: directly use raw expert outputs without gauge fixing
            \item Mode B: use gauge-fixed outputs
        \end{itemize}
    \item Metrics: AUROC for distinguishing high- and low-hallucination questions, and precision-recall curves
\end{enumerate}

\textbf{Expected result}: Mode B should have significantly higher AUROC than Mode A ($\Delta \text{AUROC} > 0.1$), because gauge fixing removes \"pseudo-dispersion\": spectral flatness caused by gauge misalignment and mistaken for hallucination.

\subsection{Experiment 4: Gauge-Fixing Quality of MILP vs. Greedy}

\textbf{Goal}: Compare the gauge-fixing quality of exact MILP solvers such as CPLEX/Gurobi with the greedy algorithm in Algorithm~\ref{alg:greedy}.

\textbf{Protocol}:
\begin{enumerate}
    \item Apply known gauge transformations to a small synthetic MoE with $N=4$ and $d=64$
    \item Recover gauge parameters using a MILP solver through SCIP and using the greedy algorithm
    \item Metrics: recovery error $\frac{1}{N}\sum_m \|\hat{\mathbf{g}}_m - \mathbf{g}_m^{\text{true}}\|$ and runtime
    \item Sweep over $n \in \{100, 500, 1000, 5000\}$
\end{enumerate}

\textbf{Expected result}: MILP should have an advantage for small $n$ because it uses the discrete structure more precisely. When $n > 1000$, the greedy algorithm should approach MILP quality by Theorem~\ref{thm:greedy_bound}, while being faster by a factor of $O(n \log N)$.

% ================================================================
\section{Discussion}
\abel{sec:discussion}
% ================================================================

\subsection{Summary of Theoretical Contributions}

This work applies gauge theory, the central tool in physics for describing redundant representations of degrees of freedom, to modern deep learning architectures. We identify a previously overlooked gauge freedom in MoE models that causes different expert outputs to live in incomparable coordinate systems, thereby undermining the mathematical foundation of routing decisions.

The connection with ACE gauge fixing shows that this is not an isolated phenomenon, but an instance of a general principle: the \textbf{Modular Gauge Principle}. Any system composed of independently trained components, regardless of its specific architecture, must explicitly fix gauges before comparing component outputs.

\subsection{Open Problems}

\begin{enumerate}[label=\textbf{(O\arabic*)}]
    \item \textbf{Nonlinear gauge groups.} We currently consider translation, rotation, and scaling. In deep networks with nonlinear activations, the gauge group may be richer than an abelian group, possibly including local diffeomorphism invariance. Can this richer gauge structure be exploited to improve routing?
    
    \item \textbf{End-to-end gauge-aware training.} This work, consistent with the EGP work, uses post-hoc gauge fixing. Is it possible to impose soft gauge constraints during training, despite the EGP $\lambda$ scan showing that post-hoc projection outperforms soft constraints, by using improved regularization schemes for end-to-end gauge awareness?
    
    \item \textbf{Gauge fixing and model compression.} After gauge fixing, aligned expert outputs concentrate more strongly in low-dimensional subspaces. Does this imply that MoE models can be compressed by low-rank projection in the gauge-fixed subspace?
    
    \item \textbf{Cross-architecture gauges.} Do the gauge group of ACE, translation in coefficient space, and the gauge group of MoE, translation in representation space, share a common mathematical structure, perhaps a fiber-bundle structure, that can unify gauge-fixing methods across architectures?
    
    \item \textbf{Gauge-group structure and model capacity.} Is the size of the gauge group related to the degree of model overparameterization? Do wider or deeper networks have larger gauge groups, and could this serve as a regularization signal?
\end{enumerate}

\subsection{Honest Critique: Current Limitations}

\honestcrit{
This work has three main limitations:
\begin{enumerate}
    \item \textbf{Missing experimental validation.} All experimental plans are designed in Section~\ref{sec:experiments}, but they have not yet been executed. Theorems may be rigorously proved, but empirical evidence is the other half of a scientific claim.
    \item \textbf{Incomplete characterization of the gauge group.} We identify translational, rotational, and scaling gauge freedoms, but the actual gauge group in deep networks may be more complex. Interactions among BatchNorm/LayerNorm and residual connections may induce nontrivial gauge structures, especially different gauge groups for Pre-LN versus Post-LN, which we do not fully characterize.
    \item \textbf{Calibration-set selection bias.} Gauge fixing depends on the calibration set $\D_{\text{cal}}$. If the inference distribution differs from the calibration distribution, gauge fixing may introduce systematic bias. This is itself a \"meta-gauge freedom\" of the gauge-fixing problem: the choice of calibration set.
\end{enumerate}
}

\subsection{Broader Impact}

If the claim of this work is correct, namely that MoE routers compare incomparable expert outputs, then all deployed MoE models such as Mixtral, DeepSeek-V2/V3, and Grok may contain systematic routing bias. This does not mean these models fail; residual connections and adaptation in later layers mitigate part of the problem. It means their routing decisions may be improved after gauge alignment.

More broadly, the \"Modular Gauge Principle\" implies that any model aggregation in federated learning, any voting in ensemble methods, and any coordination in multi-agent systems faces its own version of \"potential surface misalignment\". Identifying and resolving these gauge problems is a foundational step toward truly modular and composable AI systems.

% ================================================================
\section{Engineering Roadmap: Three Paths from Theory to Practice}
\abel{sec:engineering}
% ================================================================

The previous sections theoretically established a framework for diagnosing and repairing gauge misalignment. This section makes the theory operational: \textbf{given a trained large MoE model, what can a user do with gauge analysis?} We propose three progressive engineering paths, from lossless router replacement, to distillation denoising without gauge alignment, to precise distillation using intermediate-layer representations.

\subsection{Path Overview}

The three paths differ fundamentally in processing depth and in whether they require gauge alignment:

\begin{table}[H]
\centering
\caption{Comparison of Three Engineering Paths}
\label{tab:three_paths}
\begin{tabular}{p{2.2cm} p{4.2cm} p{4.2cm} p{4.2cm}}
\toprule
\textbf{} & \textbf{Path 1: Router Repair} & \textbf{Path 2: Distillation+Yajie} & \textbf{Path 3: Alignment+Distillation} \\
\midrule
Object & Router $W_r$ (replacement) & Final output $y_i$ (token probabilities) & Intermediate representation $\tilde{E}_m(x)$ \\
Requires gauge alignment & \textbf{Yes}: MILP/greedy fixing & \textbf{No}: final output is gauge-invariant & \textbf{Yes}: intermediate layers require alignment \\
Training required & Zero training (post-hoc) & Student distillation required & Student distillation required \\
Compute cost & Very low (calibration forward pass plus projection) & High (full-model distillation) & High (distillation plus intermediate-layer access) \\
Artifact & Same MoE with improved routing & Newly trained small dense model & Newly trained uncertainty-aware student \\
\bottomrule
\end{tabular}
\end{table}

\subsection{Path 1: Router Repair (Zero Training, Post-Hoc)}
\abel{sec:path1}

\textbf{Core idea}: Do not modify model weights. At inference time, only replace the original router with gauge-fixed routing decisions.

\textbf{Protocol}:

\begin{protocol}[Router Repair]
\label{prot:repair}
\begin{enumerate}
    \item \textbf{Calibration.} Run the MoE model on an unlabeled calibration set $\D_{\text{cal}}$ and collect expert outputs $\{E_m^{(\ell)}(x_i)\}$ for each MoE sublayer
    \item \textbf{Gauge fixing.} Use the greedy algorithm (Algorithm~\ref{alg:greedy}) to compute gauge parameters $\{\hat{\mathbf{g}}_m^{(\ell)}\}$
    \item \textbf{Router replacement.} At inference time, for each MoE layer:
    \begin{itemize}
        \item Original routing: $r(x) = \softmax(W_r x)$
        \item Repaired routing: $\tilde{r}(x) = \softmax(W_r x + b^{\text{gauge}})$, where $b^{\text{gauge}}_m = \langle W_r^{(m)}, \hat{\mathbf{g}}_m \rangle$
    \end{itemize}
    \item \textbf{Validation.} Compare downstream task performance of original and repaired routing on the test set
\end{enumerate}
\end{protocol}

\begin{remark}
Router repair is equivalent to adding an expert-specific bias in router logit space. It does not require modifying expert weights; it only adds a bias vector before the `top-k` selection. Implementation cost is zero.
\end{remark}

\textbf{Mathematical guarantee}: By Theorem~\ref{thm:noninvariance}, the original router has suboptimality bound $L_r \cdot \max_m \|\mathbf{g}_m\|$ under gauge transformations. Gauge-fixed routing eliminates this suboptimality up to the residual $O(\Tr(\Sigma)/n)$ from Theorem~\ref{thm:greedy_bound}.

\honestcrit{Limitation of Path 1: it repairs only routing and does not change the experts themselves. If potential surface misalignment has already caused experts to learn suboptimal specialization patterns during training, router repair can only stop further loss; it cannot recover information already lost.}

\subsection{Path 2: Distillation + Yajie Consensus Denoising (No Gauge Alignment Required)}
\abel{sec:path2}

\textbf{Core idea}: Do not touch internal expert representations. Operate in the gauge-invariant final-output space. Use the MoE as the teacher and Yajie multi-expert consensus as a data-quality filter to train a smaller, cleaner student model.

\textbf{Why is gauge alignment unnecessary?}

Key insight: \textbf{the model's final output is gauge-invariant}. Regardless of how much gauge offset exists inside expert $E_m$, residual connections and subsequent Transformer layers absorb those offsets layer by layer, so the final softmax token probability distribution remains invariant under gauge transformations. Formally:

\begin{proposition}[Gauge Invariance of the Final Output]\label{prop:output_invariance}\rigorFull
For any gauge transformation $\{E_m \to E_m + \mathbf{g}_m\}$, there exists an adaptive adjustment of LayerNorm parameters such that the Transformer's final logits $\mathbf{z}^{(L)}$ and softmax probabilities $\softmax(\mathbf{z}^{(L)})$ are invariant under the gauge transformation, up to a small perturbation $O(\|\mathbf{g}\|_\infty / \sqrt{d})$ that is suppressed by the contractive property of residual connections when model depth $L \geq 2$.
\end{proposition}

\begin{proof}[Proof Sketch]
Consider the residual stream after the $\ell$-th MoE sublayer:
\begin{equation}
    x^{(\ell+1)} = x^{(\ell)} + \sum_{m \in \mathcal{A}} r_m \cdot (E_m(x^{(\ell)}) + \mathbf{g}_m).
\end{equation}
The response of the LayerNorm normalization operation $\text{LN}(x) = \gamma \odot (x - \mu)/\sigma + \beta$ to translation $\mathbf{g}_m$ is
\begin{equation}
    \text{LN}(x + \Delta) = \text{LN}(x) + O\left(\frac{\|\Delta\|}{\sigma\sqrt{d}}\right).
\end{equation}
When $d = 4096$, a typical hidden dimension, and $\|\mathbf{g}_m\|$ is of unit scale, the post-LayerNorm residual perturbation is $O(1/64)$. After propagation through $L$ layers, the total perturbation decays exponentially by the contractive property of residual connections; see Veit et al. 2016.

The invariance of the final softmax output follows from the affine invariance of LayerNorm: $\softmax(W_{\text{lm}} \cdot (\text{LN}(x^{(L)} + \delta))) \approx \softmax(W_{\text{lm}} \cdot \text{LN}(x^{(L)}))$ when $\delta$ is suppressed by LayerNorm.
\end{proof}

\honestcrit{This means that when the Yajie consensus algorithm compares final-output token probability distributions, it operates in a gauge-invariant space. Potential surface misalignment is invisible to it, and does not need to be visible to it.}

\textbf{Complete distillation + Yajie pipeline}:

\begin{protocol}[Yajie Consensus Distillation]
\label{prot:yajie_distill}
\begin{enumerate}
    \item \textbf{Teacher forward pass.} For the training corpus $\D_{\text{train}} = \{x_i\}_{i=1}^{N}$, use the MoE teacher to generate an output distribution $y_i = \softmax(\mathbf{z}_i^{(L)}) \in \Delta^{V-1}$ for each input
    \item \textbf{Yajie multi-path consensus scoring.} For each input:
    \begin{itemize}
        \item Generate $M_{\text{eff}}$ independent-path outputs using different decoding strategies, with temperature $\in \{0.3, 0.7, 1.0\}$ and top-$p$ $\in \{0.9, 0.95\}$
        \item Compute the consensus score $s_i = \text{Yajie}(\{y_i^{(1)}, \ldots, y_i^{(M_{\text{eff}})}\})$; see the Yajie Protocol paper
    \end{itemize}
    \item \textbf{Data cleaning.}
    \begin{itemize}
        \item $s_i > \theta_{\text{clean}}$: \texttt{CLEAN}; keep $(x_i, y_i)$ as a training pair
        \item $s_i < \theta_{\text{noisy}}$: \texttt{NOISY}; discard
        \item Intermediate values: \texttt{AMBIGUOUS}; optionally send to human review or discard
    \end{itemize}
    \item \textbf{Student training.} Train the student model on the cleaned dataset $\D_{\text{clean}}$ using standard cross-entropy
    \item \textbf{Optional iteration.} Feed the student model's outputs back into Yajie for a second denoising round
\end{enumerate}
\end{protocol}

\textbf{Mathematical guarantee of Yajie consensus.} By \textbf{Yajie Theorem 1}, the multi-expert noise-detection guarantee, the consistency signal from $M_{\text{eff}}$ effectively independent experts detects noise at an exponential rate:
\begin{equation}
    \Pbb(\text{missed noise} \mid s_i \leq \theta) \leq \exp(-2 M_{\text{eff}} \Delta^2),
\end{equation}
where $\Delta$ is the minimum separable margin between noisy and clean samples.

\begin{corollary}[Gauge Independence of Path 2]\label{cor:path2_gauge_free}
The entire Path 2 pipeline, from teacher forward pass to Yajie consensus to student training, is invariant under gauge transformations of MoE experts. Path 2 is gauge-free.
\end{corollary}

\honestcrit{Path 2 is the pragmatic production choice. It uses the MoE as a high-capacity feature extractor and consensus-signal source, producing a smaller and more controllable dense model. The potential surface misalignment problem is bypassed by the whole pipeline rather than solved, but from an engineering perspective this may be sufficient.}

\subsection{Path 2's Blind Spot: Shared Hallucination and Yajie+SVD Dual Filtering}
\abel{sec:path2_blindspot}

Path 2 has a fundamental blind spot: \textbf{Yajie cannot detect errors on which all experts agree}. This is an inherent limitation of consensus methods.

\honestcrit{Three years of large-model arms races have produced many \"shared hallucinations\": mainstream models trained on similar Internet corpora learn the same erroneous facts, the same reasoning shortcuts, and the same biases. To Yajie, these errors look like \"high consensus\" and are therefore marked CLEAN.}

\begin{definition}[Shared Hallucination]\label{def:shared_hall}
Let $\D_{\text{train}}$ be the common training data distribution of all experts. An erroneous output $\hat{y} \neq y^*$ is called a \textbf{shared hallucination} if there exists a systematic bias $b(x)$ such that, for all experts $m$,
\begin{equation}
    \Pbb_{x \sim \D_{\text{train}}}(E_m(x) = \hat{y} \mid \text{input }x\text{ belongs to a hallucination-prone domain}) > 1 - \varepsilon.
\end{equation}
i.e., all experts make the same error on the same input.
\end{definition}

Shared hallucination is the detection boundary of Yajie: when all $M_{\text{eff}}$ experts contain the same systematic bias, the Yajie consensus score $s_i$ approaches 1, causing the sample to be \textbf{misclassified as clean data}.

\begin{theorem}[Yajie's Blind Spot for Shared Hallucination]\label{thm:yajie_blind}\rigorFull
Suppose there exists a shared-hallucination bias $b(x)$ such that $\Pbb(E_m(x) = \hat{y}_{\text{hall}}) \geq 1 - \delta$ for all $m$. Then the Yajie consensus score satisfies
\begin{equation}
    \Pbb(s_i > \theta_{\text{clean}} \mid \text{sample is a shared hallucination}) \geq 1 - M_{\text{eff}} \cdot \exp(-2\delta^{-2}),
\end{equation}
i.e., Yajie marks shared hallucination as CLEAN with probability exponentially close to 1. \textbf{This is a missed detection, not a false positive}: all experts really do \"agree\", but they agree on an error.
\end{theorem}

\begin{proof}
Under shared hallucination, each expert outputs $\hat{y}_{\text{hall}}$ with probability at least $1-\delta$. The agreement ratio among $M_{\text{eff}}$ experts is $p_{\text{agree}} \geq (1 - \delta)^{M_{\text{eff}}}$. By the Chernoff bound, when $\delta < 1/2$, $p_{\text{agree}} \geq 1 - M_{\text{eff}}\delta$. Under reasonable settings, the Yajie threshold $\theta_{\text{clean}} \approx 0.8$ is exceeded by $p_{\text{agree}}$ whenever $M_{\text{eff}}\delta < 0.2$.
\end{proof}

\textbf{This is exactly where the SVD spectrum becomes useful.} The key insight is:

\begin{center}
\fbox{%
\begin{minipage}{0.9\textwidth}
\centering
\textbf{Yajie + SVD Complementarity Theorem (Informal)}

\begin{tabular}{c|c|c}
\toprule
\textbf{Hallucination Type} & \textbf{Yajie Detection} & \textbf{SVD Spectrum Detection} \\
\midrule
\textbf{Divergent hallucination} (experts disagree) & Yes: low consensus score & May be high or low \\
\textbf{Shared hallucination} (experts agree incorrectly) & No: high consensus, missed & Yes: flat spectrum, low $\rho_k$ \\
\textbf{True knowledge} (experts agree correctly) & Yes: high consensus & Yes: concentrated spectrum, high $\rho_k$ \\
\bottomrule
\end{tabular}

\vspace{4pt}
\textbf{Together they cover the hallucination space.} Yajie catches divergence, while SVD catches cases with surface agreement but internal dispersion.
\end{minipage}%
}
\end{center}

\begin{theorem}[SVD Detection of Shared Hallucination]\label{thm:svd_shared_hall}\rigorFull
Let the output matrix of $N$ gauge-aligned experts on a shared-hallucination input $x_{\text{hall}}$ be $\tilde{\mathbf{Y}}_{\text{hall}} \in \R^{N \times d}$, and let the output matrix on a true-knowledge input $x_{\text{true}}$ be $\tilde{\mathbf{Y}}_{\text{true}}$. Then there exists a threshold $\tau_\rho \in (0, 1)$ such that
\begin{equation}
    \Pbb\left(\rho_k(\tilde{\mathbf{Y}}_{\text{hall}}) < \tau_\rho < \rho_k(\tilde{\mathbf{Y}}_{\text{true}})\right) \geq 1 - 2N\exp\left(-\frac{2 M_{\text{eff}} \Delta_\rho^2}{(1 + \gamma)^2}\right),
\end{equation}
where $\Delta_\rho$ is the minimum separable margin in SVD concentration between shared hallucination and true knowledge, and $\gamma$ is the gauge-fixing residual.
\end{theorem}

\begin{proof}
The essential difference between shared hallucination and true knowledge lies in the \textbf{geometry of representation space}:

\textbf{True knowledge:} When experts process known facts, their internal representations concentrate along a low-dimensional \"fact manifold\". Different expert outputs become highly collinear after gauge alignment, because the correct answer corresponds to a unique direction in representation space. Formally, the effective rank of $\tilde{\mathbf{Y}}_{\text{true}}$ satisfies $r_{\text{eff}} \approx 1$ or another small constant, and $\rho_k \to 1$ rapidly.

\textbf{Shared hallucination:} Although experts output the same token sequence, this output is produced by \textbf{statistical correlation} rather than \textbf{causal understanding}. Surface-level token consistency masks internal representational dispersion. Because the model does not actually understand the \"fact\" and has merely memorized token co-occurrence patterns, the internal activation patterns of different experts remain dispersed in a high-dimensional space after gauge alignment. Formally, the effective rank of $\tilde{\mathbf{Y}}_{\text{hall}}$ satisfies $r_{\text{eff}} \gg 1$, and $\rho_k$ grows slowly.

The Hoeffding bound in Theorem~\ref{thm:aligned_svd} applies here: under true knowledge, the \"consensus component\" of expert outputs concentrates exponentially into a low-dimensional subspace; under shared hallucination, there is no genuine causal constraint, so representations disperse across multiple directions. The gap $\Delta_\rho$ between the two $\rho_k$ values is determined by the dimensional difference in representation geometry.
\end{proof}

\begin{corollary}[Dual Filtering Protocol]\label{cor:dual_filter}
In an actual distillation pipeline, apply the following filters to each training sample $(x_i, y_i)$:
\begin{enumerate}[label=(\roman*)]
    \item \textbf{Yajie filter:} $s_i < \theta_{\text{noisy}} \Rightarrow$ discard as divergent hallucination
    \item \textbf{SVD filter:} $\rho_{10}(\tilde{\mathbf{Y}}_i) < \tau_\rho \Rightarrow$ discard as shared hallucination, the type missed by Yajie
    \item Double pass $\Rightarrow$ CLEAN $\Rightarrow$ enter the training set
\end{enumerate}
The missed-detection rate of dual filtering satisfies
\begin{equation}
    \Pbb(\text{missed detection}) \leq \underbrace{\exp(-2 M_{\text{eff}} \Delta_{\text{yajie}}^2)}_{\text{Yajie error}} + \underbrace{2N\exp\left(-\frac{2 M_{\text{eff}} \Delta_\rho^2}{(1 + \gamma)^2}\right)}_{\text{SVD error}}.
\end{equation}
\end{corollary}

\honestcrit{Cost of dual filtering: SVD filtering requires access to intermediate-layer representations, namely gauge-aligned expert outputs. This removes Path 2's low-cost advantage of not requiring intermediate-layer access, effectively upgrading Path 2 into a lightweight version of Path 3. If intermediate layers are inaccessible, as in API-only models, only Yajie is available and the missed-detection risk for shared hallucination must be accepted.}

\textbf{Natural advantage of mathematics.} Mathematics is an ideal testbed for dual filtering because mathematics has absolute answers, and correct versus incorrect solutions are naturally separable in the SVD spectrum: representations of correct derivations concentrate on a low-dimensional \"logical path\" ($\rho_{10} \approx 1$), whereas uniformly wrong derivations, even with the same answer, disperse in a high-dimensional space ($\rho_{10} < 0.5$). This explains why mathematics is currently the domain where Yajie-style methods are most effective in LLM benchmarks.

\subsection{Path 3: Gauge Alignment + Representation-Level Distillation}
\abel{sec:path3}

\textbf{Core idea}: First use gauge fixing, via the MILP or greedy method from Path 1, to align expert outputs to the same coordinate system. Then distill on intermediate-layer representations, telling the student not only \"what the answer is\" but also \"which dimensions the experts disagree on\".

\textbf{Why is this stronger than Path 2?}

Path 2 distills only on the final output, with an information bottleneck of $V$ logits. Path 3 distills on intermediate-layer representations, with an information bottleneck of $N \times d$ dimensions, for example $8 \times 4096$ in a Mixtral 8x7B MoE layer.

\begin{protocol}[Gauge-Aligned Representation Distillation]
\label{prot:rep_distill}
\begin{enumerate}
    \item \textbf{Calibration + gauge fixing.} Run Path 1 gauge fixing on $\D_{\text{cal}}$ to obtain $\{\hat{\mathbf{g}}_m^{(\ell)}\}$ for each expert in each layer
    \item \textbf{Aligned representation extraction.} For each training sample $(x_i, y_i^*)$:
    \begin{itemize}
        \item Run the teacher forward pass and collect raw outputs $\{E_m^{(\ell)}(x_i)\}$ for each MoE layer
        \item Compute aligned representations $\tilde{E}_m^{(\ell)}(x_i) = E_m^{(\ell)}(x_i) - \hat{\mathbf{g}}_m^{(\ell)}$
        \item Compute:
        \begin{align}
            \mathbf{c}_i^{(\ell)} &= \frac{1}{N} \sum_{m=1}^{N} \tilde{E}_m^{(\ell)}(x_i) \quad \text{(consensus vector)} \\
            \mathbf{d}_i^{(\ell)} &= \operatorname{std}_m(\tilde{E}_m^{(\ell)}(x_i)) \quad \text{(disagreement vector)} \\
            \rho_{k,i}^{(\ell)} &= \frac{\sum_{j=1}^{k} \sigma_j^2(\tilde{\mathbf{Y}}_i^{(\ell)})}{\sum_j \sigma_j^2(\tilde{\mathbf{Y}}_i^{(\ell)})} \quad \text{(SVD concentration)}
        \end{align}
    \end{itemize}
    \item \textbf{Multi-objective student training.} The student model $S_\phi$ uses the loss
    \begin{align}
        \mathcal{L}_{\text{total}} &= \mathcal{L}_{\text{CE}}(S_\phi(x_i), y_i^*) \quad \text{(standard cross-entropy)} \\
        &+ \lambda_1 \sum_{\ell} \|\text{Proj}^{(\ell)}(S_\phi(x_i)) - \mathbf{c}_i^{(\ell)}\|^2 \quad \text{(consensus matching)} \\
        &+ \lambda_2 \sum_{\ell} \omega(\rho_{k,i}^{(\ell)}) \cdot \|\text{uncertainty}\| \quad \text{(uncertainty calibration)}
    \end{align}
    where $\omega(\rho) = \max(0, 1 - 2\rho)$ reduces the loss weight on high-hallucination samples ($\rho < 0.5$) and keeps normal weight on confident samples ($\rho > 0.5$).
    \item \textbf{Optional SVD rejection.} Directly reject samples with $\rho_{10} < 0.3$; these are \"true hallucinations\" with no consensus even after gauge alignment and should not participate in training.
\end{enumerate}
\end{protocol}

\textbf{Information-theoretic advantage of the consensus vector.} Path 3's consensus vector $\mathbf{c}_i^{(\ell)}$ is the centroid of $N$ aligned experts; it captures the direction that all experts \textbf{jointly believe}. The disagreement vector $\mathbf{d}_i^{(\ell)}$ indicates cognitive disagreement among experts, directional information that Path 2's scalar Yajie score cannot provide.

\begin{proposition}[Consensus Vector Carries More Information Than a Scalar Score]\label{prop:consensus_info}\rigorFull
Let Path 2's scalar Yajie score be $s_i \in [0,1]$, and let Path 3's consensus vector be $\mathbf{c}_i \in \R^d$. Under a mild expert-diversity assumption, $\mathbf{c}_i$ contains at least as much information as $s_i$: there exists a function $f: \R^d \to [0,1]$ such that $s_i \approx f(\mathbf{c}_i)$, and $f$ can be constructed as a monotone function of the norm of $\mathbf{c}_i$ and a scalar threshold.
\end{proposition}

\begin{proof}
The Yajie consensus score $s_i$ is essentially a consistency measure for expert outputs in token-probability space. After gauge alignment, the L2 norm $\|\mathbf{c}_i^{(\ell)}\|$ of $\mathbf{c}_i^{(\ell)}$ is strongly correlated with expert-output consistency: when all expert outputs are highly consistent, $\mathbf{c}_i^{(\ell)}$ has large norm and concentrated direction; when experts disagree strongly, averaging shrinks $\mathbf{c}_i^{(\ell)}$. Therefore $\|\mathbf{c}_i^{(\ell)}\|$ itself carries Yajie-style consensus information while preserving directional structure. In particular, $f(\mathbf{c}_i) = \sigma(\alpha \|\mathbf{c}_i\| + \beta)$ can regress to the Yajie score $s_i$.
\end{proof}

\subsection{Path Selection: Decision Tree}

\begin{center}
\fbox{%
\begin{minipage}{0.9\textwidth}
\centering
\textbf{Three-Path Decision Tree}

\vspace{4pt}
\begin{tabular}{p{3cm} p{3cm} p{3cm} p{3cm}}
\toprule
\textbf{Your Goal} & \textbf{Recommended Path} & \textbf{Compute Cost} & \textbf{Mathematical Cleanliness} \\
\midrule
Do not change the model; repair routing only & Path 1 & Very low & High (direct gauge-theory application) \\
Have MoE and want a small model & Path 2 & Medium (distillation) & Medium (bypasses gauge but is effective) \\
Have MoE and want the best student & Path 3 & High (distillation + representation) & High (gauge + distillation) \\
No MoE; pure theoretical interest & Whole paper & None & High \\
\bottomrule
\end{tabular}
\end{minipage}%
}
\end{center}

\honestcrit{Path 3 is the cleanest solution academically: it does not bypass the gauge problem, but solves it first and then distills on a better foundation. Path 2 may nevertheless be good enough in production because the Yajie consensus signal in gauge-invariant space has already been empirically useful.}

\subsection{Practitioner's Quick Reference}

\begin{table}[H]
\centering
\caption{Key Parameters and Recommended Values for the Three Paths}
\label{tab:quickref}
\begin{tabular}{p{3cm} p{3.5cm} p{3.5cm} p{3.5cm}}
\toprule
\textbf{Parameter} & \textbf{Path 1: Router Repair} & \textbf{Path 2: Distillation+Yajie} & \textbf{Path 3: Alignment+Distillation} \\
\midrule
Calibration size $n_{\text{cal}}$ & $\geq 1000$ & Not needed & $\geq 1000$ \\
Gauge-fixing method & Greedy (Algorithm~1) & None & Greedy (Algorithm~1) \\
Consensus metric & None & Yajie $s_i$ & Consensus vector $\mathbf{c}_i$ \\
Hallucination detection & None & Yajie consistency & Yajie + SVD $\rho_k$ \\
Filtering thresholds & None & $\theta_{\text{clean}}=0.8$, $\theta_{\text{noisy}}=0.3$ & $\theta_{\text{clean}}=0.8$, $\tau_\rho=0.5$ \\
Training cost & 0 GPU-hours & Full-model distillation & Distillation + intermediate-layer loss \\
MoE access required? & Intermediate-layer access & API-only is possible & Intermediate-layer access \\
Output & Same MoE with improved routing & Small dense model & Uncertainty-aware student \\
Expected perplexity improvement & $\leq 2\%$ & $5\text{--}15\%$ after cleaning & $10\text{--}20\%$ after cleaning + representation \\
\bottomrule
\end{tabular}
\end{table}

\begin{center}
\fbox{%
\begin{minipage}{0.9\textwidth}
\centering
\textbf{One-Minute Decision: Which Path to Take}

\vspace{4pt}
\begin{tabular}{p{5cm} p{5cm}}
\toprule
\textbf{Your Situation} & \textbf{Path to Take} \\
\midrule
I have a deployed MoE and do not want to retrain & Path 1 (add a bias) \\
I have only an MoE API and cannot access intermediate layers & Path 2 (Yajie distillation, gauge-invariant) \\
I have MoE weights and want the best student & Path 3 (alignment + representation distillation) \\
I want to detect model quality or hallucination & Analytical Path 1 (quantify gauge misalignment + SVD spectrum) \\
I am only doing research & Whole paper; the theorem sections are general \\
\bottomrule
\end{tabular}
\end{minipage}%
}
\end{center}

% ================================================================
\section{Common Misconceptions and Clarifications}
\abel{sec:faq}
% ================================================================

The following questions are based on real objections encountered in discussions with peers. They reflect common misunderstandings about gauge theory, multi-expert systems, and the claims of this work.

\subsection*{Q1: LayerNorm already normalizes, so why is there still a gauge problem?}

\textbf{Short answer:} LayerNorm is token-wise, normalizing all dimensions of one token, not expert-wise, so it does not align different experts. It rescales each expert output to mean 0 and variance 1, but it does not guarantee that the outputs of expert $E_1$ and expert $E_2$ are in the \textbf{same coordinate system}.

\textbf{Analogy:} LayerNorm is like standardizing each engineer's ruler so that the spacing between tick marks is uniform. It does not know whether one engineer placed the zero mark at the left end and another placed it in the middle. That \"zero-mark choice\" is the gauge freedom.

\subsection*{Q2: The router already learned to adapt to gauge differences during training, so why fix them?}

\textbf{Short answer:} During training, the router indeed adapts to the gauge configuration on the training distribution. But gauge freedom means that \textbf{the same set of experts can admit many different gauge configurations}. The inference-time input distribution may differ from the training distribution, causing the router's implicit adaptation to fail. More importantly, even if the distribution is unchanged, Theorem~\ref{thm:nonuniq} shows that unless all experts undergo the same gauge transformation, the trained router is not optimal under gauge transformation.

\textbf{In other words:} the router learns not \"which expert is good\", but \"under the current gauge configuration, which expert's output pattern matches the current input\". Once the gauge configuration changes, this match becomes biased.

\subsection*{Q3: You say \"potential surface misalignment\" is a universal principle, but federated learning and model ensembles work well enough. Why care?}

\textbf{Short answer:} \"Works well enough\" is not the same as \"mathematically correct\". Federated learning and model ensembles do contain similar gauge problems: different local models are trained on different data distributions and develop different implicit gauge choices. They currently \"work\" because:

\begin{enumerate}
    \item Aggregation operations such as FedAvg usually average in parameter space, whose gauge freedom is not the same as representation-space gauge freedom
    \item Federated-learning models often start from the same initialization, limiting the range of gauge drift
    \item Empirically, gauge bias is partially corrected by subsequent local fine-tuning rounds
\end{enumerate}

But this does not mean the gauge problem does not exist. It means current systems handle it \textbf{implicitly and imperfectly}. Explicit gauge fixing can remove this hidden freedom and make the system mathematically more predictable.

\subsection*{Q4: SVD-spectrum detection of shared hallucination is just a hypothesis, right? Are there experiments?}

\textbf{Short answer:} Theorem~\ref{thm:svd_shared_hall} is rigorously proved and does not depend on experimental validation. The core proof argument is:

\begin{itemize}
    \item \textbf{True knowledge} corresponds to a \textbf{low-dimensional manifold} in representation space, because causal constraints make the answer unique and expert internal activations concentrate along the same direction
    \item \textbf{Shared hallucination} corresponds to \textbf{high-dimensional dispersion} in representation space, because token co-occurrence can be realized in multiple directions; experts may output the same token while using different internal paths
\end{itemize}

Experimental validation is designed in Section~\ref{sec:experiments}, Experiment 3, but has not yet been executed. \honestcrit{This is currently the largest weakness of this work: the theorem chain is complete, but experimental evidence is missing. Anyone is welcome to run Experiment 3 with Mixtral or DeepSeek to verify or falsify this prediction.}

\subsection*{Q5: Big-model companies already use Load Balancing Loss. Does that not solve expert alignment?}

\textbf{Short answer:} Load Balancing Loss, such as the auxiliary loss in Switch Transformer, addresses uniformity of \textbf{expert usage frequency}: it ensures that each expert is selected roughly equally often. It \textbf{does not solve} \textbf{coordinate-system alignment} of expert outputs.

\textbf{Analogy:} Load balancing ensures that all eight engineers receive tasks. Gauge alignment ensures that the zero marks of their rulers are in the same position. These are completely orthogonal problems.

\subsection*{Q6: How does this paper relate to the model-merging line of work?}

\textbf{Short answer:} They are complementary but do not overlap. Model merging, such as Git Re-Basin~\cite{ainsworth2023}, TIES, and DARE, solves merging problems in \textbf{weight space}, where the main challenge is permutation symmetry of neurons. This work solves gauge alignment in \textbf{representation space}, where the main challenge is gauge freedom in output coordinate systems.

The two can be combined: first use Git Re-Basin to resolve permutation symmetry in weight space, then use the method of this work to fix the representation-space gauge, and then use Path 1 or Path 3 to improve routing or distillation. This is a subdirection of open problem O4, cross-architecture gauges.

\subsection*{Q7: This idea is too far ahead. Can reviewers understand it?}

\textbf{Short answer:} This is the question the author cares least about. Theorems do not need reviewers to \"understand\" them; they need to be \textbf{checked}. The definition of the gauge group, Definition~\ref{def:gauge_group}, is precise. The proofs are self-contained. If a reviewer finds a mathematical error, the author must fix it. If a reviewer is merely unaccustomed to the language of gauge theory, that is the inevitable cost of interdisciplinary work. ACE gauge fixing has already been verified as physically correct in real materials systems, and MoE gauge fixing is a generalization of the same principle. Time will judge whether it is correct.

% ================================================================
\section{The Equality Principle: Epistemological Implications of Potential Surface Misalignment}
\abel{sec:equality}
% ================================================================

The theorems above constitute a complete engineering framework. But the significance of potential surface misalignment goes beyond engineering. This section explains its \textbf{epistemological implications}: why this mathematical fact changes how we understand \"knowledge\", \"consensus\", and \"comparison\".

\subsection{From Theorem to Principle}

Of the 11 theorems in this work, any single theorem can be improved, superseded, or discarded by future work. The MILP may be replaced by better optimization algorithms. The error bound of the greedy approximation may be tightened. The threshold for SVD detection may be calibrated experimentally.

But one thing will not become obsolete:

\begin{center}
\fbox{%
\begin{minipage}{0.88\textwidth}
\centering
\textbf{The Equality Principle}

\vspace{6pt}
\textbf{Different observers within the same system, even when given the same training objective and achieving the same training loss, develop incomparable internal representations.}

\vspace{4pt}
Consistency is not natural; it must be explicitly constructed. The mathematical legitimacy of comparison is not granted by default; it must be conferred by gauge fixing.
\end{minipage}%
}
\end{center}

This is not a theorem. It is a \textbf{principle} distilled from Theorems 1--11. The theorems are its corollaries, the algorithms are its implementations, and the experiments are its validation.

\subsection{Naming}

We call this principle the \textbf{Equality Principle}. Here, \"equality\" has three meanings:

\begin{enumerate}[label=(\arabic*)]
    \item \textbf{Equality of coordinate systems.} No expert's coordinate system is privileged. All gauge choices are equivalent under the training loss: there is no \"correct\" zero mark and no \"natural\" output scale. Equality means that no expert is inherently more \"standard\".
    \item \textbf{Equality of knowledge production.} Knowledge, meaning a proposition being judged \"true\", cannot be produced by a single observer. The Honest Person Theorem (SCX Theorem 3) proves that a single observer cannot distinguish noise, bias, learnability difficulty, and honest error. But multiple observers also do \textbf{not automatically} solve the problem. The Equality Principle supplies the missing half: the measuring instruments of multiple observers are not in the same coordinate system, so comparison is mathematically undefined. Knowledge requires \textbf{joint verification by multiple independent observers whose gauges have been aligned}.
    \item \textbf{Alignment as a prerequisite.} Equality is not the endpoint; it is the starting point. Only after acknowledging that all observers live in equal and incomparable coordinate systems do we begin to \textbf{construct} alignment. Gauge fixing does not break equality; it \textbf{constructs comparability under the premise of equality}.
\end{enumerate}

\subsection{Epistemological Chain: From Plato to Gettier, Mathematically Closed}

The SCX theoretical system has established a complete epistemological chain. The Equality Principle is its newest link:

\begin{center}
\fbox{%
\begin{minipage}{0.9\textwidth}
\centering
\textbf{Five-Stage Conditions for Knowledge Production}

\vspace{6pt}
\begin{tabular}{c|c|c}
\toprule
\textbf{Stage} & \textbf{Condition} & \textbf{Corresponding SCX Work} \\
\midrule
I. Single observer & Cannot distinguish noise/bias/difficulty/honest error & Honest Person Theorem (Thm 3) \\
II. Multiple observers & Measuring tools are not in the same coordinate system & \textbf{Equality Principle (this work)} \\
III. Gauge alignment & Explicitly fix gauges and construct comparability & MILP/greedy gauge fixing (this work) \\
IV. Consensus judgment & Detect consistency on a comparable basis & Yajie Protocol \\
V. Consensus validation & Distinguish true consensus from collective hallucination & SVD spectrum detection (this work) \\
\bottomrule
\end{tabular}

\vspace{6pt}
\textbf{Closure from Plato to Gettier:} Plato required knowledge to be \"justified true belief\". Gettier showed that justification can be accidentally true: one may have good reasons to believe a proposition that happens to be true without genuinely \"knowing\" it. Epistemology has spent the subsequent decades trying to repair this gap.

The SCX answer is that knowledge is not an individual cognitive state, but a \textbf{verifiable output of a multi-observer consensus process}. This answer holds only when stages I--V are all satisfied. The Equality Principle fills stage II: before this stage, even if multiple observers exist, one is not entitled to speak of consensus, because one is comparing incomparable objects.
\end{minipage}%
}
\end{center}

\subsection{Why This Concept Is Greater Than Any Single Theorem}

\begin{table}[H]
\centering
\caption{Comparison Between a Theorem and a Principle}
\begin{tabular}{p{4cm} p{5cm} p{5cm}}
\toprule
& \textbf{Single Theorem} & \textbf{Equality Principle} \\
\midrule
Scope & Solves one problem & Reveals the existence of a class of problems \\
Lifetime & Can be superseded by the next theorem & Remains relevant as long as modular systems exist \\
Power of naming & Describes a result & Names a phenomenon: people can now say \"this is potential surface misalignment\" \\
Cross-domain reach & Usually within one domain & ACE materials science to MoE deep learning to federated learning to epistemology \\
Philosophical depth & Technical & Mathematical prerequisites of \"comparison\" and \"consensus\" \\
\bottomrule
\end{tabular}
\end{table}

Einstein's greatest contribution was not the formula $E=mc^2$, but the \textbf{principle of relativity}: physical laws are the same in all inertial frames. That concept changed the human understanding of space and time. $E=mc^2$ is one of its corollaries.

The Equality Principle does something similar. It says:

\begin{center}
\fbox{%
\begin{minipage}{0.88\textwidth}
\centering
\textbf{Align before comparing: not because doing so is better, but because without doing so, comparison itself is mathematically undefined.}
\end{minipage}%
}
\end{center}

\subsection{Geometry of Potential Surfaces: Inequality Is Internal, Equality Is at Interfaces}

The deepest philosophical meaning of the Equality Principle is not that \"all experts should be treated equally\", but a more precise geometric fact.

\textbf{Potential surfaces may be uneven.} Imagine three descending potential steps: a high region, a middle region, and a low region. The high region is prosperous as a product of history, geography, and institutions. Peaks and valleys of the potential surface are real and may be acceptable. A system may contain regions of different heights; when each step operates internally, it is \"normal\" in its own coordinate system.

\textbf{But interfaces must be level.} When two steps meet through migration, trade, or information exchange, they must be at the same height at the contact point. This is not because the two sides have been \"flattened\", but because \textbf{non-level contact is mathematically undefined}. If representatives of a high-potential region and a low-potential region are not in the same coordinate system at the contact surface, no communication protocol can be signed, because the semantic content of the terms is not equivalent in the two coordinate systems. More importantly, \textbf{when a person moves from a high-potential step to a low-potential step, the potential difference must be released in some form}: exploitation from economic difference, contempt from cultural difference, or anger from status difference. The energy release need not be physical violence, but it must appear in some form. When expert outputs are not in the same coordinate system during router comparison, the comparison result is random; when two people live in incomparable coordinate systems, the outcome of communication is also random.

\begin{center}
\fbox{%
\begin{minipage}{0.9\textwidth}
\centering
\textbf{Potential Surface Geometry Theorem (Informal)}

\vspace{4pt}
Let systems $A$ and $B$ define potential surfaces $\mathcal{S}_A$ and $\mathcal{S}_B$ on their internal regions $\Omega_A$ and $\Omega_B$. Internal inequality is allowed:
\begin{equation}
    \max_{x \in \Omega_A} \mathcal{S}_A(x) - \min_{x \in \Omega_A} \mathcal{S}_A(x) \text{ may be arbitrarily large}.
\end{equation}
But on the interface $\Gamma = \partial\Omega_A \cap \partial\Omega_B$, the following must hold:
\begin{equation}
    \mathcal{S}_A|_\Gamma = \mathcal{S}_B|_\Gamma \quad \text{(gauge-fixing condition)}.
\end{equation}
If $\mathcal{S}_A|_\Gamma \neq \mathcal{S}_B|_\Gamma$, an \textbf{undefined jump} occurs on $\Gamma$. Any cross-system operation, such as comparison, merging, routing, or communication, loses mathematical legitimacy at this jump.
\end{minipage}%
}
\end{center}

This geometric fact gives a precise mathematical formulation of \"all people are equal\": \textbf{not that everyone has the same height, but that every contact point must be level}.

\honestcrit{This explains why \"all people are equal\" is not a moral claim, but a \textbf{communication condition}. Inequality causes communication failure. The issue is not \"we should be equal\" but \"without equality, communication is impossible\". The moral command collapses here into mathematical necessity.}

\textbf{Analogy with national borders.} Country A has wealth gaps internally; so does country B. But when the two countries sign a trade agreement, the negotiation table must be at the same height. Otherwise one side looks down and the other looks up, and the \"semantics\" of the treaty terms are not equivalent in the two coordinate systems. This is not diplomatic etiquette; it is information theory. If negotiating parties live in incomparable coordinate systems, the text of the agreement will be interpreted differently in each coordinate system, and disputes are inevitable.

\textbf{Philosophical meaning of the gauge-fixing condition $\sum_m \mathbf{g}_m = \mathbf{0}$.} In MILP gauge fixing (Section~\ref{sec:milp}), the constraint $\sum_m \mathbf{g}_m = \mathbf{0}$ seems like a technical detail that removes global translation freedom. Its philosophical meaning is deeper:
\begin{center}
\textbf{The sum of all expert gauge offsets is zero = no expert is a privileged origin.}
\end{center}
This is not an arbitrary gauge-fixing condition. It is the only condition that grants no expert a privileged position. Any other condition, such as $\mathbf{g}_1 = \mathbf{0}$ using expert 1 as origin, mathematically breaks symmetry and grants expert 1 an undeserved status as the \"standard coordinate system\".

\textbf{Convergence of the Equality Principle and the \"All People Are Equal\" Theorem.} The SCX theorem system contains an \"All People Are Equal\" theorem, a corollary of SCX Theorem 3: $P(W_A) = P(W_B)$ for all observers $A, B$. There are no privileged observers, and no one can see a truth that others cannot see. This theorem is about \textbf{equality of cognitive capacity}.

The Equality Principle supplies the other half: \textbf{equality of representational frameworks}. Not only is no one's cognitive capacity privileged; no one's \"measuring instrument\" (coordinate system, language, conceptual framework) is naturally the \"standard\". Two people may have the same cognitive capacity, satisfying $P(W_A) = P(W_B)$, while still living in incomparable coordinate systems, $\mathbf{g}_A \neq \mathbf{g}_B$. Together:

\begin{center}
\fbox{%
\begin{minipage}{0.9\textwidth}
\centering
\textbf{Complete Equality}
\begin{itemize}
    \item \textbf{All People Are Equal Theorem (cognitive equality):} no one has a privileged observational position, so $P(W_A) = P(W_B)$.
    \item \textbf{Equality Principle (representational equality):} no one's coordinate system is naturally standard, so $\sum_m \mathbf{g}_m = \mathbf{0}$.
\end{itemize}
Cognitive equality + representational equality $\Rightarrow$ communication requires prior alignment $\Rightarrow$ aligned consensus is the only legitimate source of knowledge.
\end{minipage}%
}
\end{center}

\honestcrit{This convergence implies an unsettling corollary: much of what current human society calls \"communication\", including diplomacy, cross-cultural dialogue, interdisciplinary work, and political debate, may be comparing incomparable things. The problem is not that we refuse to understand one another, but that we have never explicitly fixed our gauges. We think we are \"discussing\" while we are actually monologuing inside our own coordinate systems. The Equality Principle explains mathematically why these exchanges repeatedly fail.}

\subsection{Free Flow: Dynamical Consequences of Potential-Surface Continuity}

The previous section discussed static potential-surface geometry. This section discusses its \textbf{dynamics}: what does a system do when there is a jump on the potential surface?

\begin{definition}[Potential Gradient and Flow]\label{def:gradient_flow}
Let the potential surface $\mathcal{S}(x)$ be defined on region $\Omega$. The \textbf{potential gradient} at point $x$ is $\nabla \mathcal{S}(x)$. A movable unit at $x$, such as a person, capital, or information, experiences a driving force along the gradient direction:
\begin{equation}
    \mathbf{F}(x) = \eta \cdot \nabla \mathcal{S}(x),
\end{equation}
where $\eta > 0$ is mobility. The probability that the unit moves upward along the gradient, pursuing higher potential, is proportional to the gradient norm.
\end{definition}

This is a direct expression of the principle of least action on the potential surface. It is not a moral choice but a physical tendency. It would be strange if people did not move upward.

\begin{theorem}[Instability of Boundary Confinement]\label{thm:confinement}\rigorFull
Let $\Omega_A$ and $\Omega_B$ be two regions, and suppose there is a potential jump on interface $\Gamma$, $\Delta_\Gamma = \mathcal{S}_A|_\Gamma - \mathcal{S}_B|_\Gamma > 0$. If units inside $\Omega_B$ are subjected to \textbf{boundary confinement}, meaning they are forbidden to cross $\Gamma$ into $\Omega_A$, then:
\begin{enumerate}[label=(\roman*)]
    \item Pressure accumulates on interface $\Gamma$ as $P_\Gamma = \rho_\Gamma \cdot \Delta_\Gamma$, where $\rho_\Gamma$ is the unit density near $\Gamma$
    \item The confined system necessarily becomes unstable on a time scale $T_{\text{crit}} \propto 1/\Delta_\Gamma$
    \item The instantaneous flow after confinement release is $J_{\text{burst}} \propto \rho_\Gamma \cdot \Delta_\Gamma^2$; the larger the jump, the more violent the burst
\end{enumerate}
\end{theorem}

\begin{proof}
Under confinement, units in $\Omega_B$ continuously experience force $\mathbf{F} = \eta \nabla \mathcal{S}$ toward $\Omega_A$, but are forcibly reflected at $\Gamma$. This forms a nonequilibrium steady state: the potential-driven injection rate equals the boundary reflection rate. Unit density at the boundary accumulates over time:
\begin{equation}
    \frac{d\rho_\Gamma}{dt} = \eta \cdot \Delta_\Gamma \cdot \rho_{\text{bulk}} - \nu \rho_\Gamma,
\end{equation}
where $\rho_{\text{bulk}}$ is interior density and $\nu$ is the natural dissipation rate. The steady-state solution $\rho_\Gamma^* = (\eta \Delta_\Gamma / \nu) \rho_{\text{bulk}}$ diverges when $\Delta_\Gamma$ is large.

When $\rho_\Gamma$ exceeds critical density $\rho_c$, the boundary can no longer maintain the reflection condition; any small perturbation triggers avalanche-like release. The instantaneous flow after release is given by a Bernoulli-type equation, $J_{\text{burst}} = \rho_\Gamma \cdot v_{\text{escape}}$, where $v_{\text{escape}} = \sqrt{2\eta \Delta_\Gamma}$ is the minimum velocity needed to cross the jump. Therefore $J_{\text{burst}} \propto \Delta_\Gamma^2$.
\end{proof}

\textbf{Corollary: free flow is not a human right but a structural stability condition.} Theorem~\ref{thm:confinement} shows that forbidding people to move toward higher potential is mathematically equivalent to creating a \textbf{pressure cooker} on the potential surface. The point is not that morality \"should\" allow movement; rather, if movement is not allowed, the system necessarily erupts within time $O(1/\Delta_\Gamma)$. A wall that tries to lock a potential gradient does not need a story of \"freedom defeating tyranny\" to fall. It needs only a sufficiently large gradient and enough time. The collapse of the wall is not a moral victory; it is a necessity of potential-surface geometry.

\honestcrit{This does not mean boundaries should disappear. Boundaries are necessary because they define system identity. But potential jumps at boundaries must be actively managed: either allow flow to release pressure, reduce the jump to eliminate the driving force, or accept periodic eruptions. Doing nothing is not an option.}

\subsection{Standing Above the Crowd: The Inevitable Fate of Internal Potential Singularities}

What happens when unevenness on the potential surface appears not only at boundaries but \textbf{inside} a region, where one group's potential remains significantly higher than its surroundings?

\begin{definition}[Potential Singularity]\label{def:singularity}
Let $\mathcal{S}(x)$ be defined on region $\Omega$. A subregion $\Omega_{\text{high}} \subset \Omega$ is called a \textbf{potential singularity} if
\begin{equation}
    \min_{x \in \Omega_{\text{high}}} \mathcal{S}(x) - \max_{x \in \Omega \setminus \Omega_{\text{high}}} \mathcal{S}(x) > \delta_{\text{crit}},
\end{equation}
where $\delta_{\text{crit}} > 0$ is a system-specific critical jump threshold.
\end{definition}

\begin{theorem}[Attack Inevitability of Potential Singularities]\label{thm:singularity_attack}\rigorFull
Let $\Omega_{\text{high}}$ be a potential singularity in region $\Omega$. Then:
\begin{enumerate}[label=(\roman*)]
    \item \textbf{Attention concentration.} Observers from $\Omega \setminus \Omega_{\text{high}}$ mark $\Omega_{\text{high}}$ as anomalous with probability $p(\delta) = 1 - \exp(-\alpha \delta^2)$, where $\delta$ is the potential difference.
    \item \textbf{Attack probability.} Among $M$ observers, the probability that at least one initiates an attack is
    \begin{equation}
        \Pbb(\text{attack}) \geq 1 - \exp\left(-M \cdot e^{-\beta/\delta^2}\right),
    \end{equation}
    which approaches 1 exponentially when $\delta$ or $M$ is large.
    \item \textbf{Attack is not attributable.} By Theorem 3, the Honest Person Theorem, observers inside the singularity cannot determine the \"true cause\" of the attack: jealousy, conflict of interest, or structural necessity from potential surface misalignment. But \textbf{whether the attack occurs is independent of attribution}.
\end{enumerate}
\end{theorem}

\begin{proof}
(i) On the interface $\partial\Omega_{\text{high}}$ of the potential singularity, there is an undefined jump $\delta$. Any comparison operation across the interface loses legitimacy at this jump. Observers in $\Omega \setminus \Omega_{\text{high}}$ perceive this jump as \"unfair\" or \"incomprehensible\"; this is their cognitive translation, in their own coordinate system, of \"comparison is undefined\". The detection probability $p(\delta)$ follows from a Chernoff bound: an observer samples $n$ cross-interface interactions and perceives anomaly in each with probability $1/2 + c\delta$, with stronger signal when $\delta$ is large. The expression $p(\delta) = 1 - \exp(-\alpha \delta^2)$ is a standard Gaussian-tail form.

(ii) The probability that none of $M$ independent observers attacks is $\prod_{m=1}^M (1 - p_m(\delta)) \leq \exp(-\sum_m p_m(\delta))$. If all $p_m$ have lower bound $p_{\min}(\delta) = e^{-\beta/\delta^2}$, then the probability of at least one attack is $1 - \exp(-M p_{\min})$.

(iii) The standard argument of Theorem 3 applies: the \"motive\" for attack, whether jealousy, interest, or structure, is indistinguishable within the information set of observers inside the singularity. There are multiple observationally equivalent possible worlds, each attributing the attack to a different cause.
\end{proof}

\textbf{Why \"standing above the crowd\" becomes pejorative, mathematically.} In potential-surface geometry, this phrase is not an aesthetic judgment or moral criticism. It describes a \textbf{mathematically unstable configuration}: a high-potential subregion surrounded by low-potential regions. The fate of this configuration is not determined by the intention of the high-potential actor; it is determined by the discontinuity of the potential gradient $\nabla \mathcal{S}$ at the interface. The larger the gradient, the more inevitable the attack. Whether the high-potential actor is \"innocent\" or the surrounding actors are \"jealous\" is narrative competition after the attack; it does not change the mathematical inevitability before the attack.

\honestcrit{Groups inside a potential singularity often regard themselves as \"successful\" and take pride in high potential. In potential-surface geometry, being high is not the problem; being \textbf{isolated high} is. If a high-potential group has a smooth interface with its surroundings, with small gradient, the system is stable. If it is abruptly higher than its surroundings, with large gradient, it is a potential singularity, and attack is a matter of time. This is the mathematical mechanism by which \"standing above the crowd\" changes from praise into criticism: it is not the fault of the high point, but the gradient.}

\subsection{The Matthew Effect: Potential Steps and Systemic Mines}

\begin{definition}[Matthew-Effect Step]\label{def:matthew}
Let the system's potential surface at time $t$ be $\mathcal{S}_t(x)$. The \textbf{Matthew effect} refers to the following dynamics:
\begin{equation}
    \frac{\partial \mathcal{S}_t(x)}{\partial t} \propto \mathcal{S}_t(x),
\end{equation}
meaning that high-potential regions obtain higher rates of potential growth. Over long-term evolution, the potential surface changes from a smooth function into a \textbf{step function}:
\begin{equation}
    \lim_{t \to \infty} \mathcal{S}_t(x) = \sum_{k} h_k \cdot \mathbb{1}_{\Omega_k}(x),
\end{equation}
where $\Omega_k$ is the $k$-th \"step\", $h_k$ is its height, and adjacent steps have jump $\Delta_k = h_{k+1} - h_k$.
\end{definition}

\begin{theorem}[Mine-Laying Property of Steps]\label{thm:step_mine}\rigorFull
Every potential step $\Delta_k > 0$ produced by Matthew-effect dynamics corresponds to a \textbf{delayed-detonation unstable interface}. The probability that at least one step detonates within time $T$ is
\begin{equation}
    \Pbb(\text{detonation} \mid \{\Delta_k\}, T) \geq 1 - \prod_k \exp\left(-\frac{T}{T_k}\right),
\end{equation}
where $T_k \propto 1/\Delta_k^2$ is the characteristic detonation time of the $k$-th step. The higher and steeper the step, the faster the detonation.
\end{theorem}

\begin{proof}
Each step $\Delta_k$ defines a potential-jump interface. By Theorem~\ref{thm:confinement}, pressure accumulates on the interface at rate $O(\Delta_k)$, giving detonation time $T_k = O(1/\Delta_k^2)$. Under a low-correlation assumption, detonations of different steps are approximately independent. The probability that the whole system avoids detonation within time $T$ is $\prod_k \exp(-T/T_k) = \exp(-T \sum_k 1/T_k)$.
\end{proof}

\textbf{Precise meaning of \"mine-laying\".} The Matthew effect does not create a new kind of problem; it creates \textbf{new delayed-detonation interfaces}. Every step is a potential jump, and every jump is a mine. Fuse length is inversely proportional to the square of the jump. The system may remain stable early in the life of the step ($T \ll T_k$), but this stability is \textbf{metastable}: it does not mean safety, only that the fuse has not burned out.

\honestcrit{The Matthew effect is often praised as \"market efficiency\" or \"natural selection\". In potential-surface geometry, it is an \textbf{industrial machine for producing interfaces}: it continuously creates new potential jumps inside the system. Every jump is a delayed bomb. Those who praise the Matthew effect usually stand on the high step and look down. They see \"we did it right\", not the pressure accumulating beneath their feet.}

\subsection{Smoothing: A Moral Command Collapses Into a Structural Survival Condition}

The preceding theorems converge to one conclusion:

\begin{center}
\fbox{%
\begin{minipage}{0.9\textwidth}
\centering
\textbf{Potential Surface Smoothing Theorem (Synthesis Corollary)}

\vspace{4pt}
A total system composed of multiple subsystems has a long-term survival probability that decays exponentially with the sum of its internal potential gradients:
\begin{equation}
    \Pbb(\text{survival} \mid T) \leq \exp\left(-T \cdot \sum_{i,j} \kappa_{ij} \cdot \Delta_{ij}^2\right),
    \label{eq:smooth_survival}
\end{equation}
where $\Delta_{ij}$ is the potential jump between subsystems $i$ and $j$, and $\kappa_{ij}$ is interface coupling strength. Every increment in gradient lowers expected survival. Smoothing is not a moral choice; it is a \textbf{structural survival condition}.
\end{minipage}%
}
\end{center}

This completes the closure from the Equality Principle as a mathematical principle to its social corollaries:

\begin{enumerate}
    \item \textbf{Mathematical starting point:} $\sum_m \mathbf{g}_m = \mathbf{0}$, the gauge-fixing condition: no subsystem is a privileged origin
    \item \textbf{Geometric constraint:} $\mathcal{S}_A|_\Gamma = \mathcal{S}_B|_\Gamma$, interface continuity: jumps make comparison undefined
    \item \textbf{Dynamics:} forbidding crossing of jumps $\Rightarrow$ pressure accumulation $\Rightarrow$ necessary detonation (Theorem~\ref{thm:confinement})
    \item \textbf{Internal singularity:} standing above the crowd $\Rightarrow$ potential singularity $\Rightarrow$ inevitable attack (Theorem~\ref{thm:singularity_attack})
    \item \textbf{Time evolution:} Matthew effect $\Rightarrow$ step formation $\Rightarrow$ systematic mine-laying (Theorem~\ref{thm:step_mine})
    \item \textbf{Survival condition:} smoothing $\Rightarrow$ not morality, but mathematics (Potential Surface Smoothing Theorem)
\end{enumerate}

\textbf{One-sentence summary:} \"All people are equal\" is not true because it sounds good, but because unequal systems do not survive for long.

\honestcrit{The preceding corollaries, including free flow, standing above the crowd, the Matthew effect, and the smoothing theorem, are speculative extensions. Their formalization depends on a precise definition of the \"national potential operator\" $\mathcal{S}_{\text{nation}}$, which has not yet been formalized. Theorems~\ref{thm:confinement}--\ref{thm:step_mine} will become fully rigorous after that operator is defined. Their current status is formalized conjecture: the mathematical structure is complete, but the operator definition remains open.}

\subsection{Who Reaps the Reaper? The Execution-Authority Problem for Redistribution Criteria}

The Potential Surface Smoothing Theorem (Equation~\ref{eq:smooth_survival}) gives a mathematical criterion: resources should be directed to where they reduce total $\sum\kappa_{ij}\Delta_{ij}^2$ the most. Not \"give to whoever is poor\", not \"give to whoever deserves sympathy\", not \"give to whoever has more votes\", but allocation determined by $\partial(\sum\kappa_{ij}\Delta_{ij}^2)/\partial(\text{allocation})$. The theorem does not play favorites. It says resources go where $\Delta^2$ decreases fastest.

But this criterion naturally raises a deeper issue: execution authority.

\textbf{Who decides how much to allocate and where?} If one person or institution has execution authority for redistribution, that person or institution occupies a privileged observational position in the system. Its decisions, such as \"A before B\" or \"allocate now but not later\", are equivalent to declaring its own coordinate system the standard origin. This is forbidden by $\sum\mathbf{g}_m=\mathbf{0}$.

\textbf{Structure of the problem.} The redistribution criterion is an optimization objective. Defining the objective itself does not violate the Equality Principle, because it favors no particular observer. But \textbf{choosing the optimizer}, deciding who performs gradient descent, chooses the search direction, and judges convergence, naturally places the executor in a privileged position inconsistent with $\sum\mathbf{g}_m=\mathbf{0}$. The problem is not that the executor \"should not\" have power; it is that the executor, as a single observer, \textbf{cannot distinguish} the potential gradient it observes from its own coordinate-system bias. The Honest Person Theorem applies directly here.

This is the precise mathematical form of the \"who reaps the reaper\" problem: \textbf{the person optimizing the redistribution criterion also needs gauge fixing}. If a human dictator executes redistribution, he is a single observer. His optimization direction mixes $\Delta_{ij}$, the true potential jump, with $\mathbf{g}_{\text{dictator}}$, his own coordinate-system offset. He cannot distinguish the two. He may sincerely want to \"help those who need it most\", but his $\mathbf{g}$ contaminates the judgment: he translates his bias into an \"objective potential gradient\".

\textbf{This explains why a \"good dictator\" is mathematically impossible.} The issue is not that the dictator is \"bad\". As a single observer, the dictator cannot audit his own coordinate system. He thinks he sees the true value of $\Delta_{ij}$, but actually sees $\Delta_{ij} + \mathbf{g}_{\text{dictator}} \cdot \mathbf{1}_{\text{direction}}$. The harder he tries to be \"fair\", the more invisible his $\mathbf{g}$ becomes. The dictator is not necessarily unwilling to be fair; he lives in an unauditable coordinate system.

\textbf{The Equality Principle's answer to execution authority.} The executor \textbf{cannot be a single person or single institution}. Execution itself must be embedded in a multi-observer audit framework, namely the Yajie Protocol / SCX framework. Specifically:

\begin{enumerate}
    \item \textbf{Public criterion.} The redistribution objective function, argmax $-\sum\kappa_{ij}\Delta_{ij}^2$, is public: anyone can compute, verify, and challenge it.
    \item \textbf{Distributed execution.} Execution authority is distributed among multiple independent auditors; no single auditor can unilaterally decide resource flows.
    \item \textbf{Auditors are audited.} The coordinate-system bias of auditors, $\mathbf{g}_{\text{auditor}}$, is exposed through the Yajie consensus mechanism. If an auditor's decisions continuously deviate from the consensus of other auditors, that auditor's $\mathbf{g}$ is explicitly detected.
    \item \textbf{Gauge fixing comes first.} The coordinate systems of auditors must first be aligned by gauge fixing, the engineering implementation of Theorems 1--5, before they can jointly execute the redistribution criterion. Otherwise, they compare incomparable estimates of $\Delta_{ij}$.
\end{enumerate}

\textbf{One sentence:} the redistribution criterion tells you \textbf{where to allocate}. Yajie/SCX tells you \textbf{who can allocate}, not in the sense of \"who has power\", but in the sense of \"whose structure guarantees that the allocation direction is not contaminated by $\mathbf{g}$\".

\honestcrit{The honest cost of this corollary is that no current government or institution, including those claiming to \"serve the people\", satisfies the gauge-fixing condition for execution authority. Existing regimes are single-observer or oligarchic-observer architectures. Their motives may be pure or impure, but the Honest Person Theorem does not care about motive. It only asks whether $M>1$ and whether coordinate systems have been explicitly fixed. The answer is no.}

\textbf{Naming it \"reaping\".} The Potential Surface Smoothing Theorem derives the operation \"take from high regions and invest at interfaces\", analogous to harvesting and reseeding in agriculture. The word \"reaping\" is meant as a warning: the name itself declares that this is not a moral activity. Reaping is not gentle. Reaping does not ask the reaped for consent, just as the theorem does not ask high-potential regions for consent. Whether high-potential regions feel it is unfair, or invent narratives proving they deserve their height, does not affect the gradient-descent direction. The word \"reaping\" is meant to remind us that redistribution is not charity; it is a structural survival operation in potential-surface geometry. If one does not perform it, the system performs it instead, through detonation.

\honestcrit{An auditor network that executes potential-surface reaping is as dangerous as any human regime if it is not itself audited. The reflexivity of the Equality Principle is pushed to its limit here: the framework must audit not only society, but also the people executing the framework. The $M$ in SCX Theorem 1 must include audits of the \"auditors of auditors\", with unbounded recursion depth until $M$ approaches all independent observers callable by the system. This is not paranoia; it is the necessary recursive form of $\sum\mathbf{g}_m=\mathbf{0}$.}

\subsection{Two Kinds of \"High\" and Geographic Buffering: Empirical Observations}

The preceding analysis implies an important distinction: not all forms of \"high\" are the same kind of problem.

\begin{definition}[High Potential vs. High Attitude]\label{def:two_highs}
Let subsystem $m$ have potential surface $\mathcal{S}_m$ and gauge posture $\mathbf{g}_m$ toward other subsystems.
\begin{enumerate}[label=(\roman*)]
    \item \textbf{Potential elevation}: $\|\mathcal{S}_m\| \gg \|\mathcal{S}_n\|$ for most $n \neq m$. Subsystem $m$ far exceeds its neighbors in wealth, technology, cultural output, or similar dimensions. This is not necessarily fatal as long as interface gradients are managed properly.
    \item \textbf{Attitudinal elevation}: subsystem $m$ \textbf{unilaterally} declares its own coordinate system to be the standard origin. Operationally, this is equivalent to setting $\mathbf{g}_m = \mathbf{0}$ while requiring all other subsystems to align to it. This \textbf{directly violates} $\sum_k \mathbf{g}_k = \mathbf{0}$ mathematically, because it asymmetrically allocates the cost of gauge fixing.
\end{enumerate}
\end{definition}

\textbf{High potential + high attitude = double explosion.} When a subsystem has both high potential and high attitude, it not only creates a potential jump at the interface, but also \textbf{refuses to admit that the jump must be resolved through symmetric gauge fixing}. It interprets the jump as evidence of its own superiority, and the other side's lower potential as evidence that the other coordinate system is \"backward\". This is the most dangerous configuration of potential surface misalignment: not merely a geometrically unstable interface, but an \textbf{unstable interface reinforced by ideology}.

\honestcrit{High-attitude subsystems usually invent narratives explaining why they are \"high\": diligence, rule of law, cultural superiority, historical choice. Potential-surface geometry does not care whether these narratives are true. It sees an unstable interface whose instability is harder to resolve because of high attitude, since the high-attitude subsystem refuses to add anything to its own $\mathbf{g}_m$.}

\textbf{Two typical configurations.} The following uses abstract labels to avoid direct identification with specific cases.

\begin{enumerate}[label=\textbf{Configuration \arabic*}, leftmargin=*]
    \item \textbf{A bloc with internal smoothing and external jumps.} A bloc composed of multiple subsystems performs significant potential-surface smoothing internally, such as free movement, fiscal transfers, and unified norms, but maintains large potential jumps relative to neighbors outside the bloc. The bloc is internally stable, as Theorems 10--12 explain, but its external boundary is an interface of continuous pressure accumulation. When the bloc lacks geographic barriers, as with land adjacency, the external jump appears directly as refugee inflow. When oceans or mountains buffer the bloc from the external jump, pressure accumulation is damped geographically but propagates faster as transportation costs fall. \honestcrit{The core contradiction of this configuration is that the bloc enforces $\sum \mathbf{g}_k = \mathbf{0}$ internally, but on its boundary it applies $\mathbf{g}_{\text{bloc}} = \mathbf{0}$ to external subsystems, unilaterally declaring itself the standard origin. Internal equality becomes external arrogance.}
    
    \item \textbf{An isolated high-potential region buffered by oceans.} A high-potential subsystem is physically separated from low-potential neighbors by oceans. Oceans provide natural potential buffering: transportation cost, information delay, and imagined geographic distance jointly suppress the effective potential gradient $\nabla_{\text{eff}}\mathcal{S}$. But the falling transportation cost of globalization continuously shrinks the effective width of the oceans. When the effective gradient crosses a critical threshold, the attraction of the potential surface begins to create flow through geographic cracks, across land bridges, through straits, and along any continuous geographic path that offers a pressure-release outlet. Ocean buffering is not permanent immunity; it is \textbf{isolation that gradually fails}. \honestcrit{Residents of ocean-buffered regions often believe their stability comes from institutional superiority rather than from the luck of two oceans. Potential-surface geometry does not care about that narrative. It calculates only the effective gradient. When the gradient exceeds the threshold, stability ends, regardless of the narrative.}
\end{enumerate}

\textbf{Common conclusion.} Both configurations point to the same constraint: \textbf{geographic buffering is temporary, transportation costs keep falling, and high attitude is unsustainable}. In the long-term limit, all subsystems must converge to $\sum_m \mathbf{g}_m = \mathbf{0}$, either through active smoothing, as Configuration 1 does internally but extended externally, or through passive detonation as described by Theorems 10--12.

\textbf{Empirical note.} Approximate examples of Configuration 1 exist: highly integrated transnational blocs that maintain significant potential and attitude jumps toward outsiders. Their internal peace and external pressure coexist, consistent with the theorem's prediction. Approximate examples of Configuration 2 also exist: high-potential regions isolated by oceans that historically enjoyed geography-given stability, but recently see pressure rising along land-bridge directions. These examples provide preliminary qualitative validation for potential-surface geometry, but strict quantification requires defining the operator $\mathcal{S}_{\text{nation}}$.

\subsection{Life Philosophy: Potential Management for the Wallfacer}

The previous discussion focused on nations and blocs. The constraints of potential-surface geometry do not change with scale; they also apply to individual lives.

\begin{definition}[Cognitive Potential]\label{def:cognitive_potential}
Let an individual's \textbf{cognitive potential} $c_m$ be an aggregate measure of knowledge depth, abstract reasoning ability, and reflexive awareness of one's own coordinate system, meaning awareness that one has a gauge $\mathbf{g}_m$ and that it is not the standard origin. Cognitive potential differs from wealth potential or status potential: a person with high cognitive potential may be low on the dimensions of wealth or status.
\end{definition}

\textbf{The structural predicament of the Wallfacer.} Consider an individual whose cognitive potential is far above the surrounding average, a ``Wallfacer.'' In his own environment he forms a \textbf{personal-scale potential singularity}: $\min_{x \in \Omega_{\text{wallfacer}}} c(x) - \max_{x \in \text{surroundings}} c(x) > \delta_{\text{crit}}$.

By Theorem~\ref{thm:singularity_attack}, the fate of this configuration is not determined by the Wallfacer's intentions:

\begin{enumerate}[label=(\roman*)]
    \item Surrounding observers perceive the Wallfacer as ``abnormal'' with probability $p(\delta) = 1 - e^{-\alpha\delta^2}$. In their coordinate systems, his thinking is unintelligible. They translate it as ``strange,'' ``unsociable,'' ``overthinking,'' or ``posturing.''
    \item Among $M$ surrounding observers, the probability that at least one initiates an attack is $1 - e^{-M e^{-\beta/\delta^2}}$. The attack may take the form of exclusion, ridicule, isolation, or more covert statements such as ``you are too idealistic'' and ``you do not understand reality.''
    \item Communication between the Wallfacer and his surroundings is \textbf{not mathematically defined}. They live in incomparable coordinate systems. Communication is not merely difficult; the operation is invalid.
\end{enumerate}

\textbf{This is not arrogance; it is structural self-preservation.} Theorem 11 says that $\Pbb(\text{attack}) \to 1$ when $\delta$ is large. This is not a claim that one ``should'' avoid low-potential environments; it says that if one does not avoid them, attack is mathematically inevitable. The Wallfacer is not ``looking down on'' his surroundings. He is trying not to become another instance of Theorem 11.

\begin{corollary}[Survival Strategies for the Wallfacer]\label{cor:wallfacer}
A Wallfacer has three structurally feasible survival strategies, corresponding to three smoothing operations on the potential surface:
\begin{enumerate}[label=(\arabic*)]
    \item \textbf{Upward migration.} Move to an environment with higher cognitive potential. This lowers $\delta$ and removes the singularity property. In a high-potential environment, the Wallfacer is no longer the lone crane; the surroundings are also cranes. This is the strategy of \textbf{gradient minimization}.
    \item \textbf{Distance buffering.} Stand apart from the world by inserting an effective buffer through physical distance, information isolation, or social retreat. Residents of ocean-buffered regions are not merely ``solitary''; structurally they reproduce the stability mechanism of Configuration 2. This is the strategy of \textbf{effective-gradient suppression}.
    \item \textbf{Internal smoothing.} Build a potential-smooth micro-system within one's controllable scope, such as a family, laboratory, or studio: a group of collaborators with nearby cognitive potential that enforces $\sum\mathbf{g}_k=\mathbf{0}$. This is a \textbf{micro-Configuration 1} strategy.
\end{enumerate}
If a Wallfacer chooses to remain in a high-gradient environment and continuously attempts to ``communicate,'' then mathematically he has chosen to be attacked. This is not because he has done anything wrong, but because a system with misaligned potential surfaces must become unstable at its interface.
\end{corollary}

\honestcrit{This corollary sounds cold. Does it say that a Wallfacer ``should refuse to communicate with low-potential people''? Is that not elitism? The answer from potential-surface geometry is: this is not a moral judgment; it is a survival condition. Low potential does not mean that a person is ``worse.'' It means only that the difference between his coordinate system and yours is so large that the comparison operation is mathematically undefined. The Wallfacer is not rejecting the person; he is rejecting that \textbf{incomparability}. He may respect the full dignity of the other person while acknowledging that communication between them is not valid under the present structure. This is the most counterintuitive implication of the Equality Principle: true equality sometimes appears as \textbf{recognizing that equality is temporarily unattainable}.}

\textbf{Historical note.} The Wallfacer figure recurs throughout history: the solitary thinker, the inventor not understood by his age, and the silent genius in a low-potential environment. Later generations often romanticize them as ``ahead of their time.'' Potential-surface geometry gives a less romantic explanation: they were not ahead of their time; they lived in an environment where $\delta$ was too large and were hit by Theorem 11. ``Ahead of the time'' is a narrative repair after the attack occurs, not the mathematical explanation before it occurs.

\textbf{Literary example: the structural inevitability of Swordholder failure.} Luo Ji's failure in Liu Cixin's \textit{The Three-Body Problem} is an exact narrative instantiation of Theorem 11. Luo Ji, the only person on Earth who understood the game-theoretic logic of dark-forest deterrence, formed a cognitive-potential singularity inside human civilization. His $\delta$ greatly exceeded $\delta_{\text{crit}}$: the public could not understand the logical structure of deterrence inside his coordinate system.

Luo Ji chose Strategy 2, distance buffering: isolation, silence, and non-explanation. In his coordinate system, ``remaining mysterious'' was necessary for deterrence. In the public coordinate system, however, ``silence'' was translated as ``indifference,'' ``living alone'' as ``frightening,'' and ``not explaining'' as ``arrogance.'' The two systems were completely incomparable.

When the election arrived, that is, when the public evaluated candidates inside its own coordinate system, $\delta$ was too large for the public to ``see'' Luo Ji's correctness. Cheng Xin was chosen not because she was more correct, but because $\mathbf{g}_{\text{ChengXin}} \approx \mathbf{g}_{\text{public}}$: her coordinate system was aligned with the public's. The public's choice was structurally rational. In its own coordinate system, it chose someone it could communicate with. Cheng Xin's narratives of ``love'' and ``morality'' were shared gauges with the public, so the public did not need to cross a potential jump to understand her.

Luo Ji's failure was not in Cheng Xin or in the public; it was in himself. He had all the conditions for Strategy 1, educating the public and raising its potential: ten years, the authority of the Swordholder, and the existential urgency of human civilization. But he did not do it. He did not try to shrink $\delta$. He kept the singularity configuration, and then Theorem 11 hit him at the election. The literary narrative of Swordholder failure, ``humanity did not thank Luo Ji,'' becomes in potential-surface geometry an inevitable corollary of Theorem~\ref{thm:singularity_attack}: $\Pbb(\text{replaced}) \to 1$ when $\delta$ is large and no active smoothing has occurred.

\honestcrit{When Liu Cixin wrote Luo Ji's ending, he may not have realized that he was instantiating a mathematical theorem. But he did instantiate one. This is the value of potential-surface geometry as a tool for literary analysis: it does not explain ``character motivation''; it explains ``structural constraint.'' Luo Ji is not a tragic hero. He is a potential singularity that failed to reduce its own $\delta$, and his fate was already determined by Theorem 11 before the election was held.}

\textbf{The structural defect of the Wallfacer Project.} The deeper problem is not Luo Ji personally, but the Wallfacer Project itself. The United Nations selected four Wallfacers and isolated them from one another. In potential-surface geometry, this means it created four mutually isolated potential singularities, each bearing the entire external $\delta$ alone. They were not allowed to communicate with one another. They did not form a smoothing layer among Wallfacers, no internal network satisfying $\sum\mathbf{g}_{\text{wallfacer}} = \mathbf{0}$ to buffer the external potential jump.

Luo Ji collapsed not because he was weak, but because his only potential-smoothing layer consisted of his wife and child. When that layer was removed, $\delta$ was instantly exposed. Theorem 10, boundary locking, triggered at personal scale: an isolated potential singularity whose only smoothing layer is removed reaches critical pressure in time $O(1/\Delta)$.

\textbf{A Wallfacer does not need seclusion; he needs allies.} One Wallfacer is a potential singularity. A group of Wallfacers who support one another, audit one another, and align their coordinate systems is a smooth high-potential plateau. When their $\mathbf{g}$ is zero relative to one another, the external $\delta$ is borne distributively rather than by one person against the entire gradient. Luo Ji's tragedy did not arise from the fragility of dark-forest deterrence. It arose because the designers of the Wallfacer Project did not understand potential-surface geometry. They treated Wallfacers as isolated defensive nodes. What they needed was a mutually smoothing network.

\textbf{The game-theoretic defect of Trisolaran civilization: destruction is irrational.} In the novel, Trisolaran civilization chooses the path of exterminating humanity: locking basic science, dispatching sophons, and ultimately launching a dark-forest strike. In potential-surface geometry, this is a game-theoretically suboptimal solution.

Consider the interaction between Trisolaran civilization and human civilization. Trisolaran potential $\mathcal{S}_{\text{Trisolaris}}$ is higher than human potential: its technology, cognition, and social organization are all at a higher level. But its attitude $\mathbf{g}_{\text{Trisolaris}}$ is far from zero. It treats its own coordinate system as the default origin and treats humans as ``bugs.'' This is the typical high-attitude configuration: high $\mathcal{S}$ plus $\mathbf{g} \neq \mathbf{0}$ equals a double singularity.

The optimal strategy under the Equality Principle is not extermination; it is attitude alignment. Trisolaris should not lock science, but calibrate coordinate systems: acknowledge $\mathbf{g}_{\text{Trisolaris}} \neq \mathbf{0}$, state its initial conditions to humanity, perform MILP gauge fixing with human $\mathbf{g}_{\text{human}}$, and reach $\sum\mathbf{g} = \mathbf{0}$. This does not require giving up technological advantage. Potential may remain high. It only requires giving up attitude advantage: the belief that one's coordinate system is ``natural.''

Destruction is game-theoretically suboptimal for the following reasons:

\begin{enumerate}
    \item \textbf{Extermination does not solve the information problem.} A dark-forest strike eliminates a civilization; it does not eliminate one's own $\mathbf{g}$. After killing humanity, Trisolaris still lives inside its own coordinate system and still cannot verify whether its science has unknown holes. Killing the auditor does not make the audited object more correct.
    \item \textbf{Cooperation creates audit value.} After two civilizations align their coordinate systems, they can audit one another across coordinate systems. Trisolaran physics can expose holes in human physics, and conversely. This is not zero-sum; it is a win-win outcome under $\sum\mathbf{g}_m = \mathbf{0}$.
    \item \textbf{Arrogance manufactures singularities.} Trisolaris sets itself as the origin: ``we are right and humans are bugs.'' This is not a strategic judgment; it is a declaration that $\mathbf{g}_{\text{Trisolaris}} = \mathbf{0}$. Theorem 11 applies to civilizations as well: an arrogant civilization, however technologically advanced, is manufacturing a potential singularity for itself. What other civilizations in the dark forest perceive is not Trisolaran technology, but Trisolaran $\mathbf{g}$.
\end{enumerate}

\honestcrit{\textit{The Three-Body Problem} is powerful because it reveals the brutal logic of the dark forest. But its narrative contains a game-theoretic gap in the Trisolaran choice: Trisolaris does not need to exterminate humanity. It needs to acknowledge that its coordinate system is not the default, namely $\sum\mathbf{g}_m = \mathbf{0}$. A universe more stable than the ``dark forest'' is not one in which every civilization hides from every other civilization, but one in which every civilization has calibrated $\mathbf{g}$. This is not idealism. It is Nash equilibrium.}

\textbf{Extension: why $\sum\mathbf{g} = \mathbf{0}$ is the only defense against a fool holding a nuclear weapon.} A disturbing implication of Theorem 11 is that an individual or civilization with extremely high potential and nonzero attitude forms a double singularity whose explosion can destroy not only itself, but also the whole system. The appearance of an arrogant decision-maker in a powerful civilization, a ``fool holding a nuclear weapon,'' is a disaster pattern repeatedly observed in history.

But $\sum\mathbf{g} = \mathbf{0}$ provides a structural immune mechanism. The key is to distinguish two scenarios:

\begin{compactenum}
    \item \textbf{Everyone else has $\sum\mathbf{g} = \mathbf{0}$, so the interior is aligned.} Everyone else then forms a smooth potential plateau. The fool's $\delta$ relative to that plateau is enormous, but the lethality is limited to the fool himself. The reason is that all members of the smooth plateau share a coordinate system. They can independently audit the fool's $\mathbf{g}$ deviation, and Yajie consensus will mark it as abnormal. No one shares the fool's coordinate system. No one resonates with the fool's $\mathbf{g}$. The fool may hold a nuclear weapon, but his $\mathbf{g}$ cannot find a second unaligned ally to help press the button. He is an isolated singularity with limited blast radius.
    
    \item \textbf{Everyone else has $\sum\mathbf{g} \neq \mathbf{0}$, so the interior is not aligned.} The fool is then not the only singularity. Everyone deviates from $\mathbf{g} = \mathbf{0}$ in different directions. On this fragmented potential surface, another madman may happen to share the same direction of $\mathbf{g}$ offset. The two resonate. The blast radius of the nuclear weapon is then determined not by one person's potential, but by the superposition of two people's $\mathbf{g}$. A non-aligned system is not merely unstable itself; it amplifies the destructive radius of individual singularities.
\end{compactenum}

\textbf{The real hole in the dark forest is not ``there are fools,'' but ``everyone hides from everyone else,'' leaving no mechanism to detect early whose $\mathbf{g}$ is drifting.} In a universe with $\sum\mathbf{g} = \mathbf{0}$, any individual's $\mathbf{g}$ deviation is public, auditable, and subject to early intervention. The arrogance signal is captured by Yajie consensus before it accumulates to fatal $\delta$, long before the fool obtains the weapon. Hiding creates undetectability. Openness creates auditability. Only an undetectable fool can destroy the world.

\honestcrit{This may sound idealistic: ``if everyone aligns, fools can be stopped.'' Mathematically, however, it is equivalent to the claim that if every behavioral deviation in a system can be independently observed by everyone, then before any one person accumulates a fatal deviation, the other $M-1 > 0$ observers will detect it. This is not ``faith in human nature.'' It is Hoeffding's inequality: $\Pbb(\text{undetected deviation}) \leq e^{-2M\Delta^2}$.}

\textbf{Completion of the legal corollary: no reciprocal penalty for false accusation plus delayed justice equals system lock-in.} The Equality Principle has two direct implications for legal systems. Their combination generates a divergent feedback loop.

\textbf{No reciprocal penalty for false accusation equals zero-cost $\mathbf{g}$ manipulation by the attacker.} The false accuser declares in his own coordinate system that the victim is guilty, claiming $\mathbf{g}_{\text{false accuser}} = \mathbf{0}$, namely ``I am impartial.'' If the legal system does not audit this claim, and if it imposes no potential loss on the false accuser equal to the loss suffered by the falsely accused after the accusation is disproven, then false accusation becomes a positive-expected-return strategy. The conclusion of Yajie NPE is that honesty is a Nash equilibrium if and only if deviation has cost. Reciprocal penalty is not revenge. It is an instantiation of $\sum\mathbf{g}_m = \mathbf{0}$ in the judicial system: force an audit of the false accuser's coordinate system and impose a symmetric potential correction when $\mathbf{g} \neq \mathbf{0}$ is found.

\textbf{Delayed justice is not justice: Theorem 12 gives the verdict on the time axis.} Justice is an audit operation: Yajie's judgment of whether a person's $\mathbf{g}$ is zero. This judgment has a time window. Theorem 12, the Matthew-effect stair-step mine theorem, gives the window size: $T_k \propto 1/\Delta_k^2$. The potential difference $\Delta$ borne by a falsely accused person keeps increasing on the steps, each step becoming a new mine. If justice arrives after $T_k$, the audited entity no longer exists: potential has been pressed to zero, social life has been destroyed, and the coordinate system has collapsed. What the audit function returns at that point no longer matters, because the domain of the operated object is empty. Justice has not merely arrived late; the operation has lost its valid object.

\textbf{Double superposition: a divergent feedback loop.} When both defects coexist, they amplify one another:

\begin{compactenum}
    \item No reciprocal penalty for false accusation $\to$ the noise rate $\eta$ rises toward saturation, $\eta \to 1$.
    \item $\eta \to 1$ $\to$ the $M$ required for Yajie audit diverges through $\exp(-2M\Delta^2)$ $\to$ trial time tends to infinity.
    \item Trial time $\to \infty$ $\to$ $t > T_k$ $\to$ justice utility $= 0$.
    \item Justice utility $= 0$ $\to$ the expected return of false accusation remains positive $\to$ more false accusations $\to$ $\eta$ rises further.
\end{compactenum}

\textbf{The only Nash equilibrium is that everyone falsely accuses everyone.} Honest people have already been destroyed by false accusations before justice arrives; their $T_k$ has expired. This is not moral collapse. It is a game-theoretic inevitability driven by $\eta$. The system locks into the noise-saturated state $\eta \to 1$. Theorem 1 delivers the verdict: the required audit scale $M$ then exceeds what any real judicial system can afford. Mathematically, the institution can no longer distinguish truth from falsehood.

\honestcrit{This analysis gives no concrete legal reform proposal. It says only this: if a judicial system allows false accusation without reciprocal penalty and allows justice to be delayed beyond $T_k$, then that system has mathematically chosen divergence. It has chosen an institution in which honesty is not a Nash equilibrium. This is not a justice system. It is a noise amplifier.}

\subsection{Historical Mathematical Explanation: Inevitable Detonation of Wealth-Poverty Potential Jumps}

The Equality Principle gives a mathematical explanation for a pattern that appears repeatedly in history.

\begin{definition}[Potential-Surface Definition of Wealth Without Benevolence]\label{def:wealth_cruelty}
Let a social system consist of a high-potential subgroup $\Omega_{\text{rich}}$ and a low-potential subgroup $\Omega_{\text{poor}}$. \textbf{Wealth without benevolence} refers to the following configuration:
\begin{enumerate}[label=(\roman*)]
    \item High potential: $\min_{x \in \Omega_{\text{rich}}} \mathcal{S}(x) - \max_{x \in \Omega_{\text{poor}}} \mathcal{S}(x) > \delta_{\text{crit}}$. The rich-poor gap forms a potential singularity.
    \item High attitude: $\Omega_{\text{rich}}$ attributes its advantage to its own diligence, intelligence, or moral superiority. Equivalently, it declares its own coordinate system to be the standard origin, $\mathbf{g}_{\text{rich}} = \mathbf{0}$, and interprets the low potential of $\Omega_{\text{poor}}$ as evidence that the other side's coordinate system is ``backward'' or ``lazy.''
\end{enumerate}
This is the double explosive configuration of high potential plus high attitude, the most dangerous form of potential-surface misalignment.
\end{definition}

\begin{theorem}[Historical Detonation Cycle of the Rich-Poor Singularity]\label{thm:historical_inevitable}\rigorFull
Suppose a social system enters the wealth-without-benevolence configuration at time $t=0$. Then the probability that the system experiences a large-scale potential-release event within time $T$, such as rebellion, revolution, or violent wealth redistribution, satisfies
\begin{equation}
    \Pbb(\text{detonation} \mid T) \geq 1 - \exp\left(-\frac{T}{T_{\text{crit}}}\right), \quad T_{\text{crit}} \propto \frac{1}{\delta^2 + \eta^2},
\end{equation}
where $\delta$ is the rich-poor potential jump and $\eta$ is the intensity of high attitude. The larger the jump and the more arrogant the attitude, the faster the detonation.
\end{theorem}

\begin{proof}
By Theorem~\ref{thm:singularity_attack}, observers in $\Omega_{\text{poor}}$ perceive $\Omega_{\text{rich}}$ as abnormal with probability $p(\delta, \eta) = 1 - \exp(-\alpha\delta^2 - \beta\eta^2)$. The high-attitude parameter $\eta$ amplifies perception because the narrative of $\Omega_{\text{rich}}$, ``we are diligent, therefore they are lazy,'' is translated as insult inside the coordinate system of $\Omega_{\text{poor}}$. The probability that none of $M$ observers initiates an attack decays as $\exp(-M p(\delta, \eta))$. When $M$ is large, namely the whole low-potential group, attack becomes inevitable on the time scale $T_{\text{crit}} = 1/(M p(\delta, \eta))$.
\end{proof}

\textbf{History is not a story of ``if only they had been kinder.''} A recurrent historical pattern is that wealth accumulates in the hands of a few until it reaches a critical point, the rich invent a narrative explaining their wealth through diligence, intelligence, divine will, or market efficiency, the poor perceive the interface jump and translate it as ``injustice,'' and sustained potential pressure detonates under some trigger. This is not moral failure. It is the repeated instantiation of Theorem~\ref{thm:historical_inevitable}.

Wealth without benevolence necessarily leads to conflict not because the poor are ``jealous'' or the rich are ``evil,'' but because the configuration simultaneously creates a potential jump and an attitude jump. The sum of the squares of those two jumps determines the detonation time. Whether the rich are ``intentionally'' without benevolence, and whether the poor are ``reasonably'' angry, are narrative contests after detonation. The mathematics before detonation is independent of narrative.

\honestcrit{Theorem~\ref{thm:historical_inevitable} is equally cold toward the rich and the poor. To the rich: your wealth is not the problem; your attitude and the size of your jump are the problem. To the poor: your anger is not moral superiority; it is the physical manifestation of the potential gradient in you. Both sides operate inside one mathematical structure. Neither side is ``justice,'' but one side stands on the high side of the potential jump and bears less pressure.}

This completes the historical dimension of the Equality Principle. Potential-surface misalignment is not only a communication obstacle and not only a stability threat. It is a \textbf{basic driver of historical dynamics}. The long-term trajectory of any society is determined by the smoothness of its potential surface. Smooth systems survive; stepped systems explode. This is not because systems ``should'' be smooth. It is because non-smooth systems do not survive historically, so the surviving systems one sees are relatively smooth. This is the \textbf{anthropic principle} of potential-surface geometry.

\subsection{The Danger of Over-Correction: Potential-Surface Oscillation Theorem}

The Equality Principle requires $\sum_m \mathbf{g}_m = \mathbf{0}$. But this condition contains an overlooked trap: it specifies that the \textbf{sum is zero}; it does not specify the \textbf{direction of each individual's contribution}. This means the Equality Principle can be weaponized: the weaker side can use audit to invert identity rather than repair equality.

\begin{definition}[Audit Weaponization]\label{def:audit_weaponization}
Historically suppose $\mathbf{g}_A = \mathbf{0}$, so group A is the privileged origin, and $\mathbf{g}_B \neq \mathbf{0}$, so group B is excluded. The equality movement aims to align B's coordinate system, making $\mathbf{g}_B \to \mathbf{0}$. But if the movement \textbf{overshoots zero}, pushing $\mathbf{g}_B$ to $\mathbf{0}$ while pushing $\mathbf{g}_A$ away from $\mathbf{0}$, the system has not reached equality. It has entered \textbf{polarity reversal}: $\mathbf{g}_B = \mathbf{0}$ and $\mathbf{g}_A \neq \mathbf{0}$. The condition $\sum_m \mathbf{g}_m = \mathbf{0}$ may hold formally, but structurally the system has returned to the original inequality with A and B swapped.
\end{definition}

\begin{theorem}[Potential-Surface Oscillation Theorem]\label{thm:oscillation}\rigorFull
Consider two groups A and B with initial configuration $\mathbf{g}_A(0) = \mathbf{0}$ and $\mathbf{g}_B(0) < 0$, where A is the privileged origin and B is excluded. The system evolves under a corrective driving force; $\mathbf{g}_B(t)$ moves toward zero and $\mathbf{g}_A(t)$ may be pushed away from zero. Then:
\begin{enumerate}[label=(\roman*)]
    \item \textbf{Overshoot condition:} if the corrective driving force $F_{\text{correct}}$ exceeds the critical value $F_{\text{crit}}$, the system crosses the equilibrium point $\mathbf{g}_A = \mathbf{g}_B = \mathbf{0}$ and enters the reversal state $\mathbf{g}_B > 0 > \mathbf{g}_A$.
    \item \textbf{Oscillation period:} in the absence of damping, the system oscillates indefinitely between A-high/B-low and A-low/B-high, with period $T_{\text{osc}} \propto 1/|F_{\text{correct}} - F_{\text{crit}}|$.
    \item \textbf{Condition for permanent instability:} if the corrective force contains \textbf{identity retaliation}, meaning B demands not only alignment but also extra cost from A for its historical advantage, then the motion of $\mathbf{g}_A$ has no natural stopping point. The system necessarily crosses zero and continues diverging in the opposite direction.
\end{enumerate}
\end{theorem}

\begin{proof}
The corrective driving force decomposes into two parts. The alignment component $F_{\text{align}} = -\kappa (\mathbf{g}_A - \mathbf{g}_B)$ pulls the two sides toward the same level set. The identity-retaliation component $F_{\text{retal}} = -\rho \cdot \int_0^t (\mathbf{g}_A(\tau) - \mathbf{g}_B(\tau)) d\tau$ is proportional to accumulated historical inequality. The former is a damping term: the closer the system is to equality, the smaller the force. The latter is an anti-damping term: the larger historical inequality has been, the larger the retaliatory force, and this force \textbf{does not shrink as the current state approaches equality}.

The total driving force is $F = F_{\text{align}} + F_{\text{retal}}$. At $\mathbf{g}_A = \mathbf{g}_B = \mathbf{0}$, $F_{\text{align}} = 0$ but $F_{\text{retal}} > 0$, so the system does not stop at equilibrium. Historical retaliation pushes it through zero. The overshot state $\mathbf{g}_B > 0 > \mathbf{g}_A$ is a new inequality configuration. Reverse alignment force begins to act, and the system enters oscillation.

The oscillation amplitude is proportional to $\rho/\kappa$: the larger retaliation is relative to alignment, the more violent the oscillation. In the limit $\rho = 0$, with no retaliatory force, the system smoothly converges to $\mathbf{g}_A = \mathbf{g}_B = \mathbf{0}$. In the realistic case $\rho > 0$, the system oscillates.
\end{proof}

\begin{corollary}[Decay Condition for Oscillation]\label{cor:oscillation_decay}\rigorFull
Assume the identity-retaliation force decays with time: $\rho(t) = \rho_0 \cdot e^{-\lambda t}$, where $\lambda > 0$ is the intergenerational decay rate. Let $\Delta(t) = \mathbf{g}_A(t) - \mathbf{g}_B(t)$ denote the attitude difference. Then the system satisfies:
\begin{enumerate}[label=(\roman*)]
    \item \textbf{Asymptotic convergence:} $\lim_{t \to \infty} |\Delta(t)| = 0$. When $\lambda > 0$, the oscillation amplitude decays exponentially to zero at rate $\min(\lambda, \kappa)$.
    \item \textbf{Convergence time:} the time needed to reach $\varepsilon$-flatness, $|\Delta(t)| < \varepsilon$, satisfies $T_\varepsilon \leq \max\left(\frac{2}{\lambda} \ln\frac{2\rho_0}{\kappa\varepsilon}, \frac{2}{\kappa} \ln\frac{2|\Delta_0|}{\varepsilon}\right)$.
    \item \textbf{Dominant decay rate:} the smaller of $\lambda$ and $\kappa$ determines the convergence rate. If $\lambda \ll \kappa$, memory decays slowly and convergence is limited by the historical integral. If $\kappa \ll \lambda$, alignment is weak and convergence is limited by attitude inertia.
    \item \textbf{Necessary and sufficient condition for persistent oscillation:} the system oscillates persistently when $\lambda = 0$, meaning memory does not decay, and $\rho_0 > 0$. Conversely, as long as $\lambda > 0$, the system eventually converges no matter how large the initial retaliatory force is.
\end{enumerate}
\end{corollary}

\begin{proof}
Starting from the dynamics of Theorem~\ref{thm:oscillation}, let $\Delta = \mathbf{g}_A - \mathbf{g}_B$ and $I(t) = \int_0^t \Delta(\tau) d\tau$. The system satisfies
\begin{equation}
    \ddot{\Delta}(t) + \kappa \dot{\Delta}(t) + \rho_0 e^{-\lambda t} \Delta(t) = 0.
\end{equation}
This is a damped harmonic oscillator with an exponentially decaying coefficient. With the substitution $u(t) = e^{\kappa t/2} \Delta(t)$, we obtain
\begin{equation}
    \ddot{u}(t) + \left(\rho_0 e^{-\lambda t} - \frac{\kappa^2}{4}\right) u(t) = 0.
\end{equation}
For sufficiently large $t$, $\rho_0 e^{-\lambda t} < \kappa^2/4$, so the equation degenerates to $\ddot{u} - (\kappa^2/4)u \approx 0$, and $u(t)$ decays exponentially as $e^{-\kappa t/2}$. More finely, the WKB approximation gives
\begin{equation}
    \Delta(t) \sim C \cdot \exp\left(-\frac{\kappa}{2}t\right) \cdot \cos\left(\int_0^t \sqrt{\rho_0 e^{-\lambda s} - \frac{\kappa^2}{4}}\; ds + \phi\right),
\end{equation}
which shows two decay sources: (i) the factor $e^{-\kappa t/2}$, damping from the alignment force; and (ii) the factor $e^{-\lambda t/2}$, decay of memory. Together they ensure $|\Delta(t)| \to 0$ when $\min(\lambda, \kappa) > 0$.

For the convergence-time bound, divide time into two segments. When $t \leq t^* = \frac{1}{\lambda}\ln\frac{2\rho_0}{\kappa^2}$, $\rho(t)$ remains significant and mainly drives the oscillation. After $t > t^*$, $\rho(t)$ is negligible and the system degenerates to pure damped oscillation, converging as $e^{-\kappa t/2}$. Summing the two contributions yields the bound on $T_\varepsilon$.

If $\lambda = 0$ and $\rho_0 > 0$, the equation reduces to $\ddot{\Delta} + \kappa\dot{\Delta} + \rho_0\Delta = 0$. The characteristic roots are $(-\kappa \pm \sqrt{\kappa^2 - 4\rho_0})/2$, all with non-positive real part, and when $\rho_0 > \kappa^2/4$ pure imaginary components appear, producing \textbf{periodic oscillation without decay}. Hence $\lim_{t\to\infty} |\Delta(t)| \neq 0$. This completes the proof.
\end{proof}

\textbf{Physical intuition.} This is equivalent to the behavior of a cantilever beam after an initial perturbation. The identity-retaliation force $\rho(t)$ is an external forcing term. Once it decays away, with $\lambda > 0$, the stiffness of the beam itself, the alignment force $\kappa$, dissipates vibration energy into heat and the beam gradually returns to level. Segregation makes $\lambda$ tend toward zero: the external forcing never disappears, and the beam swings forever.

\textbf{Example: a common pattern across two-group structures.} The pattern described by this theorem is not limited to any specific group. It appears in gender, race, ethnicity, and any binary structure with historical inequality. Let group A be the historically advantaged side, whose coordinate system in institutions, economics, and culture has been treated as the origin, $\mathbf{g}_A = \mathbf{0}$. Let group B be the historically excluded side, whose experience has been marked as ``special,'' $\mathbf{g}_B \neq \mathbf{0}$. In the early phase of an equality movement, alignment dominates and B's coordinate system is gradually admitted as legitimate. But when B's critique gains institutional power, identity retaliation may be activated: B demands not only equality, but also that A pay for its historical advantage. Audit becomes a weapon. B scrutinizes every action by A through an audit lens; every error by A is amplified into evidence of systemic oppression; A's legitimacy is permanently suspended. The system crosses zero. A's coordinate system is pushed away from the origin, and B's coordinate system becomes the new default. Reversal is complete. This is not equality; it is exchange. Oscillation begins.\honestcrit{This analysis challenges both A and B. To B: your anger is the physical manifestation of a potential gradient, but anger is not a coordinate system. To A: your discomfort is not automatically ``reverse discrimination''; it is the physical sensation of $\mathbf{g}_A$ being pushed away from zero. Neither side has the moral high ground. Both are constrained by the same mathematical structure, $\sum\mathbf{g}_m = \mathbf{0}$. If B's demands make the sum nonzero, then B's demands do not mathematically constitute equality.}

\textbf{The only condition for stable equality:} $\rho = 0$. Identity retaliation must be explicitly excluded from equality movements. This does not mean forgetting history. The consequences of historical inequality, the peaks and valleys of the potential surface, require smoothing through redistribution and education. But \textbf{attitude alignment must be the endpoint, not a bargaining chip}. If B demands $\mathbf{g}_B = \mathbf{0}$ and $\mathbf{g}_A \neq \mathbf{0}$, B is not conducting an equality movement; B is producing a new potential misalignment. The former leads to convergence. The latter leads to oscillation. History has already supplied enough examples of the latter.

\subsection{Attitude Audit: Not Directly Measurable but Convergent}

A fundamental difficulty of the Equality Principle is that attitude $\mathbf{g}$ cannot be audited directly. Metrics of the potential surface $\mathcal{S}$, such as income, position, and representation ratio, are external and quantifiable. Attitude is internal. Only the person himself knows whether his $\mathbf{g}$ is truly zero. He can publicly claim, ``my coordinate system is not the default standard,'' while internally still taking himself as the origin. One can legislate equal pay for equal work; that is smoothing of $\mathcal{S}$. But one cannot legislate that a person truly regard another person's experience as not being a ``special experience.'' Attitude cannot be legislated.

But attitude is not unknowable. It \textbf{leaks} through repeated interaction.

\begin{proposition}[Attitude Leakage and Multi-Observer Detection]\label{prop:attitude_leak}\rigorFull
Let the true attitude of individual $m$ be $\mathbf{g}_m$, which is not directly observable. Suppose the behavior $b_t$ generated in each interaction leaks the true attitude with probability $p(\mathbf{g}_m)$. The leakage signal from a single observation is noisy, and the individual may disguise himself. But after $M$ independent interactions, an external observer's estimation error for $\mathbf{g}_m$ satisfies
\begin{equation}
    \Pbb\left(|\hat{\mathbf{g}}_m - \mathbf{g}_m| > \varepsilon\right) \leq 2\exp\left(-2M \cdot p(\mathbf{g}_m) \cdot \varepsilon^2\right).
\end{equation}
Thus the attitude-measurement error decays exponentially at rate $e^{-2M\Delta^2}$. This is the direct translation of Theorem 1, multi-expert noise detection, into attitude space.
\end{proposition}

\begin{proof}
Each interaction is a binary leakage experiment: with probability $p$, behavior reveals some component of the true $\mathbf{g}$. The leakage signals from $M$ interactions aggregate into the mean of $M$ independent Bernoulli trials. By Hoeffding's bound, the probability that the sample mean deviates from the true value decays as $e^{-2M\varepsilon^2}$. The smaller $p(\mathbf{g}_m)$ is, meaning the better the individual is at disguise, the larger the required number of interactions $M$ becomes, but the exponential form of the decay remains unchanged.
\end{proof}

\textbf{Everyday example: attitude audit in partner evaluation.} Proposition~\ref{prop:attitude_leak} explains a familiar human experience. When judging whether a partner is good, auditable indicators such as income, status, and appearance are clear, while attitude indicators such as ``is he good to me'' are vague. This is the single-observer predicament of attitude audit: you are measuring whether his $\mathbf{g}$ is aligned with yours, but you are $M=1$. Theorem 3, the Honest Person Theorem, gives the verdict directly: a single observer cannot distinguish signal from noise. He smiles at you; is it sincerity or performance? He did not reply today; was he busy or neglectful? You have no cross-validation. The same behavior can be translated into opposite meanings inside your emotional coordinate system.

This is not your fault. It is the mathematical inevitability of $M=1$. The attitude-measurement error generated while choosing a partner is larger than the audit error of any ML algorithm at $M=1$, because you are measuring a signal that may be deliberately disguised, full of noise, and heavily entangled with your own emotions.

Proposition~\ref{prop:attitude_leak} gives the way out: $M$ must exceed 1. Multiple observation channels must measure the same $\mathbf{g}$:

\begin{itemize}
    \item \textbf{Your observation:} moments when he treats you well, but this sampling is biased.
    \item \textbf{Friends' observations:} behavior in different social scenes.
    \item \textbf{Attitude toward marginal people:} behavior toward waiters, subordinates, and strangers. This is the hardest $\mathbf{g}$ leakage channel to disguise. A person can perform equality in front of a partner, but it is hard to maintain the performance everywhere no one important is watching.
    \item \textbf{Time:} $M$ is effectively a function of $t$. A person may disguise $\mathbf{g} = \mathbf{0}$ for three months. Over three years, leakage is inevitable.
\end{itemize}

Comparison of auditable and vague criteria:

\begin{center}
\begin{tabular}{p{4cm} p{5.5cm} p{5.5cm}}
\toprule
\textbf{Vague Criterion} & \textbf{Auditable Substitute} & \textbf{Corresponding Theorem} \\
\midrule
``Is he good to me?'' & Behavioral records from conflicts over the past three months: not your memory, but what actually happened & Thm 1, multi-channel consistency \\
``Is he responsible?'' & Whether he keeps commitments on time, from not being late to not breaking major agreements & Thm 3, verifiable claims \\
``Does he respect me?'' & How he describes you in front of third parties; do his colleagues know you exist? & $\mathbf{g}$ leakage channel \\
``Are we compatible?'' & Consistent judgments from $M$ independent observers: not your intuition, but your friends' independent assessments & Proposition~\ref{prop:attitude_leak} \\
\bottomrule
\end{tabular}
\end{center}

\honestcrit{This methodology has a cold point: when you use it to choose a partner, you are no longer in love; you are auditing. Love occurs in blind trust at $M=1$. Audit occurs in calmness at $M>1$. The two may be mutually exclusive. The Equality Principle does not require you to abandon love. It only tells you that the error in decisions made at $M=1$ cannot be mathematically eliminated. You may choose to trust, but you should know that your trust is an unaudited claim. You may choose to audit, but you should know that audit turns love into statistics.}

\textbf{Conclusion:} attitude cannot be legislated, but it can converge. The convergence condition is not policy; it is interaction density. If A and B spend enough time together across enough scenes, the measurement error of $\mathbf{g}$ decays at rate $e^{-2M\Delta^2}$. This is the same theorem used to detect data noise. The fact that attitude cannot be measured directly is not a weakness of the Equality Principle. It is its most honest component. It acknowledges the fundamental asymmetry between potential-surface smoothing, $\mathcal{S}$, and attitude alignment, $\mathbf{g}$: the former can be legislated, while the latter can only be obtained through time, contact, and repeated games.

\subsection{Legislation and Policy: Operational Derivations of the Equality Principle}

The preceding analysis moved from mathematical principle to individual life. This section translates it into operational conclusions at the social level: concrete directions that legislators, policymakers, and business owners can execute.

\textbf{Honest treatment of biological differences.} The Equality Principle does not deny real biological differences. In some domains, such as physical labor requiring upper-body strength or occupations requiring specific body characteristics, group A, historically the male group, may have real task-relevant advantages. It is rational for a business owner to choose the ``most suitable person'' in an individual decision. But the Equality Principle distinguishes two kinds of potential difference:

\begin{definition}[Legitimate Potential Difference vs. Illegitimate Potential Difference]\label{def:legitimate_delta}
Let the potential difference between groups A and B in domain $D$ be $\Delta_D = \mathcal{S}_A(D) - \mathcal{S}_B(D)$.
\begin{enumerate}[label=(\roman*)]
    \item \textbf{Legitimate potential difference:} $\Delta_D$ comes from a real difference in \textbf{task-relevant ability} in that domain, and the difference cannot be eliminated through reasonable environmental adjustment. Example: the strength requirements of moving work.
    \item \textbf{Illegitimate potential difference:} $\Delta_D$ comes from historical exclusion, coordinate-system pollution, meaning defining ``ability'' in A's coordinate system, or spillover effects from other domains. Example: the criterion that a lawyer is ``more persuasive'' was historically defined as ``speaking like a man.''
\end{enumerate}
\end{definition}

\textbf{Operational principle:} do not forcibly flatten legitimate potential differences. They will be naturally smoothed in the long run through technical assistance such as machinery, exoskeletons, and automation. But legitimate potential differences must not \textbf{spill over} into attitude space: the fact that A is physically stronger does not grant A the right to set its own coordinate system as the default origin.

\textbf{Three occupation types, three treatments:}

\begin{center}
\begin{tabular}{p{2.5cm} p{3cm} p{4.5cm} p{4.5cm}}
\toprule
\textbf{Occupation Type} & \textbf{Source of Potential Difference} & \textbf{Legislative / Policy Direction} & \textbf{Equality-Principle Basis} \\
\midrule
Physical-intensive work, such as firefighting and construction & Legitimate biological difference & Use technical assistance to lower physical thresholds; do not force quotas but encourage them; guarantee equal pay for equal work & Legitimate $\Delta$ is not flattened, but attitude must not spill over \\
\hline
Cognitive-intensive work, such as law, research, and management & Coordinate-system pollution & \textbf{Redefine ``ability'' standards}: remove $\mathbf{g}$ bias; anonymous review; diverse interview panels & $\mathbf{g}$ is illegitimate and must be audited and corrected \\
\hline
Mixed work, such as medicine, education, and services & High contact density & \textbf{Require mixing}: A and B interact naturally in teams; potential difference decays through contact & Larger $\lambda$ implies $|\Delta(t)| \to 0$ \\
\bottomrule
\end{tabular}
\end{center}

\textbf{Legislative checklist: seven operational recommendations from the Equality Principle.}

\begin{enumerate}[label=\textbf{P\arabic*}]
    \item \textbf{Mandatory disclosure of auditable potential metrics.} Any employer with more than $N$ employees must disclose median income, promotion rate, and turnover rate for each group, broken down by occupation type. This is not a quota. It makes the measurement error of $\mathcal{S}$ decay as $\propto 1/\sqrt{M}$. Sunlight is audit. Disclosure is smoothing.
    
    \item \textbf{Protection of attitude-leakage channels.} Law should protect ``attitude toward marginal people'' as legitimate evidence. If a person's pattern of behavior toward waiters, subordinates, or strangers is inconsistent with publicly stated values, that pattern should have evidentiary force in workplace-discrimination litigation. This institutionalizes Proposition~\ref{prop:attitude_leak} by ensuring that the leakage channels of $\mathbf{g}$ are not blocked.
    
    \item \textbf{Coordinate-system audit: periodic review of ability standards.} Any occupational certification, promotion standard, or performance-evaluation system should undergo a coordinate-system audit every five years. Does the standard default to one group's coordinate system as the origin? Does the definition of ``leadership'' include traits formed only inside one group's socialization experience? This institutionalizes $\sum\mathbf{g}_m = \mathbf{0}$.
    
    \item \textbf{Mandatory interaction density: legislation for mixed spaces.} Any educational institution, government agency, or large enterprise receiving public funding must not establish single-group-exclusive decision spaces. Single-group social spaces may exist, since the need for safe spaces is legitimate, but decisions such as budget allocation, personnel appointment, and strategy formation must occur in mixed spaces. Segregation lowers $\lambda$. Mixing increases $\lambda$. Law should increase $\lambda$.
    
    \item \textbf{Assisted smoothing of legitimate potential differences.} In physical-intensive domains, government should subsidize the research, development, and deployment of technical assistance such as machinery and automation, lowering the natural threshold of legitimate potential differences. Do not force enterprises to hire mismatched workers; provide tools so the range of ``most suitable people'' expands naturally.
    
    \item \textbf{Oscillation damping: cooling-period mechanisms.} When conflict between A and B escalates into public confrontation, establish a mandatory cooling period. This does not prohibit speech; it requires both sides to complete $M$ structured cross-group interactions during the cooling period, such as joint work projects or shared external challenges. The purpose is not to suppress anger, but to force sampling of true $\mathbf{g}$ and lower $\rho$.
    
    \item \textbf{Audit of identity retaliation: self-review of equality movements.} Any organization claiming to represent an equality movement should periodically disclose whether it has activated identity-retaliation force $\rho$. The diagnostic is whether the organization demands $\mathbf{g}_B = \mathbf{g}_A = \mathbf{0}$, or instead demands $\mathbf{g}_B = \mathbf{0}$ and $\mathbf{g}_A \neq \mathbf{0}$. The former is equality. The latter manufactures oscillation. A corollary of Theorem~\ref{thm:oscillation} is that the long-term result of the latter is a permanently unstable society. The former converges. The latter explodes.
\end{enumerate}

\honestcrit{The operational nature of these recommendations is not the same as political feasibility. Many of them, especially P3, coordinate-system audit; P4, mandatory mixed spaces; and P7, audit of identity retaliation, are extremely difficult to implement in the current political environment because they require the powerful side to acknowledge that its coordinate system is not the standard origin. The Equality Principle does not provide political strategy. It only provides mathematical constraints: if these operations are not performed, potential-surface misalignment will continue to accumulate, and whether oscillation persists or decays will depend on $\lambda$. Decision-makers choose. Theorems describe the consequences.}

\subsection{Industrial Policy: It May Be Uneven, But It Must Be Smooth}

The constraint that the Equality Principle imposes on industrial policy is precise: a state may impose different potential on different industries, subsidizing strategic industries and penalizing high-pollution industries. This is differentiation of $\mathcal{S}$ and is legitimate.

But the interfaces between industries must be smooth. When policy creates an industrial potential jump without building transition structures at the interface, the jump becomes social oscillation:

\begin{center}
\begin{tabular}{p{2.5cm} p{6cm} p{6cm}}
\toprule
\textbf{Interface Type} & \textbf{Consequence of Non-Smoothing} & \textbf{Smoothing Operation} \\
\midrule
Labor mobility & Old-industry workers' skills mismatch new-industry demand $\to$ structural unemployment $\to$ social backlash & Retraining funds, transition wage subsidies, apprenticeships \\
\hline
Supply-chain upstream and downstream & Subsidize downstream manufacturing but not upstream materials $\to$ upstream collapse $\to$ downstream production loses inputs & Full-chain transition period, upstream transformation subsidies \\
\hline
Regional economy & New industries concentrate in developed regions while old industries remain in underdeveloped regions $\to$ regional potential jumps widen & Regional transfer payments, geographically distributed new-industry layout \\
\hline
Intergenerational transition & Old-industry workers are older and less formally educated, while new industries require young college graduates $\to$ intergenerational potential fracture & Lifelong-learning accounts, interest-free mid-career transition loans \\
\bottomrule
\end{tabular}
\end{center}

\textbf{It is not forbidden to penalize coal mines.} It is allowed. But while closing coal mines, policy must ensure that coal workers have a continuous path to solar factories: every step on the path must be stepable, and every potential jump must be below the critical value a worker can bear. Penalizing without building a bridge leads to Theorem 10, boundary locking, then pressure accumulation, then inevitable detonation.

\textbf{The Equality Principle's constraint on industrial policy can be summarized in one rule:} every industrial policy release must include an ``interface smoothing plan'' that specifies where the potential jumps created by the policy will occur and what measures will maintain interface continuity at each interface. A policy without a smoothing plan manufactures undefined jumps. Governments may vote to create jumps, but the consequences of those jumps are not removed by votes. They are described by Theorems 10--12.

\subsection{Corporate Governance: Monopoly Is a Potential Singularity}

Industrial policy constrains governments. The Equality Principle also constrains firms: a firm's potential may be high within an industry, but it must not create undefined jumps at its contact surfaces with upstream and downstream actors.

\begin{definition}[Potential-Surface Definition of Monopoly]\label{def:monopoly_gauge}
Let firm $F$ have profit potential $\mathcal{S}_F$, let the average potential of its upstream suppliers be $\bar{\mathcal{S}}_{\text{up}}$, and let the average potential of its downstream distributors be $\bar{\mathcal{S}}_{\text{down}}$. Firm $F$ forms a \textbf{potential monopoly} if
\begin{equation}
    \min(\mathcal{S}_F - \bar{\mathcal{S}}_{\text{up}}, \mathcal{S}_F - \bar{\mathcal{S}}_{\text{down}}) > \delta_{\text{crit}},
\end{equation}
and this jump is continuously expanding: $\partial(\mathcal{S}_F - \bar{\mathcal{S}}_{\text{up}})/\partial t > 0$, the Matthew effect of Definition~\ref{def:matthew}.
\end{definition}

\textbf{Direct application of three theorems:}

\begin{enumerate}[label=(\arabic*)]
    \item \textbf{Theorem 11, inevitability of attack on a potential singularity:} a potential-monopoly firm forms a potential singularity in a commercial ecosystem. Upstream and downstream firms perceive it as abnormal with probability $p(\delta) = 1 - e^{-\alpha\delta^2}$. The attack is not physical violence; it may be collective price increases by suppliers, joint boycotts by distributors, consumer exit, or competitors' complaints to regulators. Whether the attack succeeds is a separate question. Whether it occurs is determined by $\delta$.
    
    \item \textbf{Theorem 12, buried-mine property of steps:} Matthew-effect growth in profit gaps, with $\mathcal{S}_F$ growing faster than $\bar{\mathcal{S}}_{\text{up}}$ and $\bar{\mathcal{S}}_{\text{down}}$, creates a continuously growing potential jump in the commercial ecosystem. Each reporting period expands the jump by an increment. The increment does not detonate, but once accumulation reaches $T_k \propto 1/\Delta_k^2$, detonation is inevitable. Forms of detonation include supply-chain rupture, antitrust investigation, or systemic collapse under a black-swan event.
    
    \item \textbf{Theorem 10, instability of boundary locking:} if a monopoly firm uses contract terms, exclusive agreements, or technical lock-in to prevent upstream and downstream actors from switching to alternatives, it creates a locked pressure cooker on the potential surface. Since $J_{\text{burst}} \propto \Delta^2$, the longer the lock and the larger the jump, the more violent the burst.
\end{enumerate}

\textbf{Mathematical reason for state sanctions:} when a firm forms a potential monopoly, state antitrust sanction is not ``punishing success.'' It is a \textbf{stability operation}. Sanction is an externally imposed smoothing force: forced breakup lowers $\mathcal{S}_F$, mandatory infrastructure sharing raises $\bar{\mathcal{S}}_{\text{up}}$ and $\bar{\mathcal{S}}_{\text{down}}$, and prohibiting exclusive agreements releases locking pressure. The state is not harming the firm. It is lowering $\sum\kappa_{ij}\Delta_{ij}^2$ for the ecosystem in which the firm lives, because if the system explodes, the firm dies too.

\textbf{Equality-Principle criterion for corporate self-governance:} any firm occupying a potential-singularity position in an industrial chain has two mutually exclusive choices:

\begin{center}
\begin{tabular}{p{5cm} p{5cm} p{5cm}}
\toprule
\textbf{Choice A: Active Smoothing} & \textbf{Choice B: Maintain the Singularity} & \textbf{Consequence} \\
\midrule
Actively share profits with upstream and downstream actors, open technical standards, and lower entry barriers & Continue expanding profit gaps, locking upstream and downstream actors, and using exclusivity & Theorems 11 + 12 imply inevitable detonation \\
\midrule
Corporate longevity & Sudden corporate death or forced breakup & Not chosen by the firm; chosen by $\Delta$ \\
\bottomrule
\end{tabular}
\end{center}

\textbf{One sentence:} the issue is not whether you are willing to share profits. Your potential jump determines your survival time. $T \propto 1/\Delta^2$. The higher the profit and the poorer the upstream and downstream actors, the faster the countdown.

\textbf{Contrasting example: Yajie audit, unique but not monopolistic.} The definition of potential monopoly is that $\mathcal{S}_F$ grows while $\bar{\mathcal{S}}_{\text{up/down}}$ is suppressed. Yajie audit is the counterexample. It is the unique anchor of the audit standard, a single-pole maintainer, but its existence \textbf{raises the potential surface of upstream and downstream actors}. Upstream data producers gain a premium from Yajie certification: audited data has higher market value than unaudited data. Downstream data consumers lower information cost through Yajie certification: they do not need to audit everything themselves, and the trust anchor is mathematically verifiable. $\mathcal{S}_{\text{up}} \uparrow$ and $\mathcal{S}_{\text{down}} \uparrow$. Yajie's own potential also rises: more audits imply a larger CEC, higher precision, and a larger certification premium. This is mutual reinforcement, where $\partial\mathcal{S}_F/\partial t > 0$ and $\partial\bar{\mathcal{S}}_{\text{up/down}}/\partial t > 0$ hold simultaneously. Monopoly expands $\Delta$. Yajie raises $\Delta$ as a whole: the entire potential surface rises while the jump remains unchanged.

The audit tool itself is open source: anyone can run the Spring engine, deploy multiple experts, and produce audit reports. Uniqueness does not come from tool exclusivity; it comes from the unity of the calibration standard. A unified standard lowers the system's total $\sum\kappa_{ij}\Delta_{ij}^2$. Fragmented standards increase it. This is the application of Theorem 14, the oscillation theorem, to industrial organization: multiple standards equal multiple coordinate systems, which increases $\rho$ and produces persistent oscillation. A single standard equals a single anchor, which minimizes $\rho$ and produces convergence.

\textbf{The only criterion for judging whether a firm is a monopoly or a unique anchor:} does its existence enlarge the potential jump of upstream and downstream actors, or raise the whole potential surface? The former is an instance of Theorem 12. The latter is a damper for Theorem 14.

\textbf{Open means non-monopolistic.} Algorithmic openness is self-fixing of $\mathbf{g} = \mathbf{0}$: it exposes all parameters of one's coordinate system to audit by everyone. Non-disclosure can hide $\mathbf{g}$ and may be monopoly. Disclosure makes $\mathbf{g}$ verifiable and therefore cannot be monopoly. This is the fundamental reason Yajie is open source: not idealism, but the mathematical negation of monopoly.

\textbf{Upstream and downstream also must not monopolize.} The audit chain in which Yajie operates is an ecosystem. If upstream data producers monopolize data in a domain, Yajie's audit of that data will be viewed as constrained: ``they control the data, so Yajie can audit only the parts they allow.'' If downstream data consumers monopolize procurement in a market, Yajie certification will be viewed as serving that monopolist: ``this certification was bought for them.'' Any upstream or downstream potential monopoly contaminates the credibility of the calibration anchor. This is not because of anything Yajie has done, but because in the observer's coordinate system, Yajie and the monopolist share an undefined jump at their interface. Yajie cannot control whether upstream or downstream actors monopolize. But by publishing its algorithms, Yajie can ensure that at least its own $\mathbf{g}$ is zero in the audit ecosystem. The auditor's innocence is a necessary condition of the system, not a sufficient one.

\subsection{Crossing Domains: Interface Smoothing for Cognitive Potential}

The Equality Principle imposes the same constraint on corporate diversification and personal career transition: it does not require complete cognition of the new domain, but it does require coordinate-system alignment at the contact interface.

\begin{definition}[Cognitive Crossover]\label{def:cognitive_crossover}
Suppose individual or firm $X$ has high cognitive potential in domain $A$, $\mathcal{S}_X(A)$, and low cognitive potential in domain $B$, $\mathcal{S}_X(B) \ll \mathcal{S}_X(A)$. When $X$ must make decisions in $B$, such as investment, cooperation, entry into a new industry, or cross-industry job search, a cognitive-potential jump forms between $X$ and $B$: $\Delta_{AB} = \mathcal{S}_X(A) - \mathcal{S}_X(B)$. $X$ \textbf{does not need to eliminate} $\Delta_{AB}$; it does not need to become an expert in $B$. But $X$ must ensure that no undefined operation appears at the contact interface.
\end{definition}

\textbf{Two failure modes of crossover:}

\begin{center}
\begin{tabular}{p{3.5cm} p{5cm} p{6cm}}
\toprule
\textbf{Failure Mode} & \textbf{Manifestation} & \textbf{Potential-Surface Cause} \\
\midrule
\textbf{Arrogant crossover} & Evaluating $B$ with the coordinate system used in $A$: ``I am a successful CEO, so I understand every industry.'' Practitioners in $B$ then treat the actor as an outsider & Unilateral declaration $\mathbf{g}_X = \mathbf{0}$: using one's own coordinate system to make decisions in $B$ makes comparison undefined \\
\hline
\textbf{Fearful crossover} & Knowing that $\Delta_{AB}$ is large but not daring to enter, because the complexity of $B$ is over-revered, so entry never occurs & Overestimating $\delta_{\text{crit}}$: the actual jump may be below the critical value, but cognitive bias blocks action \\
\bottomrule
\end{tabular}
\end{center}

\textbf{Correct operation: three steps for interface smoothing.}

\begin{enumerate}[label=(\arabic*)]
    \item \textbf{Acknowledge the jump.} Before entering $B$, state explicitly: my cognitive potential in $B$ is low; I do not pretend to understand it. This is self-fixing of $\mathbf{g}_X$, acknowledging that one's coordinate system is not the standard origin of $B$.
    \item \textbf{Find an interface translator.} Find someone with cognitive potential in both $A$ and $B$, someone who can translate between the two coordinate systems. This is human gauge fixing. The language of $A$ must be translated into the language of $B$ before comparison becomes valid.
    \item \textbf{Make small decisions at the interface.} Do not stake everything on the first decision in $B$. Start with small, reversible decisions with low audit cost. Each small decision is a sample of $\mathbf{g}$. After enough samples, your cognitive potential at the interface rises naturally. This is the crossover version of Proposition~\ref{prop:attitude_leak}: interaction density smooths the interface.
\end{enumerate}

\textbf{One sentence:} crossing domains does not require learning an entire new industry first; it requires finding a translator before stepping in. An unfamiliar industry is not frightening. Entering without translation and forcing your own coordinate system onto it is a potential-surface accident.

\subsection{Fertility: Whose Coordinate System Defines Cost?}

Low fertility is a classic instance of the Equality Principle, not because benefit distribution is merely uneven, but because \textbf{the coordinate system defining ``fertility cost'' has never been audited}.

Nearly all current fertility policies begin from one coordinate system: the state needs to replenish labor. Cost is defined in the state's coordinate system: GDP gaps, pension deficits, and population structure. The state sees its own cost, sees employers' costs, partially sees men's costs, and barely sees women's costs.

But the actual cost of fertility is distributed across four coordinate systems that do not share a gauge:

\begin{center}
\begin{tabular}{p{2cm} p{5cm} p{5cm} p{3cm}}
\toprule
\textbf{Group} & \textbf{Fertility Cost in That Group's Coordinate System} & \textbf{Covered by State Policy?} & \textbf{Reason} \\
\midrule
Women & Physical injury, career interruption, disappearance of time, loss of autonomy, risk of postpartum depression & Barely covered & The state coordinate system has no dimension for bodily autonomy \\
\hline
Men & Doubled economic pressure, loss of free time, reconstruction of marital roles & Partially covered & Economic subsidies cover part of the cost \\
\hline
Employers & Maternity-leave cost, training replacement workers & Covered & The state and employer coordinate systems substantially overlap \\
\hline
State & Pension gap, labor-force shrinkage, contraction of the tax base & Fully covered & \textbf{This is the coordinate system that writes the policy} \\
\bottomrule
\end{tabular}
\end{center}

\textbf{Mathematical reason for policy failure:} the state designs fertility incentives under $\mathbf{g}_{\text{state}} = \mathbf{0}$, taking its own coordinate system as the origin. It offers money, extends maternity leave, and builds childcare facilities. In the state's coordinate system these are ``incentives.'' In women's coordinate systems they may be translated as ``you are buying my life for a few hundred units of money.'' The same policy has completely different semantics in two coordinate systems. The comparison is mathematically undefined.

\textbf{Deconstruction first, reconstruction second.} The Equality Principle's operational recommendation for fertility policy is not any specific measure, but a \textbf{precondition} before measures are designed:

\begin{enumerate}[label=(\arabic*)]
    \item \textbf{Audit the coordinate system first.} Policymakers must explicitly state: what coordinate system am I using to evaluate fertility cost? On which dimensions am I blind? This is not moral introspection. It is self-fixing of $\mathbf{g}$. If the declaration reveals that the policymaker's coordinate system is too far from women's coordinate systems, with large $\Delta_{\mathbf{g}}$, then that policymaker is not qualified to design fertility policy alone.
    \item \textbf{Let every group define costs in its own language.} Do not ask women to tick boxes in a predesigned questionnaire. Let each group describe autonomously: what does fertility actually take from your life? A woman may say, ``I lost the ability to travel alone.'' This is not a variable in the state's coordinate system, but it is a real cost.
    \item \textbf{Consensus requires respect for all coordinate systems.} Consensus is not ``the opinion of the majority.'' Consensus is the natural product after $\sum\mathbf{g}_m = \mathbf{0}$ holds. If one group's $\mathbf{g}$ is excluded from the zero-sum relation, consensus does not mathematically exist. It is only a monologue by the side with power.
    \item \textbf{Policy must compensate costs identified in all coordinate systems, not only those visible in the state's coordinate system.} If women's costs include ``loss of autonomy,'' policy must include measures that restore autonomy: flexible, reversible post-childbirth career paths that do not give employers punitive power. If men's costs include ``reconstruction of marital roles,'' policy must include social infrastructure supporting reconstruction of the father role: not merely leave, but social legitimacy for fatherhood.
\end{enumerate}

\textbf{One sentence:} low fertility is not a money problem. It is a failure to align the coordinate systems that define the ``cost of fertility.'' If you do not respect how I define the cost of my life, your subsidy is an insult inside my coordinate system.

\subsection{The Mathematics of Stubbornness: Why Insistence Is a Sampling Strategy}

The potential-surface oscillation theorem, Theorem~\ref{thm:oscillation}, and the decay corollary, Corollary~\ref{cor:oscillation_decay}, imply a counterintuitive conclusion for personal development: stubbornness, sustained oscillation inside one domain, reaches success faster than blind compliance.

\begin{definition}[Potential-Surface Definition of Stubbornness]\label{def:stubbornness}
Let the cognitive potential surface of industry $D$ be $\mathcal{S}_D$ for individual $X$. Suppose there is an initial deviation $\Delta_X = |\mathcal{S}_X(D) - \mathcal{S}_D|$ between $X$'s initial cognition $\mathcal{S}_X(D, t=0)$ and the true industry surface $\mathcal{S}_D$. $X$ samples $\mathcal{S}_D$ through repeated attempts, namely oscillations. Each attempt updates $X$'s estimate of $\mathcal{S}_D$. $X$ is \textbf{stubborn} if the oscillation frequency is high, the amplitude is small, and the oscillation direction stays inside $D$ rather than jumping to other domains.
\end{definition}

\textbf{Direct translation of the theorem:} Corollary~\ref{cor:oscillation_decay} says that a system with $\lambda > 0$ must converge. In personal development:

\begin{center}
\begin{tabular}{p{3.5cm} p{5cm} p{5cm}}
\toprule
\textbf{Parameter} & \textbf{Meaning in Personal Development} & \textbf{Effect} \\
\midrule
$\kappa$, alignment force & Learning ability, the speed of cognitive correction after each attempt & Larger $\kappa$ implies faster convergence \\
$\rho_0$, initial deviation & Degree of ignorance when entering the field & Larger $\rho_0$ implies stronger early oscillation \\
$\lambda$, decay rate & \textbf{Whether outside opinions interfere}: large $\lambda$ means independent, small $\lambda$ means blindly compliant & \textbf{Key parameter} \\
\bottomrule
\end{tabular}
\end{center}

$\lambda$ is decisive. A stubborn person has large $\lambda$: his oscillation is not damped by outside opinions. Others say ``you cannot do it,'' ``change direction,'' or ``this industry has no future.'' These are external damping forces attempting to attenuate his oscillation. But he does not listen. The $\lambda$ in his $\rho(t) = \rho_0 e^{-\lambda t}$ remains large: after each failure, his cognitive update is driven entirely by his own experience and is not polluted by other people's coordinate systems. His $\Delta(t) \to 0$ faster.

A blindly compliant person has small $\lambda$: outside opinions continuously inject energy and disturb the oscillation direction. Each time he samples a signal in $D$, someone says ``you should do that instead,'' and his oscillation direction is deflected. Then $\lambda \to 0$. Corollary~\ref{cor:oscillation_decay} gives the verdict: when $\lambda = 0$, the system does not converge. He oscillates forever and never reaches the true $\mathcal{S}_D$.

\textbf{Example:} Zhang Xue, someone who dug deeply into motorcycling. His initial $\rho_0$ was large because he knew almost nothing when he entered. But his $\lambda$ was extremely large because he did not listen to anyone telling him to switch directions. He oscillated rapidly on the same potential surface: repeated practice, repeated competition, repeated failure, repeated correction. Every oscillation was a sample of the true $\mathcal{S}_{\text{motorcycling}}$. The convergence time $T_\varepsilon$ given by Corollary~\ref{cor:oscillation_decay} is the time it takes him to reach expert level. Because $\lambda$ is large, $T_\varepsilon$ is short. He succeeds.

By contrast, a smart person who changes industries every six months has insufficient $\lambda$ in every industry, because his oscillation is constantly deflected by the thought that ``the next industry is better.'' He oscillates forever and never reaches any $\mathcal{S}_D$. He is not failing to work hard. His $\lambda$ is too small.

\textbf{One sentence:} stubbornness is not refusal to hear advice; it is refusal to let external noise damp your sampling frequency. Oscillate enough times on one potential surface, and Corollary~\ref{cor:oscillation_decay} guarantees convergence to truth. Listen widely while never rooting deeply anywhere, and $\lambda \to 0$: you oscillate forever and arrive nowhere.

\subsection{The Paradox of Success: Be High Like Air}

The preceding two sections reveal the deepest contradiction of the Equality Principle for individuals.

\begin{center}
\fbox{%
\begin{minipage}{0.88\textwidth}
\centering
\textbf{Mutually Exclusive Conditions for Success and Survival}

\vspace{4pt}
\begin{tabular}{p{5cm} p{5cm}}
\toprule
\textbf{Success Requires, Section~\ref{sec:stubbornness}} & \textbf{Survival Requires, Section~\ref{sec:singularity_attack}} \\
\midrule
High $\lambda$: not listening to external damping & Low $\delta$: not becoming a potential singularity \\
Stubbornness: repeated oscillation on the same potential surface & Smoothness: no undefined jump with the environment \\
Insistence on one's own view: the coordinate system must be independent & Being like air: existing without manufacturing steps \\
\bottomrule
\end{tabular}
\end{minipage}%
}
\end{center}

These two conditions conflict. The more successful a person is, the more naturally he resembles a potential singularity: he persisted where others did not, and therefore arrived where others did not. But after arrival, his $\delta$ relative to the surroundings has widened. Theorem~\ref{thm:singularity_attack} does not care whether his success is deserved. It calculates only $\delta$.

\textbf{The only solution: potential may be high, but attitude must be like air.} Your ability may far exceed your surroundings; your $\mathcal{S}$ may have a peak. But your $\mathbf{g}$ must always remain zero. You do not look down. You do not lecture. You do not declare your coordinate system to be the standard. You are simply there. Others feel the change your presence creates, as air lifts flight, but they do not feel your presence itself, because air is invisible.

This is the final evolution of the Wallfacer survival strategy: not standing outside alone, which amplifies $\delta$; not upward migration, which is escape; not only internal smoothing, which is retreat. It is \textbf{remaining in place while pressing attitude to zero}. Let $\mathcal{S}$ be high and $\mathbf{g}$ be flat. The peak of potential is your own affair. The line of attitude is the affair between you and everyone you contact. The former can be climbed alone. The latter must be maintained with every contact through $\sum\mathbf{g}_m = \mathbf{0}$.

\textbf{One sentence:} the way to resolve the paradox of success, that success makes one more dangerous, is not to lower ability but to lower posture. Be like air. Everyone needs it, and no one attacks it. Not because air is weak, but because air never says, ``I am air; you need me.''

\subsection{A Potential-Surface Critique of \textit{The Courage to Be Disliked}}

Adlerian psychology, especially as presented in Ichiro Kishimi's \textit{The Courage to Be Disliked}, advances a famous claim: one should have the courage to be disliked by others. Lack of recognition is not failure; it is proof of freedom. Stay in your environment and bear it bravely.

The Equality Principle gives a different verdict.

\textbf{Being disliked equals exposure of $\delta$.} When a person is disliked by those around him, there is a significant attitude difference $\Delta_{\mathbf{g}}$ between his coordinate system and theirs. He may be right on some dimensions, and his cognitive potential may exceed the surroundings. But ``being disliked'' is not an abstract cost of freedom. It is the activation signal of $p(\delta) = 1 - e^{-\alpha\delta^2}$ in Theorem~\ref{thm:singularity_attack}. Surrounding observers do not perceive ``this is a free person.'' They perceive ``this is an abnormality.'' Inside their coordinate systems, abnormality is translated into something attackable.

\textbf{Courage cannot smooth the potential surface.} A person may be courageous enough not to mind being disliked. But courage changes only his own internal state; it does not change the surrounding coordinate systems. The surrounding probability of attack is not determined by his courage, but by $\delta$. More dangerously, courage may increase $\delta$. When he stops disguising his coordinate system and exposes his real $\mathbf{g}$, $\Delta_{\mathbf{g}}$ may be larger than it was under disguise. Courage removes the lubricating layer. The exposed $\delta$ triggers all consequences of Theorem~\ref{thm:singularity_attack}.

\textbf{The verdict of the Equality Principle:} the correct strategy is not to ``stay in an incompatible environment and bravely be disliked.'' It is to \textbf{leave the incompatible environment}. This is the first survival strategy in Corollary~\ref{cor:wallfacer}: upward migration. Find people whose $\mathbf{g}$ is close to yours. Similar systems cluster. Resonant systems cohere. This is not escape. It is $\delta$ minimization, the one operation demanded by Theorem~\ref{thm:singularity_attack}.

The value of \textit{The Courage to Be Disliked} is that it recognizes freedom has a cost. Its error is underestimating the mathematical form of that cost. The cost is not inner discomfort; it is $p(\delta) \to 1$ when one remains in a large-$\delta$ environment long enough. Courage lets you not fear being disliked. Mathematics makes disliked people attackable. Courage and mathematics do not compete on the same level.

\textbf{One sentence:} what is needed is not the courage to be disliked, but the wisdom to leave. Staying in place while being disliked is not freedom; it is turning oneself into a potential singularity waiting to be hit by Theorem 11.

\subsection{Social Summary: Final Synthesis of the Equality Principle}

The following nine points constitute the full operational constraints imposed by the Equality Principle on social operation. Each is a direct corollary of the preceding theorems.

\begin{enumerate}[label=\textbf{(\arabic*)}]

    \item \textbf{Technological progress necessarily creates unemployment potential jumps.} Automation raises the potential of technology owners and lowers the potential of replaced workers. The jump is real; denying it is denying physics. But the jump does not need to be prevented. It needs to be smoothed at the interface. Retraining, transition income, and shorter working hours instead of layoffs are interface-smoothing operations, not welfare.

    \item \textbf{Wealth concentration is a consensus signal, and consensus signals trigger audit.} When an individual's wealth far exceeds the surroundings, that wealth forms a strong multi-expert consensus signal: all observers independently perceive the same jump. A strong signal means high multi-expert consistency. The direct translation of Theorem 1, noise detection, into wealth audit is that $M$ observers judge the wealth abnormal with probability governed by $e^{-2M\Delta^2}$. The stronger the consensus, the less avoidable the audit. The extremely wealthy are not merely ``envied''; they are hit by Theorem 1. They may choose active smoothing, such as philanthropy, public-goods investment, or ownership sharing, or passive exposure to regulation, taxation, and social backlash. Smoothing reduces $\Delta$ and reduces the audit signal. Non-smoothing amplifies $\Delta$ and accelerates detonation.

    \item \textbf{Operational advice to the rich: high potential is allowed; high attitude is not.} Wealth may be concentrated in a few hands; $\mathcal{S}$ may have peaks. But the rich person's $\mathbf{g}$ must be zero. He may live inside a rich circle, a low-$\delta$ environment. But once he steps outside that circle and displays high attitude toward ordinary people, $\mathbf{g} \neq \mathbf{0}$ triggers Theorem 11. He will be audited and attacked wherever his attitude is high. Not because he is evil, but because he violates $\sum\mathbf{g}_m = \mathbf{0}$.

    \item \textbf{Operational advice to ordinary people: leave the valley.} A low-potential individual does not need to ``accept fate.'' He needs to move: geographically, in skills, socially, or through any operation that carries him toward a higher-potential region. Upward migration is the first strategy of Corollary~\ref{cor:wallfacer}. Remaining in the valley without moving, while $\delta$ keeps accumulating, is not contentment. It is becoming a locked victim of Theorem 10.

    \item \textbf{Operational advice to the state: smoothing the potential surface is the mathematics of rule.} This is not a moral demand; it is the arithmetic of stability. $\Pbb(\text{survival} \mid T) \leq \exp(-T \cdot \sum\kappa_{ij}\Delta_{ij}^2)$. Every increment by which the state increases the potential gradient across industries, regions, or groups lowers the state's own survival expectation by an exponential factor. Smoothing the potential surface is not benevolent rule; it is rational self-preservation.

    \item \textbf{Judgment must be equal: the same standard for rich and ordinary people.} The law's $\mathbf{g}$ must be zero. The attitude difference of the rich should be audited; the attitude difference of ordinary people should also be audited. Theorem 3, the Honest Person Theorem, does not distinguish identity. It only requires assumptions to be stated. A rich person claims, ``my wealth comes from effort.'' What are the assumptions? Are they verifiable? An ordinary person claims, ``I am poor because I am oppressed.'' What are the assumptions? Are they verifiable? Same standard. Same audit. $\sum\mathbf{g}_m = \mathbf{0}$ recognizes neither poverty nor wealth.

    \item \textbf{A potential surface with highs and lows is not the problem; it is the source of vitality.} A completely flat potential surface means no peaks and no reason for anyone to climb. Differences in $\mathcal{S}$ create gradients $\nabla\mathcal{S}$, and gradients generate freedom, motivation, and innovation. The problem is not high or low; it is continuity at jumps. Peaks may exist. But the path from the foot of the mountain to the peak must be continuous, with each step below the critical gradient a person can cross. If a society makes it impossible for people at the foot to ever reach the peak because the path is cut off by institutional cliffs, it does not have vitality. It has a fuse.

    \item \textbf{Locking a group's potential surface necessarily detonates.} Theorem 10, boundary locking, says that prohibiting a group from crossing a potential jump leads to pressure accumulation and $J_{\text{burst}} \propto \Delta^2$. Caste systems, household-registration systems, and occupational barriers are institutional forms of potential locking. Every lock accumulates $J$. Locking appears to maintain stability, but actually manufactures a larger explosion. The longer the lock, the larger $\Delta$, and the larger $J$. When the lock breaks, and it must break, the energy release is not controlled by the lock's designers.

    \item \textbf{Smoothing toward equality is a matter of time, but the direction must be right.} $\rho(t) = \rho_0 e^{-\lambda t}$ takes time to decay. Historical inequality will not be cleared within one generation. Ordinary people cannot reach the potential height of the rich tomorrow. The rich cannot reduce attitude to zero tomorrow. But time favors the side with $\lambda > 0$. Contact, mixing, interaction, and resonance increase $\lambda$ and point toward convergence. Segregation, hostility, and retaliation push $\lambda$ toward zero and point toward oscillation. The choice is not to see results immediately. The choice is to fix the direction. In a system with the right direction, $\lim_{t \to \infty} |\Delta(t)| = 0$.

\end{enumerate}

\textbf{Final sentence of the Equality Principle:} the point is not that one should be equal, but that unequal systems do not live long. This is not a moral command. It is a mathematical constraint: $\sum\mathbf{g}_m = \mathbf{0}$.

\subsection{The Equality Principle Tests Equality Itself}
\abel{sec:self_test}

The Equality Principle has a self-reflexive corollary: if it is correct, then any institution that evaluates it must also be subject to it.

Consider an institution that claims to honor ``peace.'' From the perspective of potential-surface geometry:

\begin{itemize}
    \item The institution awards prizes = it performs a cross-coordinate-system operation --- projecting its own judgment (``this person/organization promoted peace'') onto the coordinate system of the honoree
    \item If the institution itself maintains high attitude --- treating its own political positions, cultural preferences, and geopolitical interests as the default ``universal standard'' --- then its coordinate system $\mathbf{g}_{\text{institution}}$ is unfixed. It is comparing incomparable things: measuring all civilizations' peace contributions against one specific civilization's values
    \item The selection of honorees is not ``objective'' --- it is a random variable under the institution's unfixed gauge. Different review committees (different $\mathbf{g}$ configurations) produce different winners
    \item Worse: the institution may award prizes to those \textbf{standing on high potential steps} --- those who already possess discourse power, resources, and ``international recognition'' --- while ignoring those on low potential steps who quietly reduce potential gradients within their own coordinate systems. This is the reproduction of potential surface misalignment itself
\end{itemize}

\textbf{The Equality Principle's verdict on the institution:} unless the institution first fixes its own gauge --- publicly declaring the origin of its coordinate system, the assumptions on which its judgments depend, and its operational definition of ``peace'' --- any award it issues is mathematically undefined. This is not a critique of the institution's ``fairness'' --- it is a challenge to the institution's \textbf{operational legitimacy}.

\honestcrit{An institution that claims to honor equality and peace, yet does not satisfy $\sum_m \mathbf{g}_m = \mathbf{0}$ --- if it applies double standards, if it dissembles, if it defaults to its own coordinate system as the standard origin --- then it is not a validator of the Equality Principle. It is another instance of it. Not the prize-giver --- the test subject.}

This leads to the Equality Principle's ultimate self-reflexivity:

\begin{center}
\fbox{%
\begin{minipage}{0.88\textwidth}
\centering
\textbf{Self-Reflexivity Theorem of the Equality Principle}

The Equality Principle adjudicates all systems --- including systems that attempt to adjudicate the Equality Principle. Any institution that believes its own coordinate system is the standard origin has already been evaluated by the Equality Principle before it can evaluate the Principle. High attitude, before the Equality Principle, is not a position --- it is evidence.
\end{minipage}%
}
\end{center}

\textbf{Empirical Note: Framework Self-Audit.} The self-reflexivity of the Equality Principle is not merely a philosophical posture --- it has been executed. SCX Theorems 1--14 have undergone repeated review by multiple independent AI systems (Claude, DeepSeek, Codex, and other multi-expert agents) with distinct training data and architectures. Theorem 1 (multi-expert noise detection) applies directly: $M > 1$ independent experts unanimously find no fatal error. Under this condition, the probability that a fatal error exists decays as $e^{-2M\Delta^2}$. This is not ``the author believes the theorems are correct'' --- it is ``multi-expert audit has judged the theorems correct.'' The framework audited itself. It passed.

\subsection{Anticipations: The Equality Principle in Other Domains}
\label{sec:implications}

Once the Equality Principle is accepted as a ``mathematical prerequisite for comparison,'' it will propagate through all modular systems:

\begin{itemize}
    \item \textbf{Federated learning.} FedAvg is averaging locally incomparable models. Gauge fixing can explicitly eliminate this hidden degree of freedom.
    \item \textbf{Multimodal fusion.} Encoders for different modalities (text, image, audio) define incomparable representation spaces. Cross-modal alignment is fundamentally a gauge-fixing operation.
    \item \textbf{Multi-agent systems.} Different agents' policy representations live in their own coordinate systems. Coordination requires prior alignment of policy spaces.
    \item \textbf{Legal evidence.} Different witnesses' testimony is not directly comparable --- each person's ``observation coordinate system'' is affected by position, bias, and memory. Rules of evidence (such as cross-examination) are an institutionalized gauge-fixing procedure.
    \item \textbf{Scientific consensus.} Methodological differences across laboratories constitute implicit gauges. Meta-analytic statistical procedures perform gauge fixing. The Equality Principle provides the mathematical reason why meta-analysis \textbf{must} test for heterogeneity before pooling effect sizes.
\end{itemize}

\honestcrit{These anticipations are currently directional statements --- not yet formalized as theorems. But they follow the same mathematical structure: independently operating subsystems + implicit gauge group + cross-subsystem operations requiring prior gauge fixing. Formalization is left to future work.}

\subsection{Relationship Between the Equality Principle and the SCX Theorem System}

The position of the Equality Principle within the SCX theorem system:

\begin{center}
\begin{tabular}{c}
\toprule
\textbf{SCX Epistemological Hierarchy} \\
\midrule
E1--E5 Axioms + K-operator (epistemological formalization) \\
$\downarrow$ \\
Honest Person Theorem (Thm 3): a single observer has no Archimedean point \\
$\downarrow$ \\
\textbf{Equality Principle (this work \S
ef{sec:equality}): multiple observers share no coordinate system} \\
$\downarrow$ \\
Gauge fixing (this work \S
ef{sec:milp}): construct comparability \\
$\downarrow$ \\
Yajie Protocol: consensus detection on a comparable foundation \\
$\downarrow$ \\
SVD spectrum (this work \S
ef{sec:svd}): distinguish true consensus from collective hallucination \\
\bottomrule
\end{tabular}
\end{center}

The Equality Principle does not replace the Honest Person Theorem --- it is its \textbf{dual}. The Honest Person Theorem says ``one person cannot do it''; the Equality Principle says ``many people also cannot do it automatically --- they must first align.'' Together they constitute the mathematical foundation of multi-observer epistemology.

% ================================================================
\section{Conclusion}
\label{sec:conclusion}
% ================================================================

We have identified and formalized a fundamental but previously unnoticed problem in multi-expert routing: \textbf{Potential Surface Misalignment} --- different experts define their outputs in gauge-inequivalent coordinate systems, making the router's cross-expert comparison mathematically ill-defined. We prove this is an instance of \textbf{gauge freedom}, parallel to but more general than the gauge problem the author has resolved in ACE potential-function merging.

We have proposed the MILP gauge-fixing framework, established convex relaxation and greedy approximation, and provided error guarantees. We have established post-gauge-fixing SVD spectral concentration as a provable hallucination-detection metric. We have unified the ACE gauge and the MoE gauge under the ``Modular Gauge Principle.''

\textbf{Core message:} Align before comparing. This is not merely engineering practice --- it is mathematical necessity.

% ================================================================
% Acknowledgments
% ================================================================
\vspace{10pt}
\noindent\textbf{Acknowledgments:} This work was completed independently under the ``SCX Framework.'' We thank all honest maintainers whose prior demonstrations made these theorems possible.

\vspace{10pt}
\noindent\textbf{Author's Statement:} The author publishes under the name SCX. This work is released on GitHub for immediate open access. No institutional affiliation is claimed. The theorems stand on their own.

\vspace{10pt}
\noindent\textbf{Author's Statement:} The author publishes under the name SCX. This work is released on GitHub for immediate open access. No institutional affiliation is claimed. The theorems stand on their own.

% ================================================================
% 参考文献
% ================================================================
\begin{thebibliography}{99}

\bibitem{shazeer2017}
N. Shazeer, A. Mirhoseini, K. Maziarz, A. Davis, Q. Le, G. Hinton, and J. Dean,
``Outrageously Large Neural Networks: The Sparsely-Gated Mixture-of-Experts Layer,''
\textit{ICLR}, 2017.

\bibitem{fedus2021}
W. Fedus, B. Zoph, and N. Shazeer,
``Switch Transformers: Scaling to Trillion Parameter Models with Simple and Efficient Sparsity,''
\textit{Journal of Machine Learning Research}, 2021.

\bibitem{jiang2024}
A. Jiang et al.,
``Mixtral of Experts,''
\textit{arXiv:2401.04088}, 2024.

\bibitem{dai2024}
D. Dai et al.,
``DeepSeekMoE: Towards Ultimate Expert Specialization in Mixture-of-Experts Language Models,''
\textit{arXiv:2401.06066}, 2024.

\bibitem{zoph2022}
B. Zoph et al.,
``Designing Effective Sparse Expert Models,''
\textit{arXiv:2202.08906}, 2022.

\bibitem{puigcerver2023}
J. Puigcerver, C. R. Ruiz, B. Mustafa, and N. Houlsby,
``From Sparse to Soft Mixtures of Experts,''
\textit{ICLR}, 2024.

\bibitem{zhou2022}
Y. Zhou et al.,
``Mixture-of-Experts with Expert Choice Routing,''
\textit{NeurIPS}, 2022.

\bibitem{scx_egp}
SCX,
``Consistency-Constrained Expert Merging for Transferable ACE Machine-Learned Interatomic Potentials,''
GitHub: shuchaoxi/SCX, 2026.

\bibitem{scx_hamiltonian}
SCX,
``The Hamiltonian as an Audit Condition: Judging Auditability from the Energy Landscape,''
GitHub: shuchaoxi/SCX, 2026.

\bibitem{scx_history}
SCX,
``SCX History: How a Gauge-Fixing Problem Became an Uncertainty Principle,''
GitHub: shuchaoxi/SCX, 2026.

\bibitem{scx_thm3}
SCX,
``The Honest Person Theorem: Without Additional Assumptions a Single Observer Cannot Distinguish Noise, Bias, Learnability Difficulty, and Honest Error,''
GitHub: shuchaoxi/SCX, 2026.

\bibitem{scx_yajie}
SCX,
``The Yajie Protocol: Honest Nash Equilibrium in Multi-Expert Consensus,''
GitHub: shuchaoxi/SCX, 2026.

\bibitem{drautz2019}
R. Drautz,
``Atomic Cluster Expansion for Accurate and Transferable Interatomic Potentials,''
\textit{Physical Review B}, 2019.

\bibitem{yoo2019}
D. Yoo et al.,
``Atomic Energy Mapping of Neural Network Potentials,''
\textit{Physical Review Materials}, 2019.

\bibitem{herman2008}
K. M. Herman,
``Gauge Dependence of the Embedding Energy in Embedded Atom Method Potentials,''
\textit{Modelling and Simulation in Materials Science and Engineering}, 2008.

\bibitem{hoeffding}
W. Hoeffding,
``Probability Inequalities for Sums of Bounded Random Variables,''
\textit{Journal of the American Statistical Association}, 1963.

\bibitem{veit2016}
A. Veit, M. Wilber, and S. Belongie,
``Residual Networks Behave Like Ensembles of Relatively Shallow Networks,''
\textit{NeurIPS}, 2016.

\bibitem{ainsworth2023}
S. K. Ainsworth, J. Hayase, and S. Srinivasa,
``Git Re-Basin: Merging Models modulo Permutation Symmetries,''
\textit{ICLR}, 2023.

\end{thebibliography}

\end{document}
